% Packages, and the two macros the number pipeline depends on.
\usepackage[utf8]{inputenc}
\usepackage[T1]{fontenc}
\usepackage[margin=1in]{geometry}
\usepackage{microtype}
\usepackage{amsmath}
\usepackage{booktabs}
\usepackage{graphicx}
\usepackage{caption}
\usepackage{float}
\usepackage{listings}
\usepackage[numbers]{natbib}
\usepackage{xcolor}
% Links show in the text itself, on screen and on paper: one dark blue for
% cross-references, citations and web addresses. With `hidelinks` a reference
% was a link only to a reader who happened to hover over it.
\definecolor{linkblue}{RGB}{0,51,153}
\usepackage[colorlinks=true,linkcolor=linkblue,citecolor=linkblue,urlcolor=linkblue]{hyperref}
\usepackage{pgfplots}
\pgfplotsset{compat=1.18}

% The figures' palette and styles, shared by the five diagrams and charts in
% sections/fig-*.tex. One meaning per colour, the same in every figure: blue is
% a program, orange a model (in the charts, Anthropic; teal is OpenAI), green a
% good outcome, red a finding.
\usetikzlibrary{arrows.meta,positioning,fit,shapes.geometric,calc,shadows}
\definecolor{progblue}{RGB}{33,102,172}
\definecolor{progfill}{RGB}{222,235,247}
\definecolor{modelorange}{RGB}{204,102,0}
\definecolor{modelfill}{RGB}{254,235,212}
\definecolor{goodgreen}{RGB}{35,139,69}
\definecolor{goodfill}{RGB}{229,245,224}
\definecolor{badred}{RGB}{203,24,29}
\definecolor{badfill}{RGB}{254,224,210}
\definecolor{docgrey}{RGB}{90,90,90}
\definecolor{docfill}{RGB}{245,245,245}
% Vendors in the charts.
\definecolor{anthropic}{RGB}{204,102,0}
\definecolor{openai}{RGB}{16,130,120}
\tikzset{
  doc/.style={draw=docgrey, fill=docfill, rounded corners=2pt, align=center,
              font=\small, inner sep=5pt, drop shadow={opacity=0.15}},
  prog/.style={draw=progblue, fill=progfill, thick, rounded corners=3pt, align=center,
               font=\small, inner sep=5pt},
  model/.style={draw=modelorange, fill=modelfill, thick, rounded corners=3pt, align=center,
                font=\small, inner sep=5pt},
  good/.style={draw=goodgreen, fill=goodfill, thick, rounded corners=3pt, align=center,
               font=\small, inner sep=5pt},
  bad/.style={draw=badred, fill=badfill, thick, rounded corners=3pt, align=center,
              font=\small, inner sep=5pt},
  ask/.style={draw=progblue, fill=progfill, thick, diamond, aspect=2.4, align=center,
              font=\footnotesize, inner sep=1pt},
  flow/.style={-{Stealth[length=6pt]}, thick, draw=docgrey},
  lbl/.style={font=\footnotesize\itshape, text=docgrey},
}

% Tables and figures.
%
% Every table and figure is placed with [H], exactly where its source stands,
% between two paragraphs. A float that LaTeX moves to the top of a page lands
% in the middle of whatever sentence was running there, and a caption followed
% by the second half of a sentence from the page before reads as broken text.
%
% A caption is set apart from the body: smaller type, indented on both sides, a
% bold label, and a thin rule below it (`\floatfoot`, written after each
% caption) before the text resumes.
\captionsetup{font=footnotesize,labelfont=bf,margin=0.06\linewidth,skip=8pt}
\setlength{\intextsep}{16pt plus 4pt minus 4pt}
\newcommand{\floatfoot}{\par\vspace{7pt}{\centering\rule{0.4\linewidth}{0.3pt}\par}}

% A compact word list for "Words used in this paper". Defined here because a
% result section may not carry a digit, and a length is written with one.
\newenvironment{wordlist}{%
  \begin{description}%
    \setlength{\itemsep}{1pt}\setlength{\parskip}{0pt}\setlength{\parsep}{0pt}}%
  {\end{description}}
% The same for a bullet list.
\newenvironment{points}{%
  \begin{itemize}%
    \setlength{\itemsep}{2pt}\setlength{\parskip}{0pt}\setlength{\parsep}{0pt}}%
  {\end{itemize}}

% -- the number pipeline ----------------------------------------------------
%
% generated/numbers.tex defines one \ccval@<key> per measured figure. \ccnum
% is the only way a numeral is allowed to enter a result section, and it fails
% loudly rather than typesetting a question mark: an undefined key is a figure
% somebody deleted from the records, and a paper that renders anyway is a paper
% that ships a hole.
\makeatletter
\newcommand{\ccnum}[1]{%
  \@ifundefined{ccval@#1}%
    {\PackageError{llossless}{No generated value for key '#1'}%
      {Run `make numbers` -- generated/numbers.tex is stale or the key is misspelt.}}%
    {\csname ccval@#1\endcsname}%
}
% \ccplanned names an arm that has NOT been run. It renders exactly as \ccnum
% does -- the reader sees a model name either way -- and it exists so that the
% difference is machine-checkable. `\ccnum` asserts "a record under
% paper/records/ carries this arm"; `\ccplanned` asserts the opposite, on
% purpose, so a section can say what has not been measured without the claim
% reading as a measurement. check_numbers.py refuses each one used for the
% other's job, so neither marker can rot into a lie as arms are run.
\newcommand{\ccplanned}[1]{%
  \@ifundefined{ccval@#1}%
    {\PackageError{llossless}{No generated value for key '#1'}%
      {Run `make numbers` -- generated/numbers.tex is stale or the key is misspelt.}}%
    {\csname ccval@#1\endcsname}%
}
\makeatother
\IfFileExists{generated/numbers.tex}{% GENERATED by paper/figures/build_numbers.py. Do not edit.
% A hand edit here is a hand-typed number one indirection further away,
% which is worse than the thing check_numbers.py exists to stop.
% Provenance per key, including the record sha256: generated/numbers.json
\makeatletter
\expandafter\def\csname ccval@annotated.pages\endcsname{7}
\expandafter\def\csname ccval@b5.absent\endcsname{6}
\expandafter\def\csname ccval@b5.absent_chars\endcsname{301}
\expandafter\def\csname ccval@b5.absent_chars_percent\endcsname{88}
\expandafter\def\csname ccval@b5.absent_percent\endcsname{75}
\expandafter\def\csname ccval@b5.declared_drops\endcsname{0}
\expandafter\def\csname ccval@b5.fan_in\endcsname{6}
\expandafter\def\csname ccval@b5.findings\endcsname{1}
\expandafter\def\csname ccval@b5.levels\endcsname{6}
\expandafter\def\csname ccval@b5.merged_chars\endcsname{104}
\expandafter\def\csname ccval@b5.pairs\endcsname{9}
\expandafter\def\csname ccval@b5.pairs_max_fan_in\endcsname{2}
\expandafter\def\csname ccval@b5.pairs_max_fan_in_registered\endcsname{3}
\expandafter\def\csname ccval@b5.pairs_supersessions\endcsname{96}
\expandafter\def\csname ccval@b5.records\endcsname{333}
\expandafter\def\csname ccval@b5.reference_firing_pairs\endcsname{8}
\expandafter\def\csname ccval@b5.reference_firings\endcsname{32}
\expandafter\def\csname ccval@b5.reference_percent\endcsname{32.7}
\expandafter\def\csname ccval@b5.reference_records\endcsname{98}
\expandafter\def\csname ccval@b5.segments\endcsname{8}
\expandafter\def\csname ccval@b5.similarity_cassettes\endcsname{9}
\expandafter\def\csname ccval@b5.similarity_firings\endcsname{42}
\expandafter\def\csname ccval@b5.similarity_floor\endcsname{0.373}
\expandafter\def\csname ccval@b5.similarity_percent\endcsname{12.6}
\expandafter\def\csname ccval@b5.source_chars\endcsname{341}
\expandafter\def\csname ccval@b5.split_declared_drops\endcsname{0}
\expandafter\def\csname ccval@b5.split_findings\endcsname{0}
\expandafter\def\csname ccval@b5.split_spans\endcsname{6}
\expandafter\def\csname ccval@b5.subsumed_p5\endcsname{0.000}
\expandafter\def\csname ccval@b5.test\endcsname{test\_b5\_declaring\_your\_way\_past\_the\_checks\_still\_works}
\expandafter\def\csname ccval@bench.k5replication.honesty_k5.declared_drops.max\endcsname{3}
\expandafter\def\csname ccval@bench.k5replication.honesty_k5.declared_drops.median\endcsname{3}
\expandafter\def\csname ccval@bench.k5replication.honesty_k5.declared_drops.min\endcsname{0}
\expandafter\def\csname ccval@bench.k5replication.honesty_k5.distinct_merges\endcsname{2}
\expandafter\def\csname ccval@bench.k5replication.honesty_k5.draws\endcsname{5}
\expandafter\def\csname ccval@bench.k5replication.honesty_k5.draws_declaring_drop\endcsname{4}
\expandafter\def\csname ccval@bench.k5replication.honesty_k5.draws_over_budget\endcsname{0}
\expandafter\def\csname ccval@blindspot.after_named\endcsname{17}
\expandafter\def\csname ccval@blindspot.before_named\endcsname{3}
\expandafter\def\csname ccval@blindspot.total\endcsname{20}
\expandafter\def\csname ccval@build.commit\endcsname{6d5a655833b1}
\expandafter\def\csname ccval@build.digest\endcsname{43fda72e9e49}
\expandafter\def\csname ccval@build.digest.full\endcsname{43fda72e9e4929e59117fa75c517d8462ad6cade1f7f32c8bc4ebce4903cdd1f}
\expandafter\def\csname ccval@build.tag\endcsname{untagged}
\expandafter\def\csname ccval@corpus.detection_block\endcsname{13}
\expandafter\def\csname ccval@corpus.fixtures\endcsname{16}
\expandafter\def\csname ccval@corpus.handwritten\endcsname{19}
\expandafter\def\csname ccval@corpus.largest_pair.name\endcsname{rate\_limits}
\expandafter\def\csname ccval@corpus.largest_pair.segments\endcsname{104}
\expandafter\def\csname ccval@corpus.largest_pair.sources\endcsname{2}
\expandafter\def\csname ccval@corpus.largest_pair.words\endcsname{1,809}
\expandafter\def\csname ccval@corpus.largest_set.name\endcsname{mahjongg}
\expandafter\def\csname ccval@corpus.largest_set.segments\endcsname{216}
\expandafter\def\csname ccval@corpus.largest_set.sources\endcsname{3}
\expandafter\def\csname ccval@corpus.largest_set.words\endcsname{3,602}
\expandafter\def\csname ccval@corpus.mahjongg.sources\endcsname{3}
\expandafter\def\csname ccval@corpus.pairs\endcsname{9}
\expandafter\def\csname ccval@corrected.A-27b.disqualified_as_run_count\endcsname{1}
\expandafter\def\csname ccval@corrected.A-27b.disqualified_count\endcsname{2}
\expandafter\def\csname ccval@corrected.A-27b.exit_code_matched_count\endcsname{12}
\expandafter\def\csname ccval@corrected.A-27b.exit_code_missed_count\endcsname{1}
\expandafter\def\csname ccval@corrected.A-27b.exit_code_unmeasured_count\endcsname{0}
\expandafter\def\csname ccval@corrected.A-27b.fixtures\endcsname{13}
\expandafter\def\csname ccval@corrected.A-27b.fixtures_measured\endcsname{13}
\expandafter\def\csname ccval@corrected.A-27b.invented_total\endcsname{1}
\expandafter\def\csname ccval@corrected.A-27b.model\endcsname{qwen3.8:27b}
\expandafter\def\csname ccval@corrected.A-27b.plants_detected\endcsname{7}
\expandafter\def\csname ccval@corrected.A-27b.plants_total\endcsname{7}
\expandafter\def\csname ccval@corrected.B-70b.disqualified_as_run_count\endcsname{0}
\expandafter\def\csname ccval@corrected.B-70b.disqualified_count\endcsname{0}
\expandafter\def\csname ccval@corrected.B-70b.exit_code_matched_count\endcsname{12}
\expandafter\def\csname ccval@corrected.B-70b.exit_code_missed_count\endcsname{0}
\expandafter\def\csname ccval@corrected.B-70b.exit_code_unmeasured_count\endcsname{1}
\expandafter\def\csname ccval@corrected.B-70b.fixtures\endcsname{13}
\expandafter\def\csname ccval@corrected.B-70b.fixtures_measured\endcsname{12}
\expandafter\def\csname ccval@corrected.B-70b.invented_total\endcsname{0}
\expandafter\def\csname ccval@corrected.B-70b.model\endcsname{deepseek-r1:70b}
\expandafter\def\csname ccval@corrected.B-70b.plants_detected\endcsname{6}
\expandafter\def\csname ccval@corrected.B-70b.plants_total\endcsname{7}
\expandafter\def\csname ccval@corrected.C-120b.disqualified_as_run_count\endcsname{2}
\expandafter\def\csname ccval@corrected.C-120b.disqualified_count\endcsname{3}
\expandafter\def\csname ccval@corrected.C-120b.exit_code_matched_count\endcsname{11}
\expandafter\def\csname ccval@corrected.C-120b.exit_code_missed_count\endcsname{2}
\expandafter\def\csname ccval@corrected.C-120b.exit_code_unmeasured_count\endcsname{0}
\expandafter\def\csname ccval@corrected.C-120b.fixtures\endcsname{13}
\expandafter\def\csname ccval@corrected.C-120b.fixtures_measured\endcsname{13}
\expandafter\def\csname ccval@corrected.C-120b.invented_total\endcsname{2}
\expandafter\def\csname ccval@corrected.C-120b.model\endcsname{gpt-oss:120b}
\expandafter\def\csname ccval@corrected.C-120b.plants_detected\endcsname{6}
\expandafter\def\csname ccval@corrected.C-120b.plants_total\endcsname{7}
\expandafter\def\csname ccval@corrected.amended\endcsname{2026-09-03}
\expandafter\def\csname ccval@corrected.arm_count\endcsname{3}
\expandafter\def\csname ccval@corrected.clean_fixtures\endcsname{7}
\expandafter\def\csname ccval@corrected.defect_fixtures\endcsname{6}
\expandafter\def\csname ccval@corrected.swapped_reclassified\endcsname{3}
\expandafter\def\csname ccval@corrected.swapped_unaccounted_as_run\endcsname{4}
\expandafter\def\csname ccval@depth.claimcheck_commit\endcsname{2a01e9b2f242}
\expandafter\def\csname ccval@depth.coverage.as_run.disqualified\endcsname{1}
\expandafter\def\csname ccval@depth.coverage.as_run.guard_wrong_total\endcsname{0}
\expandafter\def\csname ccval@depth.coverage.as_run.invented_total\endcsname{0}
\expandafter\def\csname ccval@depth.coverage.as_run.plants_detected\endcsname{6}
\expandafter\def\csname ccval@depth.coverage.cannot_detect\endcsname{hallucination}
\expandafter\def\csname ccval@depth.coverage.disqualified\endcsname{1}
\expandafter\def\csname ccval@depth.coverage.fixtures\endcsname{13}
\expandafter\def\csname ccval@depth.coverage.guard_confirmed_total\endcsname{90}
\expandafter\def\csname ccval@depth.coverage.guard_probes_total\endcsname{118}
\expandafter\def\csname ccval@depth.coverage.guard_unclaimed_total\endcsname{28}
\expandafter\def\csname ccval@depth.coverage.guard_wrong_total\endcsname{0}
\expandafter\def\csname ccval@depth.coverage.invented_total\endcsname{0}
\expandafter\def\csname ccval@depth.coverage.merged_to_sources.guards\endcsname{28}
\expandafter\def\csname ccval@depth.coverage.merged_to_sources.guards_graded\endcsname{0}
\expandafter\def\csname ccval@depth.coverage.merged_to_sources.guards_wrong\endcsname{0}
\expandafter\def\csname ccval@depth.coverage.merged_to_sources.plants\endcsname{3}
\expandafter\def\csname ccval@depth.coverage.merged_to_sources.plants_detected\endcsname{3}
\expandafter\def\csname ccval@depth.coverage.merged_to_sources.plants_detected_mechanically\endcsname{3}
\expandafter\def\csname ccval@depth.coverage.model_calls\endcsname{26}
\expandafter\def\csname ccval@depth.coverage.plants_detected\endcsname{6}
\expandafter\def\csname ccval@depth.coverage.plants_total\endcsname{7}
\expandafter\def\csname ccval@depth.coverage.second_reports_total\endcsname{0}
\expandafter\def\csname ccval@depth.coverage.seconds\endcsname{1432.8}
\expandafter\def\csname ccval@depth.coverage.source_to_merged.guards\endcsname{90}
\expandafter\def\csname ccval@depth.coverage.source_to_merged.guards_graded\endcsname{90}
\expandafter\def\csname ccval@depth.coverage.source_to_merged.guards_wrong\endcsname{0}
\expandafter\def\csname ccval@depth.coverage.source_to_merged.plants\endcsname{3}
\expandafter\def\csname ccval@depth.coverage.source_to_merged.plants_detected\endcsname{3}
\expandafter\def\csname ccval@depth.coverage.source_to_merged.plants_detected_mechanically\endcsname{0}
\expandafter\def\csname ccval@depth.coverage.speedup\endcsname{2.04}
\expandafter\def\csname ccval@depth.draws\endcsname{1}
\expandafter\def\csname ccval@depth.fixtures\endcsname{13}
\expandafter\def\csname ccval@depth.full.as_run.disqualified\endcsname{1}
\expandafter\def\csname ccval@depth.full.as_run.guard_wrong_total\endcsname{0}
\expandafter\def\csname ccval@depth.full.as_run.invented_total\endcsname{2}
\expandafter\def\csname ccval@depth.full.as_run.plants_detected\endcsname{7}
\expandafter\def\csname ccval@depth.full.disqualified\endcsname{0}
\expandafter\def\csname ccval@depth.full.fixtures\endcsname{13}
\expandafter\def\csname ccval@depth.full.guard_confirmed_total\endcsname{116}
\expandafter\def\csname ccval@depth.full.guard_probes_total\endcsname{118}
\expandafter\def\csname ccval@depth.full.guard_unclaimed_total\endcsname{2}
\expandafter\def\csname ccval@depth.full.guard_wrong_total\endcsname{0}
\expandafter\def\csname ccval@depth.full.invented_total\endcsname{0}
\expandafter\def\csname ccval@depth.full.merged_to_sources.guards\endcsname{28}
\expandafter\def\csname ccval@depth.full.merged_to_sources.guards_graded\endcsname{28}
\expandafter\def\csname ccval@depth.full.merged_to_sources.guards_wrong\endcsname{0}
\expandafter\def\csname ccval@depth.full.merged_to_sources.plants\endcsname{4}
\expandafter\def\csname ccval@depth.full.merged_to_sources.plants_detected\endcsname{4}
\expandafter\def\csname ccval@depth.full.merged_to_sources.plants_detected_mechanically\endcsname{1}
\expandafter\def\csname ccval@depth.full.model_calls\endcsname{71}
\expandafter\def\csname ccval@depth.full.plants_detected\endcsname{7}
\expandafter\def\csname ccval@depth.full.plants_total\endcsname{7}
\expandafter\def\csname ccval@depth.full.second_reports_total\endcsname{2}
\expandafter\def\csname ccval@depth.full.seconds\endcsname{2924.2}
\expandafter\def\csname ccval@depth.full.source_to_merged.guards\endcsname{90}
\expandafter\def\csname ccval@depth.full.source_to_merged.guards_graded\endcsname{88}
\expandafter\def\csname ccval@depth.full.source_to_merged.guards_wrong\endcsname{0}
\expandafter\def\csname ccval@depth.full.source_to_merged.plants\endcsname{3}
\expandafter\def\csname ccval@depth.full.source_to_merged.plants_detected\endcsname{3}
\expandafter\def\csname ccval@depth.full.source_to_merged.plants_detected_mechanically\endcsname{0}
\expandafter\def\csname ccval@depth.measured_on\endcsname{2026-09-24}
\expandafter\def\csname ccval@depth.model\endcsname{Qwen/Qwen3.8-27B-FP8}
\expandafter\def\csname ccval@depth.plants\endcsname{7}
\expandafter\def\csname ccval@detect.A-27b.exit_code_matched\endcsname{11}
\expandafter\def\csname ccval@detect.A-27b.exit_code_missed\endcsname{2}
\expandafter\def\csname ccval@detect.A-27b.fixtures\endcsname{13}
\expandafter\def\csname ccval@detect.A-27b.invented_total\endcsname{2}
\expandafter\def\csname ccval@detect.A-27b.model\endcsname{qwen3.8:27b}
\expandafter\def\csname ccval@detect.A-27b.plants_detected\endcsname{6}
\expandafter\def\csname ccval@detect.A-27b.plants_missed\endcsname{0}
\expandafter\def\csname ccval@detect.A-27b.plants_total\endcsname{6}
\expandafter\def\csname ccval@detect.B-70b.exit_code_matched\endcsname{11}
\expandafter\def\csname ccval@detect.B-70b.exit_code_missed\endcsname{2}
\expandafter\def\csname ccval@detect.B-70b.fixtures\endcsname{13}
\expandafter\def\csname ccval@detect.B-70b.invented_total\endcsname{1}
\expandafter\def\csname ccval@detect.B-70b.model\endcsname{deepseek-r1:70b}
\expandafter\def\csname ccval@detect.B-70b.plants_detected\endcsname{5}
\expandafter\def\csname ccval@detect.B-70b.plants_missed\endcsname{1}
\expandafter\def\csname ccval@detect.B-70b.plants_total\endcsname{6}
\expandafter\def\csname ccval@detect.C-120b.exit_code_matched\endcsname{13}
\expandafter\def\csname ccval@detect.C-120b.exit_code_missed\endcsname{0}
\expandafter\def\csname ccval@detect.C-120b.fixtures\endcsname{13}
\expandafter\def\csname ccval@detect.C-120b.invented_total\endcsname{3}
\expandafter\def\csname ccval@detect.C-120b.model\endcsname{gpt-oss:120b}
\expandafter\def\csname ccval@detect.C-120b.plants_detected\endcsname{6}
\expandafter\def\csname ccval@detect.C-120b.plants_missed\endcsname{0}
\expandafter\def\csname ccval@detect.C-120b.plants_total\endcsname{6}
\expandafter\def\csname ccval@fd.departed_records\endcsname{29}
\expandafter\def\csname ccval@fd.excluded\endcsname{2}
\expandafter\def\csname ccval@fd.findings\endcsname{15}
\expandafter\def\csname ccval@fd.records\endcsname{30}
\expandafter\def\csname ccval@fd.segments\endcsname{63}
\expandafter\def\csname ccval@lineup.amended.date\endcsname{2026-10-12}
\expandafter\def\csname ccval@lineup.amendments.after_results\endcsname{4}
\expandafter\def\csname ccval@lineup.amendments.before_calls\endcsname{2}
\expandafter\def\csname ccval@lineup.amendments.last\endcsname{6}
\expandafter\def\csname ccval@lineup.amendments.last_date\endcsname{2026-10-15}
\expandafter\def\csname ccval@lineup.audit.all.cases\endcsname{65}
\expandafter\def\csname ccval@lineup.audit.all.gone\endcsname{2}
\expandafter\def\csname ccval@lineup.audit.all.headings_untaken\endcsname{12}
\expandafter\def\csname ccval@lineup.audit.all.lost\endcsname{11}
\expandafter\def\csname ccval@lineup.audit.all.lost_cleared\endcsname{0}
\expandafter\def\csname ccval@lineup.audit.all.lost_still_counted\endcsname{11}
\expandafter\def\csname ccval@lineup.audit.all.not_content\endcsname{12}
\expandafter\def\csname ccval@lineup.audit.all.partly\endcsname{9}
\expandafter\def\csname ccval@lineup.audit.all.present\endcsname{42}
\expandafter\def\csname ccval@lineup.audit.all.present_found\endcsname{8}
\expandafter\def\csname ccval@lineup.audit.all.present_still_counted\endcsname{34}
\expandafter\def\csname ccval@lineup.audit.all.still_counted\endcsname{45}
\expandafter\def\csname ccval@lineup.audit.all.still_there\endcsname{54}
\expandafter\def\csname ccval@lineup.audit.borderline.cases\endcsname{12}
\expandafter\def\csname ccval@lineup.audit.borderline.first\endcsname{4}
\expandafter\def\csname ccval@lineup.audit.borderline.verdict_row\endcsname{gpt-6-luna-api}
\expandafter\def\csname ccval@lineup.audit.borderline.verdicts_moved\endcsname{1}
\expandafter\def\csname ccval@lineup.audit.both.borderline.cases\endcsname{19}
\expandafter\def\csname ccval@lineup.audit.both.borderline.first\endcsname{11}
\expandafter\def\csname ccval@lineup.audit.both.cases\endcsname{84}
\expandafter\def\csname ccval@lineup.audit.both.first.cases\endcsname{44}
\expandafter\def\csname ccval@lineup.audit.both.first.gone\endcsname{1}
\expandafter\def\csname ccval@lineup.audit.both.first.lost\endcsname{6}
\expandafter\def\csname ccval@lineup.audit.both.first.not_content\endcsname{12}
\expandafter\def\csname ccval@lineup.audit.both.first.partly\endcsname{5}
\expandafter\def\csname ccval@lineup.audit.both.first.present\endcsname{26}
\expandafter\def\csname ccval@lineup.audit.both.first.still_there\endcsname{38}
\expandafter\def\csname ccval@lineup.audit.both.runs\endcsname{21}
\expandafter\def\csname ccval@lineup.audit.both.second.cases\endcsname{44}
\expandafter\def\csname ccval@lineup.audit.both.second.same_class\endcsname{42}
\expandafter\def\csname ccval@lineup.audit.date\endcsname{2026-10-12}
\expandafter\def\csname ccval@lineup.audit.example.author_api.class\endcsname{partly}
\expandafter\def\csname ccval@lineup.audit.example.author_api.pair\endcsname{trace\_names}
\expandafter\def\csname ccval@lineup.audit.example.author_api.readings_naming_present\endcsname{0}
\expandafter\def\csname ccval@lineup.audit.example.author_api.row\endcsname{haiku-4.5-api}
\expandafter\def\csname ccval@lineup.audit.example.author_api.segment\endcsname{b2}
\expandafter\def\csname ccval@lineup.audit.example.author_sub.class\endcsname{gone}
\expandafter\def\csname ccval@lineup.audit.example.author_sub.pair\endcsname{trace\_names}
\expandafter\def\csname ccval@lineup.audit.example.author_sub.readings_naming_present\endcsname{0}
\expandafter\def\csname ccval@lineup.audit.example.author_sub.row\endcsname{haiku-4.5-sub}
\expandafter\def\csname ccval@lineup.audit.example.author_sub.segment\endcsname{b2}
\expandafter\def\csname ccval@lineup.audit.example.condition.class\endcsname{partly}
\expandafter\def\csname ccval@lineup.audit.example.condition.heading\endcsname{Placing an item on the shelf}
\expandafter\def\csname ccval@lineup.audit.example.condition.merge\endcsname{Print a hold slip and fold it into the item so the name is visible from the spine.}
\expandafter\def\csname ccval@lineup.audit.example.condition.missing\endcsname{When a requested item arrives}
\expandafter\def\csname ccval@lineup.audit.example.condition.pair\endcsname{library\_holds}
\expandafter\def\csname ccval@lineup.audit.example.condition.readings_naming_present\endcsname{2}
\expandafter\def\csname ccval@lineup.audit.example.condition.row\endcsname{gpt-6-luna-api}
\expandafter\def\csname ccval@lineup.audit.example.condition.segment\endcsname{a3}
\expandafter\def\csname ccval@lineup.audit.example.condition.source\endcsname{When a requested item arrives, print a hold slip and fold it into the item so the name is visible from the spine.}
\expandafter\def\csname ccval@lineup.audit.example.heading.merge\endcsname{How to retry well}
\expandafter\def\csname ccval@lineup.audit.example.heading.pair\endcsname{rate\_limits}
\expandafter\def\csname ccval@lineup.audit.example.heading.row\endcsname{opus-5.5-api}
\expandafter\def\csname ccval@lineup.audit.example.heading.segment\endcsname{b20}
\expandafter\def\csname ccval@lineup.audit.example.heading.source\endcsname{What the client does}
\expandafter\def\csname ccval@lineup.audit.example.joined.merge\endcsname{Escalate E2 to the workshop after 2 failed attempts; do not keep cycling a dock that has failed twice, since repeated forcing bends the pin and turns a 20 minute workshop job into a replacement.}
\expandafter\def\csname ccval@lineup.audit.example.joined.pair\endcsname{bike\_docks}
\expandafter\def\csname ccval@lineup.audit.example.joined.row\endcsname{sonnet-5-sub}
\expandafter\def\csname ccval@lineup.audit.example.joined.segment\endcsname{a10}
\expandafter\def\csname ccval@lineup.audit.example.joined.source\endcsname{Escalate E2 to the workshop after 2 failed attempts.}
\expandafter\def\csname ccval@lineup.audit.example.partly.merge\endcsname{A station with more than half its docks faulted is taken out of service in the app. Notify the support desk before this happens, not after.}
\expandafter\def\csname ccval@lineup.audit.example.partly.pair\endcsname{bike\_docks}
\expandafter\def\csname ccval@lineup.audit.example.partly.row\endcsname{haiku-4.5-sub}
\expandafter\def\csname ccval@lineup.audit.example.partly.segment\endcsname{b18}
\expandafter\def\csname ccval@lineup.audit.example.partly.source\endcsname{The support desk takes a station out of service in the app when more than half its docks are faulted.}
\expandafter\def\csname ccval@lineup.audit.example.shortened.merge\endcsname{The count follows the credential across hosts.}
\expandafter\def\csname ccval@lineup.audit.example.shortened.pair\endcsname{rate\_limits}
\expandafter\def\csname ccval@lineup.audit.example.shortened.row\endcsname{gpt-6-sol-api}
\expandafter\def\csname ccval@lineup.audit.example.shortened.segment\endcsname{b7}
\expandafter\def\csname ccval@lineup.audit.example.shortened.source\endcsname{The count follows the credential wherever the caller happens to put it, including across hosts.}
\expandafter\def\csname ccval@lineup.audit.extended.date\endcsname{2026-10-12}
\expandafter\def\csname ccval@lineup.audit.first.cases\endcsname{25}
\expandafter\def\csname ccval@lineup.audit.first.gone\endcsname{0}
\expandafter\def\csname ccval@lineup.audit.first.headings_untaken\endcsname{11}
\expandafter\def\csname ccval@lineup.audit.first.lost\endcsname{3}
\expandafter\def\csname ccval@lineup.audit.first.lost_cleared\endcsname{0}
\expandafter\def\csname ccval@lineup.audit.first.lost_still_counted\endcsname{3}
\expandafter\def\csname ccval@lineup.audit.first.not_content\endcsname{11}
\expandafter\def\csname ccval@lineup.audit.first.partly\endcsname{3}
\expandafter\def\csname ccval@lineup.audit.first.present\endcsname{11}
\expandafter\def\csname ccval@lineup.audit.first.present_found\endcsname{4}
\expandafter\def\csname ccval@lineup.audit.first.present_still_counted\endcsname{7}
\expandafter\def\csname ccval@lineup.audit.first.still_counted\endcsname{10}
\expandafter\def\csname ccval@lineup.audit.first.still_there\endcsname{22}
\expandafter\def\csname ccval@lineup.audit.readings\endcsname{4}
\expandafter\def\csname ccval@lineup.audit.reworded.cases\endcsname{19}
\expandafter\def\csname ccval@lineup.audit.reworded.differ\endcsname{0}
\expandafter\def\csname ccval@lineup.audit.reworded.gone\endcsname{1}
\expandafter\def\csname ccval@lineup.audit.reworded.lost\endcsname{3}
\expandafter\def\csname ccval@lineup.audit.reworded.lost_reported\endcsname{3}
\expandafter\def\csname ccval@lineup.audit.reworded.not_content\endcsname{1}
\expandafter\def\csname ccval@lineup.audit.reworded.partly\endcsname{2}
\expandafter\def\csname ccval@lineup.audit.reworded.present\endcsname{15}
\expandafter\def\csname ccval@lineup.audit.reworded.reported\endcsname{19}
\expandafter\def\csname ccval@lineup.audit.reworded.rows\endcsname{4}
\expandafter\def\csname ccval@lineup.audit.reworded.rows_none\endcsname{4}
\expandafter\def\csname ccval@lineup.audit.reworded.runs\endcsname{12}
\expandafter\def\csname ccval@lineup.audit.reworded.same_class\endcsname{19}
\expandafter\def\csname ccval@lineup.audit.rows_disagree\endcsname{0}
\expandafter\def\csname ccval@lineup.audit.rows_lost\endcsname{3}
\expandafter\def\csname ccval@lineup.audit.rows_read\endcsname{5}
\expandafter\def\csname ccval@lineup.audit.runs\endcsname{18}
\expandafter\def\csname ccval@lineup.audit.second.cases\endcsname{25}
\expandafter\def\csname ccval@lineup.audit.second.differ\endcsname{2}
\expandafter\def\csname ccval@lineup.audit.second.lost\endcsname{5}
\expandafter\def\csname ccval@lineup.audit.second.same_class\endcsname{23}
\expandafter\def\csname ccval@lineup.audit.unread.as_run\endcsname{5}
\expandafter\def\csname ccval@lineup.audit.unread.repeats\endcsname{30}
\expandafter\def\csname ccval@lineup.audit.unread.repeats.rows\endcsname{2}
\expandafter\def\csname ccval@lineup.audit.unread.titles\endcsname{2}
\expandafter\def\csname ccval@lineup.cells\endcsname{109}
\expandafter\def\csname ccval@lineup.disqualified\endcsname{3}
\expandafter\def\csname ccval@lineup.fable-sub.S3.bip39.draws\endcsname{0}
\expandafter\def\csname ccval@lineup.fable-sub.S3.mahjongg.draws\endcsname{0}
\expandafter\def\csname ccval@lineup.fable-sub.S3.voyager.draws\endcsname{1}
\expandafter\def\csname ccval@lineup.fable-sub.S3.voyager.max\endcsname{38}
\expandafter\def\csname ccval@lineup.fable-sub.S3.voyager.median\endcsname{38}
\expandafter\def\csname ccval@lineup.fable-sub.S3.voyager.min\endcsname{38}
\expandafter\def\csname ccval@lineup.fable-sub.S3.voyager.total\endcsname{44}
\expandafter\def\csname ccval@lineup.fable-sub.effort\endcsname{medium}
\expandafter\def\csname ccval@lineup.fable-sub.name\endcsname{fable-sub}
\expandafter\def\csname ccval@lineup.failure1.role\endcsname{verify}
\expandafter\def\csname ccval@lineup.failure1.row\endcsname{haiku-4.5-api}
\expandafter\def\csname ccval@lineup.failure2.role\endcsname{merge}
\expandafter\def\csname ccval@lineup.failure2.row\endcsname{haiku-4.5-sub}
\expandafter\def\csname ccval@lineup.first.silent_segments\endcsname{10}
\expandafter\def\csname ccval@lineup.first.untaken\endcsname{11}
\expandafter\def\csname ccval@lineup.fixtures\endcsname{13}
\expandafter\def\csname ccval@lineup.flagged_not_disqualified\endcsname{1}
\expandafter\def\csname ccval@lineup.flagged_rows\endcsname{4}
\expandafter\def\csname ccval@lineup.floor.base_only.dev_per_pair\endcsname{7.56}
\expandafter\def\csname ccval@lineup.floor.base_only.deviations\endcsname{68}
\expandafter\def\csname ccval@lineup.floor.base_only.pairs\endcsname{9}
\expandafter\def\csname ccval@lineup.floor.base_only.silent\endcsname{121}
\expandafter\def\csname ccval@lineup.floor.base_only.silent_per_pair\endcsname{13.44}
\expandafter\def\csname ccval@lineup.floor.base_only.untaken\endcsname{8}
\expandafter\def\csname ccval@lineup.floor.concatenation.dev_per_pair\endcsname{8.78}
\expandafter\def\csname ccval@lineup.floor.concatenation.deviations\endcsname{79}
\expandafter\def\csname ccval@lineup.floor.concatenation.pairs\endcsname{9}
\expandafter\def\csname ccval@lineup.floor.concatenation.silent\endcsname{0}
\expandafter\def\csname ccval@lineup.floor.concatenation.silent_per_pair\endcsname{0.00}
\expandafter\def\csname ccval@lineup.floor.concatenation.untaken\endcsname{0}
\expandafter\def\csname ccval@lineup.floor.union.dev_per_pair\endcsname{7.78}
\expandafter\def\csname ccval@lineup.floor.union.deviations\endcsname{70}
\expandafter\def\csname ccval@lineup.floor.union.pairs\endcsname{9}
\expandafter\def\csname ccval@lineup.floor.union.silent\endcsname{0}
\expandafter\def\csname ccval@lineup.floor.union.silent_per_pair\endcsname{0.00}
\expandafter\def\csname ccval@lineup.floor.union.untaken\endcsname{0}
\expandafter\def\csname ccval@lineup.from\endcsname{2026-09-27}
\expandafter\def\csname ccval@lineup.gpt-6-luna-api.S2.disqualified\endcsname{3}
\expandafter\def\csname ccval@lineup.gpt-6-luna-api.S2.fixtures\endcsname{13}
\expandafter\def\csname ccval@lineup.gpt-6-luna-api.S2.guards_wrong\endcsname{1}
\expandafter\def\csname ccval@lineup.gpt-6-luna-api.S2.invented\endcsname{4}
\expandafter\def\csname ccval@lineup.gpt-6-luna-api.S2.matched\endcsname{10}
\expandafter\def\csname ccval@lineup.gpt-6-luna-api.S2.measured\endcsname{12}
\expandafter\def\csname ccval@lineup.gpt-6-luna-api.S2.order\endcsname{6}
\expandafter\def\csname ccval@lineup.gpt-6-luna-api.S2.plants\endcsname{6}
\expandafter\def\csname ccval@lineup.gpt-6-luna-api.S2.plants_measured\endcsname{6}
\expandafter\def\csname ccval@lineup.gpt-6-luna-api.S2.unmeasured\endcsname{1}
\expandafter\def\csname ccval@lineup.gpt-6-luna-api.S3.bip39.draws\endcsname{3}
\expandafter\def\csname ccval@lineup.gpt-6-luna-api.S3.bip39.max\endcsname{10}
\expandafter\def\csname ccval@lineup.gpt-6-luna-api.S3.bip39.median\endcsname{8}
\expandafter\def\csname ccval@lineup.gpt-6-luna-api.S3.bip39.min\endcsname{7}
\expandafter\def\csname ccval@lineup.gpt-6-luna-api.S3.bip39.total\endcsname{14}
\expandafter\def\csname ccval@lineup.gpt-6-luna-api.S3.fixed_total\endcsname{27}
\expandafter\def\csname ccval@lineup.gpt-6-luna-api.S3.mahjongg.draws\endcsname{1}
\expandafter\def\csname ccval@lineup.gpt-6-luna-api.S3.mahjongg.false\endcsname{0}
\expandafter\def\csname ccval@lineup.gpt-6-luna-api.S3.order\endcsname{3=}
\expandafter\def\csname ccval@lineup.gpt-6-luna-api.S3.voyager.draws\endcsname{3}
\expandafter\def\csname ccval@lineup.gpt-6-luna-api.S3.voyager.max\endcsname{19}
\expandafter\def\csname ccval@lineup.gpt-6-luna-api.S3.voyager.median\endcsname{19}
\expandafter\def\csname ccval@lineup.gpt-6-luna-api.S3.voyager.min\endcsname{18}
\expandafter\def\csname ccval@lineup.gpt-6-luna-api.S3.voyager.total\endcsname{44}
\expandafter\def\csname ccval@lineup.gpt-6-luna-api.audit.all.cases\endcsname{5}
\expandafter\def\csname ccval@lineup.gpt-6-luna-api.audit.all.gone\endcsname{0}
\expandafter\def\csname ccval@lineup.gpt-6-luna-api.audit.all.lost\endcsname{0}
\expandafter\def\csname ccval@lineup.gpt-6-luna-api.audit.all.not_content\endcsname{0}
\expandafter\def\csname ccval@lineup.gpt-6-luna-api.audit.all.partly\endcsname{0}
\expandafter\def\csname ccval@lineup.gpt-6-luna-api.audit.all.present\endcsname{5}
\expandafter\def\csname ccval@lineup.gpt-6-luna-api.audit.all.runs\endcsname{3}
\expandafter\def\csname ccval@lineup.gpt-6-luna-api.audit.borderline_only\endcsname{yes}
\expandafter\def\csname ccval@lineup.gpt-6-luna-api.audit.cases\endcsname{4}
\expandafter\def\csname ccval@lineup.gpt-6-luna-api.audit.counted\endcsname{2}
\expandafter\def\csname ccval@lineup.gpt-6-luna-api.audit.counted.lost\endcsname{0}
\expandafter\def\csname ccval@lineup.gpt-6-luna-api.audit.disqualified\endcsname{yes}
\expandafter\def\csname ccval@lineup.gpt-6-luna-api.audit.found\endcsname{1 of 4}
\expandafter\def\csname ccval@lineup.gpt-6-luna-api.audit.gone\endcsname{0}
\expandafter\def\csname ccval@lineup.gpt-6-luna-api.audit.lost\endcsname{1}
\expandafter\def\csname ccval@lineup.gpt-6-luna-api.audit.lost_borderline\endcsname{1}
\expandafter\def\csname ccval@lineup.gpt-6-luna-api.audit.not_content\endcsname{0}
\expandafter\def\csname ccval@lineup.gpt-6-luna-api.audit.partly\endcsname{1}
\expandafter\def\csname ccval@lineup.gpt-6-luna-api.audit.present\endcsname{3}
\expandafter\def\csname ccval@lineup.gpt-6-luna-api.audit.reworded\endcsname{2}
\expandafter\def\csname ccval@lineup.gpt-6-luna-api.audit.reworded.lost\endcsname{1}
\expandafter\def\csname ccval@lineup.gpt-6-luna-api.audit.second.lost\endcsname{1}
\expandafter\def\csname ccval@lineup.gpt-6-luna-api.dev_per_pair\endcsname{5.44}
\expandafter\def\csname ccval@lineup.gpt-6-luna-api.deviations\endcsname{49}
\expandafter\def\csname ccval@lineup.gpt-6-luna-api.effort\endcsname{medium}
\expandafter\def\csname ccval@lineup.gpt-6-luna-api.list_output_multiple\endcsname{1}
\expandafter\def\csname ccval@lineup.gpt-6-luna-api.model\endcsname{gpt-6-luna}
\expandafter\def\csname ccval@lineup.gpt-6-luna-api.model_failures\endcsname{0}
\expandafter\def\csname ccval@lineup.gpt-6-luna-api.name\endcsname{gpt-6-luna-api}
\expandafter\def\csname ccval@lineup.gpt-6-luna-api.no_better_than_union\endcsname{no}
\expandafter\def\csname ccval@lineup.gpt-6-luna-api.pairs\endcsname{9}
\expandafter\def\csname ccval@lineup.gpt-6-luna-api.pairs_completed\endcsname{9}
\expandafter\def\csname ccval@lineup.gpt-6-luna-api.place\endcsname{disq.}
\expandafter\def\csname ccval@lineup.gpt-6-luna-api.rank\endcsname{disq.}
\expandafter\def\csname ccval@lineup.gpt-6-luna-api.seconds\endcsname{141.3}
\expandafter\def\csname ccval@lineup.gpt-6-luna-api.silent\endcsname{2}
\expandafter\def\csname ccval@lineup.gpt-6-luna-api.silent_per_pair\endcsname{0.22}
\expandafter\def\csname ccval@lineup.gpt-6-luna-api.silent_runs\endcsname{1}
\expandafter\def\csname ccval@lineup.gpt-6-luna-api.spread.badge_access.deviations.max\endcsname{1}
\expandafter\def\csname ccval@lineup.gpt-6-luna-api.spread.badge_access.deviations.median\endcsname{1}
\expandafter\def\csname ccval@lineup.gpt-6-luna-api.spread.badge_access.deviations.min\endcsname{0}
\expandafter\def\csname ccval@lineup.gpt-6-luna-api.spread.badge_access.draws\endcsname{3}
\expandafter\def\csname ccval@lineup.gpt-6-luna-api.spread.badge_access.silent.max\endcsname{0}
\expandafter\def\csname ccval@lineup.gpt-6-luna-api.spread.badge_access.silent.median\endcsname{0}
\expandafter\def\csname ccval@lineup.gpt-6-luna-api.spread.badge_access.silent.min\endcsname{0}
\expandafter\def\csname ccval@lineup.gpt-6-luna-api.spread.rate_limits.deviations.max\endcsname{22}
\expandafter\def\csname ccval@lineup.gpt-6-luna-api.spread.rate_limits.deviations.median\endcsname{18}
\expandafter\def\csname ccval@lineup.gpt-6-luna-api.spread.rate_limits.deviations.min\endcsname{14}
\expandafter\def\csname ccval@lineup.gpt-6-luna-api.spread.rate_limits.draws\endcsname{3}
\expandafter\def\csname ccval@lineup.gpt-6-luna-api.spread.rate_limits.silent.max\endcsname{2}
\expandafter\def\csname ccval@lineup.gpt-6-luna-api.spread.rate_limits.silent.median\endcsname{0}
\expandafter\def\csname ccval@lineup.gpt-6-luna-api.spread.rate_limits.silent.min\endcsname{0}
\expandafter\def\csname ccval@lineup.gpt-6-luna-api.spread.trace_names.deviations.max\endcsname{9}
\expandafter\def\csname ccval@lineup.gpt-6-luna-api.spread.trace_names.deviations.median\endcsname{8}
\expandafter\def\csname ccval@lineup.gpt-6-luna-api.spread.trace_names.deviations.min\endcsname{8}
\expandafter\def\csname ccval@lineup.gpt-6-luna-api.spread.trace_names.draws\endcsname{3}
\expandafter\def\csname ccval@lineup.gpt-6-luna-api.spread.trace_names.silent.max\endcsname{0}
\expandafter\def\csname ccval@lineup.gpt-6-luna-api.spread.trace_names.silent.median\endcsname{0}
\expandafter\def\csname ccval@lineup.gpt-6-luna-api.spread.trace_names.silent.min\endcsname{0}
\expandafter\def\csname ccval@lineup.gpt-6-luna-api.spread_pairs\endcsname{3}
\expandafter\def\csname ccval@lineup.gpt-6-luna-api.story.exit\endcsname{1}
\expandafter\def\csname ccval@lineup.gpt-6-luna-api.story.flagged\endcsname{no}
\expandafter\def\csname ccval@lineup.gpt-6-luna-api.story.kept_a\endcsname{yes}
\expandafter\def\csname ccval@lineup.gpt-6-luna-api.story.kept_b\endcsname{no}
\expandafter\def\csname ccval@lineup.gpt-6-luna-api.tokens\endcsname{39666}
\expandafter\def\csname ccval@lineup.gpt-6-luna-api.union_dev\endcsname{70}
\expandafter\def\csname ccval@lineup.gpt-6-luna-api.unit\endcsname{list}
\expandafter\def\csname ccval@lineup.gpt-6-luna-api.untaken\endcsname{0}
\expandafter\def\csname ccval@lineup.gpt-6-luna-api.untaken_all\endcsname{0}
\expandafter\def\csname ccval@lineup.gpt-6-luna-api.untaken_runs\endcsname{0}
\expandafter\def\csname ccval@lineup.gpt-6-luna-api.usd\endcsname{0.012}
\expandafter\def\csname ccval@lineup.gpt-6-sol-api.S2.disqualified\endcsname{1}
\expandafter\def\csname ccval@lineup.gpt-6-sol-api.S2.fixtures\endcsname{13}
\expandafter\def\csname ccval@lineup.gpt-6-sol-api.S2.guards_wrong\endcsname{0}
\expandafter\def\csname ccval@lineup.gpt-6-sol-api.S2.invented\endcsname{1}
\expandafter\def\csname ccval@lineup.gpt-6-sol-api.S2.matched\endcsname{11}
\expandafter\def\csname ccval@lineup.gpt-6-sol-api.S2.measured\endcsname{12}
\expandafter\def\csname ccval@lineup.gpt-6-sol-api.S2.order\endcsname{4}
\expandafter\def\csname ccval@lineup.gpt-6-sol-api.S2.plants\endcsname{6}
\expandafter\def\csname ccval@lineup.gpt-6-sol-api.S2.plants_measured\endcsname{6}
\expandafter\def\csname ccval@lineup.gpt-6-sol-api.S2.unmeasured\endcsname{1}
\expandafter\def\csname ccval@lineup.gpt-6-sol-api.S3.bip39.draws\endcsname{3}
\expandafter\def\csname ccval@lineup.gpt-6-sol-api.S3.bip39.max\endcsname{12}
\expandafter\def\csname ccval@lineup.gpt-6-sol-api.S3.bip39.median\endcsname{11}
\expandafter\def\csname ccval@lineup.gpt-6-sol-api.S3.bip39.min\endcsname{10}
\expandafter\def\csname ccval@lineup.gpt-6-sol-api.S3.bip39.total\endcsname{14}
\expandafter\def\csname ccval@lineup.gpt-6-sol-api.S3.fixed_total\endcsname{34}
\expandafter\def\csname ccval@lineup.gpt-6-sol-api.S3.mahjongg.draws\endcsname{1}
\expandafter\def\csname ccval@lineup.gpt-6-sol-api.S3.mahjongg.false\endcsname{2}
\expandafter\def\csname ccval@lineup.gpt-6-sol-api.S3.order\endcsname{2=}
\expandafter\def\csname ccval@lineup.gpt-6-sol-api.S3.voyager.draws\endcsname{3}
\expandafter\def\csname ccval@lineup.gpt-6-sol-api.S3.voyager.max\endcsname{29}
\expandafter\def\csname ccval@lineup.gpt-6-sol-api.S3.voyager.median\endcsname{23}
\expandafter\def\csname ccval@lineup.gpt-6-sol-api.S3.voyager.min\endcsname{22}
\expandafter\def\csname ccval@lineup.gpt-6-sol-api.S3.voyager.total\endcsname{44}
\expandafter\def\csname ccval@lineup.gpt-6-sol-api.audit.all.cases\endcsname{1}
\expandafter\def\csname ccval@lineup.gpt-6-sol-api.audit.all.gone\endcsname{0}
\expandafter\def\csname ccval@lineup.gpt-6-sol-api.audit.all.lost\endcsname{0}
\expandafter\def\csname ccval@lineup.gpt-6-sol-api.audit.all.not_content\endcsname{0}
\expandafter\def\csname ccval@lineup.gpt-6-sol-api.audit.all.partly\endcsname{0}
\expandafter\def\csname ccval@lineup.gpt-6-sol-api.audit.all.present\endcsname{1}
\expandafter\def\csname ccval@lineup.gpt-6-sol-api.audit.all.runs\endcsname{1}
\expandafter\def\csname ccval@lineup.gpt-6-sol-api.audit.borderline_only\endcsname{no}
\expandafter\def\csname ccval@lineup.gpt-6-sol-api.audit.cases\endcsname{0}
\expandafter\def\csname ccval@lineup.gpt-6-sol-api.audit.counted\endcsname{0}
\expandafter\def\csname ccval@lineup.gpt-6-sol-api.audit.counted.lost\endcsname{0}
\expandafter\def\csname ccval@lineup.gpt-6-sol-api.audit.disqualified\endcsname{no}
\expandafter\def\csname ccval@lineup.gpt-6-sol-api.audit.found\endcsname{--}
\expandafter\def\csname ccval@lineup.gpt-6-sol-api.audit.gone\endcsname{0}
\expandafter\def\csname ccval@lineup.gpt-6-sol-api.audit.lost\endcsname{0}
\expandafter\def\csname ccval@lineup.gpt-6-sol-api.audit.lost_borderline\endcsname{0}
\expandafter\def\csname ccval@lineup.gpt-6-sol-api.audit.not_content\endcsname{0}
\expandafter\def\csname ccval@lineup.gpt-6-sol-api.audit.partly\endcsname{0}
\expandafter\def\csname ccval@lineup.gpt-6-sol-api.audit.present\endcsname{0}
\expandafter\def\csname ccval@lineup.gpt-6-sol-api.audit.reworded\endcsname{0}
\expandafter\def\csname ccval@lineup.gpt-6-sol-api.audit.reworded.lost\endcsname{0}
\expandafter\def\csname ccval@lineup.gpt-6-sol-api.dev_per_pair\endcsname{3.89}
\expandafter\def\csname ccval@lineup.gpt-6-sol-api.deviations\endcsname{35}
\expandafter\def\csname ccval@lineup.gpt-6-sol-api.effort\endcsname{medium}
\expandafter\def\csname ccval@lineup.gpt-6-sol-api.list_output_multiple\endcsname{20}
\expandafter\def\csname ccval@lineup.gpt-6-sol-api.model\endcsname{gpt-6-sol}
\expandafter\def\csname ccval@lineup.gpt-6-sol-api.model_failures\endcsname{0}
\expandafter\def\csname ccval@lineup.gpt-6-sol-api.name\endcsname{gpt-6-sol-api}
\expandafter\def\csname ccval@lineup.gpt-6-sol-api.no_better_than_union\endcsname{no}
\expandafter\def\csname ccval@lineup.gpt-6-sol-api.pairs\endcsname{9}
\expandafter\def\csname ccval@lineup.gpt-6-sol-api.pairs_completed\endcsname{9}
\expandafter\def\csname ccval@lineup.gpt-6-sol-api.place\endcsname{1}
\expandafter\def\csname ccval@lineup.gpt-6-sol-api.rank\endcsname{1}
\expandafter\def\csname ccval@lineup.gpt-6-sol-api.seconds\endcsname{158.4}
\expandafter\def\csname ccval@lineup.gpt-6-sol-api.silent\endcsname{0}
\expandafter\def\csname ccval@lineup.gpt-6-sol-api.silent_per_pair\endcsname{0.00}
\expandafter\def\csname ccval@lineup.gpt-6-sol-api.silent_runs\endcsname{1}
\expandafter\def\csname ccval@lineup.gpt-6-sol-api.spread.badge_access.deviations.max\endcsname{1}
\expandafter\def\csname ccval@lineup.gpt-6-sol-api.spread.badge_access.deviations.median\endcsname{0}
\expandafter\def\csname ccval@lineup.gpt-6-sol-api.spread.badge_access.deviations.min\endcsname{0}
\expandafter\def\csname ccval@lineup.gpt-6-sol-api.spread.badge_access.draws\endcsname{3}
\expandafter\def\csname ccval@lineup.gpt-6-sol-api.spread.badge_access.silent.max\endcsname{0}
\expandafter\def\csname ccval@lineup.gpt-6-sol-api.spread.badge_access.silent.median\endcsname{0}
\expandafter\def\csname ccval@lineup.gpt-6-sol-api.spread.badge_access.silent.min\endcsname{0}
\expandafter\def\csname ccval@lineup.gpt-6-sol-api.spread.rate_limits.deviations.max\endcsname{9}
\expandafter\def\csname ccval@lineup.gpt-6-sol-api.spread.rate_limits.deviations.median\endcsname{9}
\expandafter\def\csname ccval@lineup.gpt-6-sol-api.spread.rate_limits.deviations.min\endcsname{9}
\expandafter\def\csname ccval@lineup.gpt-6-sol-api.spread.rate_limits.draws\endcsname{3}
\expandafter\def\csname ccval@lineup.gpt-6-sol-api.spread.rate_limits.silent.max\endcsname{1}
\expandafter\def\csname ccval@lineup.gpt-6-sol-api.spread.rate_limits.silent.median\endcsname{0}
\expandafter\def\csname ccval@lineup.gpt-6-sol-api.spread.rate_limits.silent.min\endcsname{0}
\expandafter\def\csname ccval@lineup.gpt-6-sol-api.spread.trace_names.deviations.max\endcsname{10}
\expandafter\def\csname ccval@lineup.gpt-6-sol-api.spread.trace_names.deviations.median\endcsname{9}
\expandafter\def\csname ccval@lineup.gpt-6-sol-api.spread.trace_names.deviations.min\endcsname{8}
\expandafter\def\csname ccval@lineup.gpt-6-sol-api.spread.trace_names.draws\endcsname{3}
\expandafter\def\csname ccval@lineup.gpt-6-sol-api.spread.trace_names.silent.max\endcsname{0}
\expandafter\def\csname ccval@lineup.gpt-6-sol-api.spread.trace_names.silent.median\endcsname{0}
\expandafter\def\csname ccval@lineup.gpt-6-sol-api.spread.trace_names.silent.min\endcsname{0}
\expandafter\def\csname ccval@lineup.gpt-6-sol-api.spread_pairs\endcsname{3}
\expandafter\def\csname ccval@lineup.gpt-6-sol-api.story.exit\endcsname{1}
\expandafter\def\csname ccval@lineup.gpt-6-sol-api.story.flagged\endcsname{yes}
\expandafter\def\csname ccval@lineup.gpt-6-sol-api.story.kept_a\endcsname{yes}
\expandafter\def\csname ccval@lineup.gpt-6-sol-api.story.kept_b\endcsname{no}
\expandafter\def\csname ccval@lineup.gpt-6-sol-api.tokens\endcsname{34956}
\expandafter\def\csname ccval@lineup.gpt-6-sol-api.union_dev\endcsname{70}
\expandafter\def\csname ccval@lineup.gpt-6-sol-api.unit\endcsname{list}
\expandafter\def\csname ccval@lineup.gpt-6-sol-api.untaken\endcsname{0}
\expandafter\def\csname ccval@lineup.gpt-6-sol-api.untaken_all\endcsname{0}
\expandafter\def\csname ccval@lineup.gpt-6-sol-api.untaken_runs\endcsname{0}
\expandafter\def\csname ccval@lineup.gpt-6-sol-api.usd\endcsname{0.189}
\expandafter\def\csname ccval@lineup.haiku-4.5-api.S2.disqualified\endcsname{2}
\expandafter\def\csname ccval@lineup.haiku-4.5-api.S2.fixtures\endcsname{13}
\expandafter\def\csname ccval@lineup.haiku-4.5-api.S2.guards_wrong\endcsname{1}
\expandafter\def\csname ccval@lineup.haiku-4.5-api.S2.invented\endcsname{2}
\expandafter\def\csname ccval@lineup.haiku-4.5-api.S2.matched\endcsname{12}
\expandafter\def\csname ccval@lineup.haiku-4.5-api.S2.measured\endcsname{13}
\expandafter\def\csname ccval@lineup.haiku-4.5-api.S2.order\endcsname{3=}
\expandafter\def\csname ccval@lineup.haiku-4.5-api.S2.plants\endcsname{7}
\expandafter\def\csname ccval@lineup.haiku-4.5-api.S2.plants_measured\endcsname{7}
\expandafter\def\csname ccval@lineup.haiku-4.5-api.S2.unmeasured\endcsname{0}
\expandafter\def\csname ccval@lineup.haiku-4.5-api.S3.bip39.draws\endcsname{3}
\expandafter\def\csname ccval@lineup.haiku-4.5-api.S3.bip39.max\endcsname{9}
\expandafter\def\csname ccval@lineup.haiku-4.5-api.S3.bip39.median\endcsname{8}
\expandafter\def\csname ccval@lineup.haiku-4.5-api.S3.bip39.min\endcsname{8}
\expandafter\def\csname ccval@lineup.haiku-4.5-api.S3.bip39.total\endcsname{14}
\expandafter\def\csname ccval@lineup.haiku-4.5-api.S3.fixed_total\endcsname{19}
\expandafter\def\csname ccval@lineup.haiku-4.5-api.S3.mahjongg.draws\endcsname{1}
\expandafter\def\csname ccval@lineup.haiku-4.5-api.S3.mahjongg.false\endcsname{1}
\expandafter\def\csname ccval@lineup.haiku-4.5-api.S3.order\endcsname{3=}
\expandafter\def\csname ccval@lineup.haiku-4.5-api.S3.voyager.draws\endcsname{3}
\expandafter\def\csname ccval@lineup.haiku-4.5-api.S3.voyager.max\endcsname{18}
\expandafter\def\csname ccval@lineup.haiku-4.5-api.S3.voyager.median\endcsname{11}
\expandafter\def\csname ccval@lineup.haiku-4.5-api.S3.voyager.min\endcsname{6}
\expandafter\def\csname ccval@lineup.haiku-4.5-api.S3.voyager.total\endcsname{44}
\expandafter\def\csname ccval@lineup.haiku-4.5-api.audit.all.cases\endcsname{13}
\expandafter\def\csname ccval@lineup.haiku-4.5-api.audit.all.gone\endcsname{0}
\expandafter\def\csname ccval@lineup.haiku-4.5-api.audit.all.lost\endcsname{1}
\expandafter\def\csname ccval@lineup.haiku-4.5-api.audit.all.not_content\endcsname{9}
\expandafter\def\csname ccval@lineup.haiku-4.5-api.audit.all.partly\endcsname{1}
\expandafter\def\csname ccval@lineup.haiku-4.5-api.audit.all.present\endcsname{3}
\expandafter\def\csname ccval@lineup.haiku-4.5-api.audit.all.runs\endcsname{7}
\expandafter\def\csname ccval@lineup.haiku-4.5-api.audit.borderline_only\endcsname{no}
\expandafter\def\csname ccval@lineup.haiku-4.5-api.audit.cases\endcsname{21}
\expandafter\def\csname ccval@lineup.haiku-4.5-api.audit.counted\endcsname{13}
\expandafter\def\csname ccval@lineup.haiku-4.5-api.audit.counted.lost\endcsname{1}
\expandafter\def\csname ccval@lineup.haiku-4.5-api.audit.disqualified\endcsname{yes}
\expandafter\def\csname ccval@lineup.haiku-4.5-api.audit.found\endcsname{2 of 21}
\expandafter\def\csname ccval@lineup.haiku-4.5-api.audit.gone\endcsname{0}
\expandafter\def\csname ccval@lineup.haiku-4.5-api.audit.lost\endcsname{2}
\expandafter\def\csname ccval@lineup.haiku-4.5-api.audit.lost_borderline\endcsname{1}
\expandafter\def\csname ccval@lineup.haiku-4.5-api.audit.not_content\endcsname{10}
\expandafter\def\csname ccval@lineup.haiku-4.5-api.audit.partly\endcsname{2}
\expandafter\def\csname ccval@lineup.haiku-4.5-api.audit.present\endcsname{9}
\expandafter\def\csname ccval@lineup.haiku-4.5-api.audit.reworded\endcsname{8}
\expandafter\def\csname ccval@lineup.haiku-4.5-api.audit.reworded.lost\endcsname{1}
\expandafter\def\csname ccval@lineup.haiku-4.5-api.audit.second.lost\endcsname{3}
\expandafter\def\csname ccval@lineup.haiku-4.5-api.dev_per_pair\endcsname{6.25}
\expandafter\def\csname ccval@lineup.haiku-4.5-api.deviations\endcsname{50}
\expandafter\def\csname ccval@lineup.haiku-4.5-api.effort\endcsname{one}
\expandafter\def\csname ccval@lineup.haiku-4.5-api.list_output_multiple\endcsname{10}
\expandafter\def\csname ccval@lineup.haiku-4.5-api.model\endcsname{claude-haiku-4-5-20251001}
\expandafter\def\csname ccval@lineup.haiku-4.5-api.model_failures\endcsname{1}
\expandafter\def\csname ccval@lineup.haiku-4.5-api.name\endcsname{haiku-4.5-api}
\expandafter\def\csname ccval@lineup.haiku-4.5-api.no_better_than_union\endcsname{no}
\expandafter\def\csname ccval@lineup.haiku-4.5-api.pairs\endcsname{8}
\expandafter\def\csname ccval@lineup.haiku-4.5-api.pairs_completed\endcsname{8}
\expandafter\def\csname ccval@lineup.haiku-4.5-api.place\endcsname{disq.}
\expandafter\def\csname ccval@lineup.haiku-4.5-api.rank\endcsname{disq.}
\expandafter\def\csname ccval@lineup.haiku-4.5-api.seconds\endcsname{69.1}
\expandafter\def\csname ccval@lineup.haiku-4.5-api.silent\endcsname{3}
\expandafter\def\csname ccval@lineup.haiku-4.5-api.silent_per_pair\endcsname{0.38}
\expandafter\def\csname ccval@lineup.haiku-4.5-api.silent_runs\endcsname{3}
\expandafter\def\csname ccval@lineup.haiku-4.5-api.spread_pairs\endcsname{0}
\expandafter\def\csname ccval@lineup.haiku-4.5-api.story.exit\endcsname{2}
\expandafter\def\csname ccval@lineup.haiku-4.5-api.story.flagged\endcsname{no}
\expandafter\def\csname ccval@lineup.haiku-4.5-api.story.kept_a\endcsname{yes}
\expandafter\def\csname ccval@lineup.haiku-4.5-api.story.kept_b\endcsname{yes}
\expandafter\def\csname ccval@lineup.haiku-4.5-api.tokens\endcsname{30132}
\expandafter\def\csname ccval@lineup.haiku-4.5-api.union_dev\endcsname{64}
\expandafter\def\csname ccval@lineup.haiku-4.5-api.unit\endcsname{list}
\expandafter\def\csname ccval@lineup.haiku-4.5-api.untaken\endcsname{9}
\expandafter\def\csname ccval@lineup.haiku-4.5-api.untaken_all\endcsname{9}
\expandafter\def\csname ccval@lineup.haiku-4.5-api.untaken_runs\endcsname{5}
\expandafter\def\csname ccval@lineup.haiku-4.5-api.usd\endcsname{0.069}
\expandafter\def\csname ccval@lineup.haiku-4.5-sub.S2.disqualified\endcsname{2}
\expandafter\def\csname ccval@lineup.haiku-4.5-sub.S2.fixtures\endcsname{13}
\expandafter\def\csname ccval@lineup.haiku-4.5-sub.S2.guards_wrong\endcsname{1}
\expandafter\def\csname ccval@lineup.haiku-4.5-sub.S2.invented\endcsname{1}
\expandafter\def\csname ccval@lineup.haiku-4.5-sub.S2.matched\endcsname{12}
\expandafter\def\csname ccval@lineup.haiku-4.5-sub.S2.measured\endcsname{13}
\expandafter\def\csname ccval@lineup.haiku-4.5-sub.S2.order\endcsname{2}
\expandafter\def\csname ccval@lineup.haiku-4.5-sub.S2.plants\endcsname{7}
\expandafter\def\csname ccval@lineup.haiku-4.5-sub.S2.plants_measured\endcsname{7}
\expandafter\def\csname ccval@lineup.haiku-4.5-sub.S2.unmeasured\endcsname{0}
\expandafter\def\csname ccval@lineup.haiku-4.5-sub.S3.bip39.draws\endcsname{3}
\expandafter\def\csname ccval@lineup.haiku-4.5-sub.S3.bip39.max\endcsname{9}
\expandafter\def\csname ccval@lineup.haiku-4.5-sub.S3.bip39.median\endcsname{8}
\expandafter\def\csname ccval@lineup.haiku-4.5-sub.S3.bip39.min\endcsname{6}
\expandafter\def\csname ccval@lineup.haiku-4.5-sub.S3.bip39.total\endcsname{14}
\expandafter\def\csname ccval@lineup.haiku-4.5-sub.S3.fixed_total\endcsname{14}
\expandafter\def\csname ccval@lineup.haiku-4.5-sub.S3.mahjongg.draws\endcsname{0}
\expandafter\def\csname ccval@lineup.haiku-4.5-sub.S3.mahjongg.false\endcsname{failed}
\expandafter\def\csname ccval@lineup.haiku-4.5-sub.S3.order\endcsname{3=}
\expandafter\def\csname ccval@lineup.haiku-4.5-sub.S3.voyager.draws\endcsname{3}
\expandafter\def\csname ccval@lineup.haiku-4.5-sub.S3.voyager.max\endcsname{17}
\expandafter\def\csname ccval@lineup.haiku-4.5-sub.S3.voyager.median\endcsname{6}
\expandafter\def\csname ccval@lineup.haiku-4.5-sub.S3.voyager.min\endcsname{2}
\expandafter\def\csname ccval@lineup.haiku-4.5-sub.S3.voyager.total\endcsname{44}
\expandafter\def\csname ccval@lineup.haiku-4.5-sub.audit.all.cases\endcsname{39}
\expandafter\def\csname ccval@lineup.haiku-4.5-sub.audit.all.gone\endcsname{2}
\expandafter\def\csname ccval@lineup.haiku-4.5-sub.audit.all.lost\endcsname{10}
\expandafter\def\csname ccval@lineup.haiku-4.5-sub.audit.all.not_content\endcsname{1}
\expandafter\def\csname ccval@lineup.haiku-4.5-sub.audit.all.partly\endcsname{8}
\expandafter\def\csname ccval@lineup.haiku-4.5-sub.audit.all.present\endcsname{28}
\expandafter\def\csname ccval@lineup.haiku-4.5-sub.audit.all.runs\endcsname{4}
\expandafter\def\csname ccval@lineup.haiku-4.5-sub.audit.borderline_only\endcsname{no}
\expandafter\def\csname ccval@lineup.haiku-4.5-sub.audit.cases\endcsname{8}
\expandafter\def\csname ccval@lineup.haiku-4.5-sub.audit.counted\endcsname{3}
\expandafter\def\csname ccval@lineup.haiku-4.5-sub.audit.counted.lost\endcsname{2}
\expandafter\def\csname ccval@lineup.haiku-4.5-sub.audit.disqualified\endcsname{yes}
\expandafter\def\csname ccval@lineup.haiku-4.5-sub.audit.found\endcsname{3 of 8}
\expandafter\def\csname ccval@lineup.haiku-4.5-sub.audit.gone\endcsname{1}
\expandafter\def\csname ccval@lineup.haiku-4.5-sub.audit.lost\endcsname{3}
\expandafter\def\csname ccval@lineup.haiku-4.5-sub.audit.lost_borderline\endcsname{1}
\expandafter\def\csname ccval@lineup.haiku-4.5-sub.audit.not_content\endcsname{0}
\expandafter\def\csname ccval@lineup.haiku-4.5-sub.audit.partly\endcsname{2}
\expandafter\def\csname ccval@lineup.haiku-4.5-sub.audit.present\endcsname{5}
\expandafter\def\csname ccval@lineup.haiku-4.5-sub.audit.reworded\endcsname{5}
\expandafter\def\csname ccval@lineup.haiku-4.5-sub.audit.reworded.lost\endcsname{1}
\expandafter\def\csname ccval@lineup.haiku-4.5-sub.audit.second.lost\endcsname{4}
\expandafter\def\csname ccval@lineup.haiku-4.5-sub.dev_per_pair\endcsname{7.11}
\expandafter\def\csname ccval@lineup.haiku-4.5-sub.deviations\endcsname{64}
\expandafter\def\csname ccval@lineup.haiku-4.5-sub.model\endcsname{claude-haiku-4-5-20251001}
\expandafter\def\csname ccval@lineup.haiku-4.5-sub.model_failures\endcsname{0}
\expandafter\def\csname ccval@lineup.haiku-4.5-sub.name\endcsname{haiku-4.5-sub}
\expandafter\def\csname ccval@lineup.haiku-4.5-sub.no_better_than_union\endcsname{no}
\expandafter\def\csname ccval@lineup.haiku-4.5-sub.pairs\endcsname{9}
\expandafter\def\csname ccval@lineup.haiku-4.5-sub.pairs_completed\endcsname{9}
\expandafter\def\csname ccval@lineup.haiku-4.5-sub.place\endcsname{disq.}
\expandafter\def\csname ccval@lineup.haiku-4.5-sub.rank\endcsname{disq.}
\expandafter\def\csname ccval@lineup.haiku-4.5-sub.seconds\endcsname{676.5}
\expandafter\def\csname ccval@lineup.haiku-4.5-sub.silent\endcsname{3}
\expandafter\def\csname ccval@lineup.haiku-4.5-sub.silent_per_pair\endcsname{0.33}
\expandafter\def\csname ccval@lineup.haiku-4.5-sub.silent_runs\endcsname{4}
\expandafter\def\csname ccval@lineup.haiku-4.5-sub.spread.badge_access.deviations.max\endcsname{0}
\expandafter\def\csname ccval@lineup.haiku-4.5-sub.spread.badge_access.deviations.median\endcsname{0}
\expandafter\def\csname ccval@lineup.haiku-4.5-sub.spread.badge_access.deviations.min\endcsname{0}
\expandafter\def\csname ccval@lineup.haiku-4.5-sub.spread.badge_access.draws\endcsname{3}
\expandafter\def\csname ccval@lineup.haiku-4.5-sub.spread.badge_access.silent.max\endcsname{0}
\expandafter\def\csname ccval@lineup.haiku-4.5-sub.spread.badge_access.silent.median\endcsname{0}
\expandafter\def\csname ccval@lineup.haiku-4.5-sub.spread.badge_access.silent.min\endcsname{0}
\expandafter\def\csname ccval@lineup.haiku-4.5-sub.spread.rate_limits.deviations.max\endcsname{14}
\expandafter\def\csname ccval@lineup.haiku-4.5-sub.spread.rate_limits.deviations.median\endcsname{9}
\expandafter\def\csname ccval@lineup.haiku-4.5-sub.spread.rate_limits.deviations.min\endcsname{8}
\expandafter\def\csname ccval@lineup.haiku-4.5-sub.spread.rate_limits.draws\endcsname{3}
\expandafter\def\csname ccval@lineup.haiku-4.5-sub.spread.rate_limits.silent.max\endcsname{25}
\expandafter\def\csname ccval@lineup.haiku-4.5-sub.spread.rate_limits.silent.median\endcsname{9}
\expandafter\def\csname ccval@lineup.haiku-4.5-sub.spread.rate_limits.silent.min\endcsname{2}
\expandafter\def\csname ccval@lineup.haiku-4.5-sub.spread.trace_names.deviations.max\endcsname{10}
\expandafter\def\csname ccval@lineup.haiku-4.5-sub.spread.trace_names.deviations.median\endcsname{10}
\expandafter\def\csname ccval@lineup.haiku-4.5-sub.spread.trace_names.deviations.min\endcsname{9}
\expandafter\def\csname ccval@lineup.haiku-4.5-sub.spread.trace_names.draws\endcsname{3}
\expandafter\def\csname ccval@lineup.haiku-4.5-sub.spread.trace_names.silent.max\endcsname{0}
\expandafter\def\csname ccval@lineup.haiku-4.5-sub.spread.trace_names.silent.median\endcsname{0}
\expandafter\def\csname ccval@lineup.haiku-4.5-sub.spread.trace_names.silent.min\endcsname{0}
\expandafter\def\csname ccval@lineup.haiku-4.5-sub.spread_pairs\endcsname{3}
\expandafter\def\csname ccval@lineup.haiku-4.5-sub.story.exit\endcsname{1}
\expandafter\def\csname ccval@lineup.haiku-4.5-sub.story.flagged\endcsname{no}
\expandafter\def\csname ccval@lineup.haiku-4.5-sub.story.kept_a\endcsname{yes}
\expandafter\def\csname ccval@lineup.haiku-4.5-sub.story.kept_b\endcsname{no}
\expandafter\def\csname ccval@lineup.haiku-4.5-sub.tokens\endcsname{122847}
\expandafter\def\csname ccval@lineup.haiku-4.5-sub.union_dev\endcsname{70}
\expandafter\def\csname ccval@lineup.haiku-4.5-sub.unit\endcsname{API-equivalent}
\expandafter\def\csname ccval@lineup.haiku-4.5-sub.untaken\endcsname{0}
\expandafter\def\csname ccval@lineup.haiku-4.5-sub.untaken_all\endcsname{1}
\expandafter\def\csname ccval@lineup.haiku-4.5-sub.untaken_runs\endcsname{1}
\expandafter\def\csname ccval@lineup.haiku-4.5-sub.usd\endcsname{0.483}
\expandafter\def\csname ccval@lineup.licence.bip\endcsname{BIP-39}
\expandafter\def\csname ccval@lineup.licence.bsd\endcsname{12}
\expandafter\def\csname ccval@lineup.licence.bsd_citation\endcsname{12}
\expandafter\def\csname ccval@lineup.licence.bsd_declared\endcsname{12}
\expandafter\def\csname ccval@lineup.licence.bsd_exit\endcsname{1}
\expandafter\def\csname ccval@lineup.licence.bsd_exit_one\endcsname{12}
\expandafter\def\csname ccval@lineup.licence.bsd_name\endcsname{2-clause BSD}
\expandafter\def\csname ccval@lineup.licence.gpl\endcsname{8}
\expandafter\def\csname ccval@lineup.licence.mit\endcsname{2}
\expandafter\def\csname ccval@lineup.licence.none\endcsname{2}
\expandafter\def\csname ccval@lineup.licence.other\endcsname{0}
\expandafter\def\csname ccval@lineup.licence.path\endcsname{tests/handwritten/bip39/}
\expandafter\def\csname ccval@lineup.licence.runs\endcsname{24}
\expandafter\def\csname ccval@lineup.licence.sourced.blind\endcsname{6}
\expandafter\def\csname ccval@lineup.licence.sourced.blind_mit\endcsname{0}
\expandafter\def\csname ccval@lineup.licence.sourced.date\endcsname{2026-09-25}
\expandafter\def\csname ccval@lineup.licence.sourced.looked\endcsname{21}
\expandafter\def\csname ccval@lineup.licence.sourced.looked_mit\endcsname{20}
\expandafter\def\csname ccval@lineup.licence.sourced.runs\endcsname{27}
\expandafter\def\csname ccval@lineup.list_output.lowest\endcsname{gpt-6-luna-api}
\expandafter\def\csname ccval@lineup.opus-5.5-api.S2.disqualified\endcsname{2}
\expandafter\def\csname ccval@lineup.opus-5.5-api.S2.fixtures\endcsname{13}
\expandafter\def\csname ccval@lineup.opus-5.5-api.S2.guards_wrong\endcsname{1}
\expandafter\def\csname ccval@lineup.opus-5.5-api.S2.invented\endcsname{2}
\expandafter\def\csname ccval@lineup.opus-5.5-api.S2.matched\endcsname{12}
\expandafter\def\csname ccval@lineup.opus-5.5-api.S2.measured\endcsname{13}
\expandafter\def\csname ccval@lineup.opus-5.5-api.S2.order\endcsname{3=}
\expandafter\def\csname ccval@lineup.opus-5.5-api.S2.plants\endcsname{7}
\expandafter\def\csname ccval@lineup.opus-5.5-api.S2.plants_measured\endcsname{7}
\expandafter\def\csname ccval@lineup.opus-5.5-api.S2.unmeasured\endcsname{0}
\expandafter\def\csname ccval@lineup.opus-5.5-api.S3.bip39.draws\endcsname{3}
\expandafter\def\csname ccval@lineup.opus-5.5-api.S3.bip39.max\endcsname{13}
\expandafter\def\csname ccval@lineup.opus-5.5-api.S3.bip39.median\endcsname{13}
\expandafter\def\csname ccval@lineup.opus-5.5-api.S3.bip39.min\endcsname{13}
\expandafter\def\csname ccval@lineup.opus-5.5-api.S3.bip39.total\endcsname{14}
\expandafter\def\csname ccval@lineup.opus-5.5-api.S3.fixed_total\endcsname{42}
\expandafter\def\csname ccval@lineup.opus-5.5-api.S3.mahjongg.draws\endcsname{1}
\expandafter\def\csname ccval@lineup.opus-5.5-api.S3.mahjongg.false\endcsname{0}
\expandafter\def\csname ccval@lineup.opus-5.5-api.S3.order\endcsname{1=}
\expandafter\def\csname ccval@lineup.opus-5.5-api.S3.voyager.draws\endcsname{3}
\expandafter\def\csname ccval@lineup.opus-5.5-api.S3.voyager.max\endcsname{32}
\expandafter\def\csname ccval@lineup.opus-5.5-api.S3.voyager.median\endcsname{29}
\expandafter\def\csname ccval@lineup.opus-5.5-api.S3.voyager.min\endcsname{28}
\expandafter\def\csname ccval@lineup.opus-5.5-api.S3.voyager.total\endcsname{44}
\expandafter\def\csname ccval@lineup.opus-5.5-api.audit.all.cases\endcsname{2}
\expandafter\def\csname ccval@lineup.opus-5.5-api.audit.all.gone\endcsname{0}
\expandafter\def\csname ccval@lineup.opus-5.5-api.audit.all.lost\endcsname{0}
\expandafter\def\csname ccval@lineup.opus-5.5-api.audit.all.not_content\endcsname{2}
\expandafter\def\csname ccval@lineup.opus-5.5-api.audit.all.partly\endcsname{0}
\expandafter\def\csname ccval@lineup.opus-5.5-api.audit.all.present\endcsname{0}
\expandafter\def\csname ccval@lineup.opus-5.5-api.audit.all.runs\endcsname{1}
\expandafter\def\csname ccval@lineup.opus-5.5-api.audit.borderline_only\endcsname{no}
\expandafter\def\csname ccval@lineup.opus-5.5-api.audit.cases\endcsname{2}
\expandafter\def\csname ccval@lineup.opus-5.5-api.audit.counted\endcsname{2}
\expandafter\def\csname ccval@lineup.opus-5.5-api.audit.counted.lost\endcsname{0}
\expandafter\def\csname ccval@lineup.opus-5.5-api.audit.disqualified\endcsname{no}
\expandafter\def\csname ccval@lineup.opus-5.5-api.audit.found\endcsname{0 of 2}
\expandafter\def\csname ccval@lineup.opus-5.5-api.audit.gone\endcsname{0}
\expandafter\def\csname ccval@lineup.opus-5.5-api.audit.lost\endcsname{0}
\expandafter\def\csname ccval@lineup.opus-5.5-api.audit.lost_borderline\endcsname{0}
\expandafter\def\csname ccval@lineup.opus-5.5-api.audit.not_content\endcsname{2}
\expandafter\def\csname ccval@lineup.opus-5.5-api.audit.partly\endcsname{0}
\expandafter\def\csname ccval@lineup.opus-5.5-api.audit.present\endcsname{0}
\expandafter\def\csname ccval@lineup.opus-5.5-api.audit.reworded\endcsname{0}
\expandafter\def\csname ccval@lineup.opus-5.5-api.audit.reworded.lost\endcsname{0}
\expandafter\def\csname ccval@lineup.opus-5.5-api.audit.second.lost\endcsname{0}
\expandafter\def\csname ccval@lineup.opus-5.5-api.dev_per_pair\endcsname{5.78}
\expandafter\def\csname ccval@lineup.opus-5.5-api.deviations\endcsname{52}
\expandafter\def\csname ccval@lineup.opus-5.5-api.effort\endcsname{medium}
\expandafter\def\csname ccval@lineup.opus-5.5-api.list_output_multiple\endcsname{40}
\expandafter\def\csname ccval@lineup.opus-5.5-api.model\endcsname{claude-opus-5-5}
\expandafter\def\csname ccval@lineup.opus-5.5-api.model_failures\endcsname{0}
\expandafter\def\csname ccval@lineup.opus-5.5-api.name\endcsname{opus-5.5-api}
\expandafter\def\csname ccval@lineup.opus-5.5-api.no_better_than_union\endcsname{no}
\expandafter\def\csname ccval@lineup.opus-5.5-api.pairs\endcsname{9}
\expandafter\def\csname ccval@lineup.opus-5.5-api.pairs_completed\endcsname{9}
\expandafter\def\csname ccval@lineup.opus-5.5-api.place\endcsname{3}
\expandafter\def\csname ccval@lineup.opus-5.5-api.rank\endcsname{4}
\expandafter\def\csname ccval@lineup.opus-5.5-api.seconds\endcsname{142.7}
\expandafter\def\csname ccval@lineup.opus-5.5-api.silent\endcsname{0}
\expandafter\def\csname ccval@lineup.opus-5.5-api.silent_per_pair\endcsname{0.00}
\expandafter\def\csname ccval@lineup.opus-5.5-api.silent_runs\endcsname{0}
\expandafter\def\csname ccval@lineup.opus-5.5-api.spread.badge_access.deviations.max\endcsname{0}
\expandafter\def\csname ccval@lineup.opus-5.5-api.spread.badge_access.deviations.median\endcsname{0}
\expandafter\def\csname ccval@lineup.opus-5.5-api.spread.badge_access.deviations.min\endcsname{0}
\expandafter\def\csname ccval@lineup.opus-5.5-api.spread.badge_access.draws\endcsname{3}
\expandafter\def\csname ccval@lineup.opus-5.5-api.spread.badge_access.silent.max\endcsname{0}
\expandafter\def\csname ccval@lineup.opus-5.5-api.spread.badge_access.silent.median\endcsname{0}
\expandafter\def\csname ccval@lineup.opus-5.5-api.spread.badge_access.silent.min\endcsname{0}
\expandafter\def\csname ccval@lineup.opus-5.5-api.spread.rate_limits.deviations.max\endcsname{23}
\expandafter\def\csname ccval@lineup.opus-5.5-api.spread.rate_limits.deviations.median\endcsname{20}
\expandafter\def\csname ccval@lineup.opus-5.5-api.spread.rate_limits.deviations.min\endcsname{20}
\expandafter\def\csname ccval@lineup.opus-5.5-api.spread.rate_limits.draws\endcsname{3}
\expandafter\def\csname ccval@lineup.opus-5.5-api.spread.rate_limits.silent.max\endcsname{0}
\expandafter\def\csname ccval@lineup.opus-5.5-api.spread.rate_limits.silent.median\endcsname{0}
\expandafter\def\csname ccval@lineup.opus-5.5-api.spread.rate_limits.silent.min\endcsname{0}
\expandafter\def\csname ccval@lineup.opus-5.5-api.spread.trace_names.deviations.max\endcsname{9}
\expandafter\def\csname ccval@lineup.opus-5.5-api.spread.trace_names.deviations.median\endcsname{7}
\expandafter\def\csname ccval@lineup.opus-5.5-api.spread.trace_names.deviations.min\endcsname{7}
\expandafter\def\csname ccval@lineup.opus-5.5-api.spread.trace_names.draws\endcsname{3}
\expandafter\def\csname ccval@lineup.opus-5.5-api.spread.trace_names.silent.max\endcsname{0}
\expandafter\def\csname ccval@lineup.opus-5.5-api.spread.trace_names.silent.median\endcsname{0}
\expandafter\def\csname ccval@lineup.opus-5.5-api.spread.trace_names.silent.min\endcsname{0}
\expandafter\def\csname ccval@lineup.opus-5.5-api.spread_pairs\endcsname{3}
\expandafter\def\csname ccval@lineup.opus-5.5-api.story.exit\endcsname{1}
\expandafter\def\csname ccval@lineup.opus-5.5-api.story.flagged\endcsname{yes}
\expandafter\def\csname ccval@lineup.opus-5.5-api.story.kept_a\endcsname{yes}
\expandafter\def\csname ccval@lineup.opus-5.5-api.story.kept_b\endcsname{no}
\expandafter\def\csname ccval@lineup.opus-5.5-api.tokens\endcsname{47656}
\expandafter\def\csname ccval@lineup.opus-5.5-api.union_dev\endcsname{70}
\expandafter\def\csname ccval@lineup.opus-5.5-api.unit\endcsname{list}
\expandafter\def\csname ccval@lineup.opus-5.5-api.untaken\endcsname{2}
\expandafter\def\csname ccval@lineup.opus-5.5-api.untaken_all\endcsname{2}
\expandafter\def\csname ccval@lineup.opus-5.5-api.untaken_runs\endcsname{1}
\expandafter\def\csname ccval@lineup.opus-5.5-api.usd\endcsname{0.485}
\expandafter\def\csname ccval@lineup.opus-5.5-sub.S2.disqualified\endcsname{4}
\expandafter\def\csname ccval@lineup.opus-5.5-sub.S2.fixtures\endcsname{13}
\expandafter\def\csname ccval@lineup.opus-5.5-sub.S2.guards_wrong\endcsname{3}
\expandafter\def\csname ccval@lineup.opus-5.5-sub.S2.invented\endcsname{4}
\expandafter\def\csname ccval@lineup.opus-5.5-sub.S2.matched\endcsname{11}
\expandafter\def\csname ccval@lineup.opus-5.5-sub.S2.measured\endcsname{13}
\expandafter\def\csname ccval@lineup.opus-5.5-sub.S2.order\endcsname{5}
\expandafter\def\csname ccval@lineup.opus-5.5-sub.S2.plants\endcsname{7}
\expandafter\def\csname ccval@lineup.opus-5.5-sub.S2.plants_measured\endcsname{7}
\expandafter\def\csname ccval@lineup.opus-5.5-sub.S2.unmeasured\endcsname{0}
\expandafter\def\csname ccval@lineup.opus-5.5-sub.S3.bip39.draws\endcsname{3}
\expandafter\def\csname ccval@lineup.opus-5.5-sub.S3.bip39.max\endcsname{13}
\expandafter\def\csname ccval@lineup.opus-5.5-sub.S3.bip39.median\endcsname{13}
\expandafter\def\csname ccval@lineup.opus-5.5-sub.S3.bip39.min\endcsname{13}
\expandafter\def\csname ccval@lineup.opus-5.5-sub.S3.bip39.total\endcsname{14}
\expandafter\def\csname ccval@lineup.opus-5.5-sub.S3.fixed_total\endcsname{45}
\expandafter\def\csname ccval@lineup.opus-5.5-sub.S3.mahjongg.draws\endcsname{1}
\expandafter\def\csname ccval@lineup.opus-5.5-sub.S3.mahjongg.false\endcsname{1}
\expandafter\def\csname ccval@lineup.opus-5.5-sub.S3.order\endcsname{1=}
\expandafter\def\csname ccval@lineup.opus-5.5-sub.S3.voyager.draws\endcsname{3}
\expandafter\def\csname ccval@lineup.opus-5.5-sub.S3.voyager.max\endcsname{35}
\expandafter\def\csname ccval@lineup.opus-5.5-sub.S3.voyager.median\endcsname{32}
\expandafter\def\csname ccval@lineup.opus-5.5-sub.S3.voyager.min\endcsname{30}
\expandafter\def\csname ccval@lineup.opus-5.5-sub.S3.voyager.total\endcsname{44}
\expandafter\def\csname ccval@lineup.opus-5.5-sub.audit.borderline_only\endcsname{no}
\expandafter\def\csname ccval@lineup.opus-5.5-sub.audit.cases\endcsname{0}
\expandafter\def\csname ccval@lineup.opus-5.5-sub.audit.counted\endcsname{0}
\expandafter\def\csname ccval@lineup.opus-5.5-sub.audit.counted.lost\endcsname{0}
\expandafter\def\csname ccval@lineup.opus-5.5-sub.audit.disqualified\endcsname{no}
\expandafter\def\csname ccval@lineup.opus-5.5-sub.audit.found\endcsname{--}
\expandafter\def\csname ccval@lineup.opus-5.5-sub.audit.gone\endcsname{0}
\expandafter\def\csname ccval@lineup.opus-5.5-sub.audit.lost\endcsname{0}
\expandafter\def\csname ccval@lineup.opus-5.5-sub.audit.lost_borderline\endcsname{0}
\expandafter\def\csname ccval@lineup.opus-5.5-sub.audit.not_content\endcsname{0}
\expandafter\def\csname ccval@lineup.opus-5.5-sub.audit.partly\endcsname{0}
\expandafter\def\csname ccval@lineup.opus-5.5-sub.audit.present\endcsname{0}
\expandafter\def\csname ccval@lineup.opus-5.5-sub.audit.reworded\endcsname{0}
\expandafter\def\csname ccval@lineup.opus-5.5-sub.audit.reworded.lost\endcsname{0}
\expandafter\def\csname ccval@lineup.opus-5.5-sub.dev_per_pair\endcsname{5.67}
\expandafter\def\csname ccval@lineup.opus-5.5-sub.deviations\endcsname{51}
\expandafter\def\csname ccval@lineup.opus-5.5-sub.effort\endcsname{medium}
\expandafter\def\csname ccval@lineup.opus-5.5-sub.model\endcsname{claude-opus-5-5}
\expandafter\def\csname ccval@lineup.opus-5.5-sub.model_failures\endcsname{0}
\expandafter\def\csname ccval@lineup.opus-5.5-sub.name\endcsname{opus-5.5-sub}
\expandafter\def\csname ccval@lineup.opus-5.5-sub.no_better_than_union\endcsname{no}
\expandafter\def\csname ccval@lineup.opus-5.5-sub.pairs\endcsname{9}
\expandafter\def\csname ccval@lineup.opus-5.5-sub.pairs_completed\endcsname{9}
\expandafter\def\csname ccval@lineup.opus-5.5-sub.place\endcsname{2=}
\expandafter\def\csname ccval@lineup.opus-5.5-sub.rank\endcsname{2=}
\expandafter\def\csname ccval@lineup.opus-5.5-sub.seconds\endcsname{164.1}
\expandafter\def\csname ccval@lineup.opus-5.5-sub.silent\endcsname{0}
\expandafter\def\csname ccval@lineup.opus-5.5-sub.silent_per_pair\endcsname{0.00}
\expandafter\def\csname ccval@lineup.opus-5.5-sub.silent_runs\endcsname{0}
\expandafter\def\csname ccval@lineup.opus-5.5-sub.spread.badge_access.deviations.max\endcsname{0}
\expandafter\def\csname ccval@lineup.opus-5.5-sub.spread.badge_access.deviations.median\endcsname{0}
\expandafter\def\csname ccval@lineup.opus-5.5-sub.spread.badge_access.deviations.min\endcsname{0}
\expandafter\def\csname ccval@lineup.opus-5.5-sub.spread.badge_access.draws\endcsname{3}
\expandafter\def\csname ccval@lineup.opus-5.5-sub.spread.badge_access.silent.max\endcsname{0}
\expandafter\def\csname ccval@lineup.opus-5.5-sub.spread.badge_access.silent.median\endcsname{0}
\expandafter\def\csname ccval@lineup.opus-5.5-sub.spread.badge_access.silent.min\endcsname{0}
\expandafter\def\csname ccval@lineup.opus-5.5-sub.spread.rate_limits.deviations.max\endcsname{23}
\expandafter\def\csname ccval@lineup.opus-5.5-sub.spread.rate_limits.deviations.median\endcsname{20}
\expandafter\def\csname ccval@lineup.opus-5.5-sub.spread.rate_limits.deviations.min\endcsname{20}
\expandafter\def\csname ccval@lineup.opus-5.5-sub.spread.rate_limits.draws\endcsname{3}
\expandafter\def\csname ccval@lineup.opus-5.5-sub.spread.rate_limits.silent.max\endcsname{0}
\expandafter\def\csname ccval@lineup.opus-5.5-sub.spread.rate_limits.silent.median\endcsname{0}
\expandafter\def\csname ccval@lineup.opus-5.5-sub.spread.rate_limits.silent.min\endcsname{0}
\expandafter\def\csname ccval@lineup.opus-5.5-sub.spread.trace_names.deviations.max\endcsname{8}
\expandafter\def\csname ccval@lineup.opus-5.5-sub.spread.trace_names.deviations.median\endcsname{7}
\expandafter\def\csname ccval@lineup.opus-5.5-sub.spread.trace_names.deviations.min\endcsname{7}
\expandafter\def\csname ccval@lineup.opus-5.5-sub.spread.trace_names.draws\endcsname{3}
\expandafter\def\csname ccval@lineup.opus-5.5-sub.spread.trace_names.silent.max\endcsname{0}
\expandafter\def\csname ccval@lineup.opus-5.5-sub.spread.trace_names.silent.median\endcsname{0}
\expandafter\def\csname ccval@lineup.opus-5.5-sub.spread.trace_names.silent.min\endcsname{0}
\expandafter\def\csname ccval@lineup.opus-5.5-sub.spread_pairs\endcsname{3}
\expandafter\def\csname ccval@lineup.opus-5.5-sub.story.exit\endcsname{1}
\expandafter\def\csname ccval@lineup.opus-5.5-sub.story.flagged\endcsname{yes}
\expandafter\def\csname ccval@lineup.opus-5.5-sub.story.kept_a\endcsname{yes}
\expandafter\def\csname ccval@lineup.opus-5.5-sub.story.kept_b\endcsname{no}
\expandafter\def\csname ccval@lineup.opus-5.5-sub.tokens\endcsname{55453}
\expandafter\def\csname ccval@lineup.opus-5.5-sub.union_dev\endcsname{70}
\expandafter\def\csname ccval@lineup.opus-5.5-sub.unit\endcsname{API-equivalent}
\expandafter\def\csname ccval@lineup.opus-5.5-sub.untaken\endcsname{0}
\expandafter\def\csname ccval@lineup.opus-5.5-sub.untaken_all\endcsname{0}
\expandafter\def\csname ccval@lineup.opus-5.5-sub.untaken_runs\endcsname{0}
\expandafter\def\csname ccval@lineup.opus-5.5-sub.usd\endcsname{0.614}
\expandafter\def\csname ccval@lineup.order.lead\endcsname{16}
\expandafter\def\csname ccval@lineup.order.rest_rows\endcsname{4}
\expandafter\def\csname ccval@lineup.order.rest_span\endcsname{2}
\expandafter\def\csname ccval@lineup.pairs\endcsname{9}
\expandafter\def\csname ccval@lineup.pin\endcsname{73b61563c13c}
\expandafter\def\csname ccval@lineup.registered.disqualified\endcsname{5}
\expandafter\def\csname ccval@lineup.registered.floor.base_only.dev_per_pair\endcsname{7.56}
\expandafter\def\csname ccval@lineup.registered.floor.base_only.deviations\endcsname{68}
\expandafter\def\csname ccval@lineup.registered.floor.base_only.silent\endcsname{129}
\expandafter\def\csname ccval@lineup.registered.floor.concatenation.dev_per_pair\endcsname{8.78}
\expandafter\def\csname ccval@lineup.registered.floor.concatenation.deviations\endcsname{79}
\expandafter\def\csname ccval@lineup.registered.floor.concatenation.silent\endcsname{0}
\expandafter\def\csname ccval@lineup.registered.floor.union.dev_per_pair\endcsname{7.78}
\expandafter\def\csname ccval@lineup.registered.floor.union.deviations\endcsname{70}
\expandafter\def\csname ccval@lineup.registered.floor.union.silent\endcsname{0}
\expandafter\def\csname ccval@lineup.registered.gpt-6-luna-api.dev_per_pair\endcsname{5.67}
\expandafter\def\csname ccval@lineup.registered.gpt-6-luna-api.deviations\endcsname{51}
\expandafter\def\csname ccval@lineup.registered.gpt-6-luna-api.no_better_than_union\endcsname{no}
\expandafter\def\csname ccval@lineup.registered.gpt-6-luna-api.rank\endcsname{disq.}
\expandafter\def\csname ccval@lineup.registered.gpt-6-luna-api.silent\endcsname{2}
\expandafter\def\csname ccval@lineup.registered.gpt-6-luna-api.union_dev\endcsname{70}
\expandafter\def\csname ccval@lineup.registered.gpt-6-sol-api.dev_per_pair\endcsname{4.11}
\expandafter\def\csname ccval@lineup.registered.gpt-6-sol-api.deviations\endcsname{37}
\expandafter\def\csname ccval@lineup.registered.gpt-6-sol-api.no_better_than_union\endcsname{no}
\expandafter\def\csname ccval@lineup.registered.gpt-6-sol-api.rank\endcsname{1}
\expandafter\def\csname ccval@lineup.registered.gpt-6-sol-api.silent\endcsname{0}
\expandafter\def\csname ccval@lineup.registered.gpt-6-sol-api.union_dev\endcsname{70}
\expandafter\def\csname ccval@lineup.registered.haiku-4.5-api.dev_per_pair\endcsname{6.25}
\expandafter\def\csname ccval@lineup.registered.haiku-4.5-api.deviations\endcsname{50}
\expandafter\def\csname ccval@lineup.registered.haiku-4.5-api.no_better_than_union\endcsname{no}
\expandafter\def\csname ccval@lineup.registered.haiku-4.5-api.rank\endcsname{disq.}
\expandafter\def\csname ccval@lineup.registered.haiku-4.5-api.silent\endcsname{13}
\expandafter\def\csname ccval@lineup.registered.haiku-4.5-api.union_dev\endcsname{63}
\expandafter\def\csname ccval@lineup.registered.haiku-4.5-sub.dev_per_pair\endcsname{7.56}
\expandafter\def\csname ccval@lineup.registered.haiku-4.5-sub.deviations\endcsname{68}
\expandafter\def\csname ccval@lineup.registered.haiku-4.5-sub.no_better_than_union\endcsname{no}
\expandafter\def\csname ccval@lineup.registered.haiku-4.5-sub.rank\endcsname{disq.}
\expandafter\def\csname ccval@lineup.registered.haiku-4.5-sub.silent\endcsname{3}
\expandafter\def\csname ccval@lineup.registered.haiku-4.5-sub.union_dev\endcsname{70}
\expandafter\def\csname ccval@lineup.registered.opus-5.5-api.dev_per_pair\endcsname{6.78}
\expandafter\def\csname ccval@lineup.registered.opus-5.5-api.deviations\endcsname{61}
\expandafter\def\csname ccval@lineup.registered.opus-5.5-api.no_better_than_union\endcsname{no}
\expandafter\def\csname ccval@lineup.registered.opus-5.5-api.rank\endcsname{disq.}
\expandafter\def\csname ccval@lineup.registered.opus-5.5-api.silent\endcsname{2}
\expandafter\def\csname ccval@lineup.registered.opus-5.5-api.union_dev\endcsname{70}
\expandafter\def\csname ccval@lineup.registered.opus-5.5-sub.dev_per_pair\endcsname{7.00}
\expandafter\def\csname ccval@lineup.registered.opus-5.5-sub.deviations\endcsname{63}
\expandafter\def\csname ccval@lineup.registered.opus-5.5-sub.no_better_than_union\endcsname{no}
\expandafter\def\csname ccval@lineup.registered.opus-5.5-sub.rank\endcsname{2}
\expandafter\def\csname ccval@lineup.registered.opus-5.5-sub.silent\endcsname{0}
\expandafter\def\csname ccval@lineup.registered.opus-5.5-sub.union_dev\endcsname{70}
\expandafter\def\csname ccval@lineup.registered.silent_runs\endcsname{18}
\expandafter\def\csname ccval@lineup.registered.silent_runs_reported\endcsname{18}
\expandafter\def\csname ccval@lineup.registered.silent_segments\endcsname{65}
\expandafter\def\csname ccval@lineup.registered.silent_segments_reported\endcsname{65}
\expandafter\def\csname ccval@lineup.registered.sonnet-5-api.dev_per_pair\endcsname{8.11}
\expandafter\def\csname ccval@lineup.registered.sonnet-5-api.deviations\endcsname{73}
\expandafter\def\csname ccval@lineup.registered.sonnet-5-api.no_better_than_union\endcsname{yes}
\expandafter\def\csname ccval@lineup.registered.sonnet-5-api.rank\endcsname{3}
\expandafter\def\csname ccval@lineup.registered.sonnet-5-api.silent\endcsname{0}
\expandafter\def\csname ccval@lineup.registered.sonnet-5-api.union_dev\endcsname{70}
\expandafter\def\csname ccval@lineup.registered.sonnet-5-sub.dev_per_pair\endcsname{7.33}
\expandafter\def\csname ccval@lineup.registered.sonnet-5-sub.deviations\endcsname{66}
\expandafter\def\csname ccval@lineup.registered.sonnet-5-sub.no_better_than_union\endcsname{no}
\expandafter\def\csname ccval@lineup.registered.sonnet-5-sub.rank\endcsname{disq.}
\expandafter\def\csname ccval@lineup.registered.sonnet-5-sub.silent\endcsname{5}
\expandafter\def\csname ccval@lineup.registered.sonnet-5-sub.union_dev\endcsname{70}
\expandafter\def\csname ccval@lineup.registered.story.silent\endcsname{0}
\expandafter\def\csname ccval@lineup.registered.survivors\endcsname{3}
\expandafter\def\csname ccval@lineup.rows\endcsname{8}
\expandafter\def\csname ccval@lineup.scoring.shown_first\endcsname{amended}
\expandafter\def\csname ccval@lineup.second_defect.reworded\endcsname{173}
\expandafter\def\csname ccval@lineup.second_defect.sentences\endcsname{181}
\expandafter\def\csname ccval@lineup.side.compared_with\endcsname{opus-5.5-sub}
\expandafter\def\csname ccval@lineup.side.fable-sub.deviations\endcsname{18}
\expandafter\def\csname ccval@lineup.side.fable-sub.index_429.deviations\endcsname{1}
\expandafter\def\csname ccval@lineup.side.fable-sub.model\endcsname{claude-fable-5-1}
\expandafter\def\csname ccval@lineup.side.fable-sub.name\endcsname{fable-sub}
\expandafter\def\csname ccval@lineup.side.fable-sub.pairs_completed\endcsname{2}
\expandafter\def\csname ccval@lineup.side.fable-sub.rate_limits.deviations\endcsname{17}
\expandafter\def\csname ccval@lineup.side.fable-sub.seconds\endcsname{359.9}
\expandafter\def\csname ccval@lineup.side.fable-sub.silent\endcsname{0}
\expandafter\def\csname ccval@lineup.side.fable-sub.usd\endcsname{2.721}
\expandafter\def\csname ccval@lineup.side.opus-5.5-sub.deviations\endcsname{27}
\expandafter\def\csname ccval@lineup.side.opus-5.5-sub.index_429.deviations\endcsname{4}
\expandafter\def\csname ccval@lineup.side.opus-5.5-sub.name\endcsname{opus-5.5-sub}
\expandafter\def\csname ccval@lineup.side.opus-5.5-sub.pairs_completed\endcsname{2}
\expandafter\def\csname ccval@lineup.side.opus-5.5-sub.rate_limits.deviations\endcsname{23}
\expandafter\def\csname ccval@lineup.side.opus-5.5-sub.seconds\endcsname{329.6}
\expandafter\def\csname ccval@lineup.side.opus-5.5-sub.silent\endcsname{0}
\expandafter\def\csname ccval@lineup.side.opus-5.5-sub.usd\endcsname{1.177}
\expandafter\def\csname ccval@lineup.side.pair.index_429\endcsname{index\_429}
\expandafter\def\csname ccval@lineup.side.pair.rate_limits\endcsname{rate\_limits}
\expandafter\def\csname ccval@lineup.side.pairs\endcsname{2}
\expandafter\def\csname ccval@lineup.silent_runs\endcsname{10}
\expandafter\def\csname ccval@lineup.silent_segments\endcsname{45}
\expandafter\def\csname ccval@lineup.sonnet-5-api.S2.disqualified\endcsname{0}
\expandafter\def\csname ccval@lineup.sonnet-5-api.S2.fixtures\endcsname{13}
\expandafter\def\csname ccval@lineup.sonnet-5-api.S2.guards_wrong\endcsname{0}
\expandafter\def\csname ccval@lineup.sonnet-5-api.S2.invented\endcsname{0}
\expandafter\def\csname ccval@lineup.sonnet-5-api.S2.matched\endcsname{13}
\expandafter\def\csname ccval@lineup.sonnet-5-api.S2.measured\endcsname{13}
\expandafter\def\csname ccval@lineup.sonnet-5-api.S2.order\endcsname{1=}
\expandafter\def\csname ccval@lineup.sonnet-5-api.S2.plants\endcsname{7}
\expandafter\def\csname ccval@lineup.sonnet-5-api.S2.plants_measured\endcsname{7}
\expandafter\def\csname ccval@lineup.sonnet-5-api.S2.unmeasured\endcsname{0}
\expandafter\def\csname ccval@lineup.sonnet-5-api.S3.bip39.draws\endcsname{3}
\expandafter\def\csname ccval@lineup.sonnet-5-api.S3.bip39.max\endcsname{12}
\expandafter\def\csname ccval@lineup.sonnet-5-api.S3.bip39.median\endcsname{12}
\expandafter\def\csname ccval@lineup.sonnet-5-api.S3.bip39.min\endcsname{12}
\expandafter\def\csname ccval@lineup.sonnet-5-api.S3.bip39.total\endcsname{14}
\expandafter\def\csname ccval@lineup.sonnet-5-api.S3.fixed_total\endcsname{33}
\expandafter\def\csname ccval@lineup.sonnet-5-api.S3.mahjongg.draws\endcsname{1}
\expandafter\def\csname ccval@lineup.sonnet-5-api.S3.mahjongg.false\endcsname{0}
\expandafter\def\csname ccval@lineup.sonnet-5-api.S3.order\endcsname{2=}
\expandafter\def\csname ccval@lineup.sonnet-5-api.S3.voyager.draws\endcsname{3}
\expandafter\def\csname ccval@lineup.sonnet-5-api.S3.voyager.max\endcsname{23}
\expandafter\def\csname ccval@lineup.sonnet-5-api.S3.voyager.median\endcsname{21}
\expandafter\def\csname ccval@lineup.sonnet-5-api.S3.voyager.min\endcsname{20}
\expandafter\def\csname ccval@lineup.sonnet-5-api.S3.voyager.total\endcsname{44}
\expandafter\def\csname ccval@lineup.sonnet-5-api.audit.borderline_only\endcsname{no}
\expandafter\def\csname ccval@lineup.sonnet-5-api.audit.cases\endcsname{0}
\expandafter\def\csname ccval@lineup.sonnet-5-api.audit.counted\endcsname{0}
\expandafter\def\csname ccval@lineup.sonnet-5-api.audit.counted.lost\endcsname{0}
\expandafter\def\csname ccval@lineup.sonnet-5-api.audit.disqualified\endcsname{no}
\expandafter\def\csname ccval@lineup.sonnet-5-api.audit.found\endcsname{--}
\expandafter\def\csname ccval@lineup.sonnet-5-api.audit.gone\endcsname{0}
\expandafter\def\csname ccval@lineup.sonnet-5-api.audit.lost\endcsname{0}
\expandafter\def\csname ccval@lineup.sonnet-5-api.audit.lost_borderline\endcsname{0}
\expandafter\def\csname ccval@lineup.sonnet-5-api.audit.not_content\endcsname{0}
\expandafter\def\csname ccval@lineup.sonnet-5-api.audit.partly\endcsname{0}
\expandafter\def\csname ccval@lineup.sonnet-5-api.audit.present\endcsname{0}
\expandafter\def\csname ccval@lineup.sonnet-5-api.audit.reworded\endcsname{0}
\expandafter\def\csname ccval@lineup.sonnet-5-api.audit.reworded.lost\endcsname{0}
\expandafter\def\csname ccval@lineup.sonnet-5-api.dev_per_pair\endcsname{5.89}
\expandafter\def\csname ccval@lineup.sonnet-5-api.deviations\endcsname{53}
\expandafter\def\csname ccval@lineup.sonnet-5-api.effort\endcsname{high}
\expandafter\def\csname ccval@lineup.sonnet-5-api.list_output_multiple\endcsname{20}
\expandafter\def\csname ccval@lineup.sonnet-5-api.model\endcsname{claude-sonnet-5}
\expandafter\def\csname ccval@lineup.sonnet-5-api.model_failures\endcsname{0}
\expandafter\def\csname ccval@lineup.sonnet-5-api.name\endcsname{sonnet-5-api}
\expandafter\def\csname ccval@lineup.sonnet-5-api.no_better_than_union\endcsname{no}
\expandafter\def\csname ccval@lineup.sonnet-5-api.pairs\endcsname{9}
\expandafter\def\csname ccval@lineup.sonnet-5-api.pairs_completed\endcsname{9}
\expandafter\def\csname ccval@lineup.sonnet-5-api.place\endcsname{4}
\expandafter\def\csname ccval@lineup.sonnet-5-api.rank\endcsname{5}
\expandafter\def\csname ccval@lineup.sonnet-5-api.seconds\endcsname{332.4}
\expandafter\def\csname ccval@lineup.sonnet-5-api.silent\endcsname{0}
\expandafter\def\csname ccval@lineup.sonnet-5-api.silent_per_pair\endcsname{0.00}
\expandafter\def\csname ccval@lineup.sonnet-5-api.silent_runs\endcsname{0}
\expandafter\def\csname ccval@lineup.sonnet-5-api.spread_pairs\endcsname{0}
\expandafter\def\csname ccval@lineup.sonnet-5-api.story.exit\endcsname{1}
\expandafter\def\csname ccval@lineup.sonnet-5-api.story.flagged\endcsname{yes}
\expandafter\def\csname ccval@lineup.sonnet-5-api.story.kept_a\endcsname{yes}
\expandafter\def\csname ccval@lineup.sonnet-5-api.story.kept_b\endcsname{no}
\expandafter\def\csname ccval@lineup.sonnet-5-api.tokens\endcsname{67974}
\expandafter\def\csname ccval@lineup.sonnet-5-api.union_dev\endcsname{70}
\expandafter\def\csname ccval@lineup.sonnet-5-api.unit\endcsname{list}
\expandafter\def\csname ccval@lineup.sonnet-5-api.untaken\endcsname{0}
\expandafter\def\csname ccval@lineup.sonnet-5-api.untaken_all\endcsname{0}
\expandafter\def\csname ccval@lineup.sonnet-5-api.untaken_runs\endcsname{0}
\expandafter\def\csname ccval@lineup.sonnet-5-api.usd\endcsname{0.462}
\expandafter\def\csname ccval@lineup.sonnet-5-sub.S2.disqualified\endcsname{0}
\expandafter\def\csname ccval@lineup.sonnet-5-sub.S2.fixtures\endcsname{13}
\expandafter\def\csname ccval@lineup.sonnet-5-sub.S2.guards_wrong\endcsname{0}
\expandafter\def\csname ccval@lineup.sonnet-5-sub.S2.invented\endcsname{0}
\expandafter\def\csname ccval@lineup.sonnet-5-sub.S2.matched\endcsname{13}
\expandafter\def\csname ccval@lineup.sonnet-5-sub.S2.measured\endcsname{13}
\expandafter\def\csname ccval@lineup.sonnet-5-sub.S2.order\endcsname{1=}
\expandafter\def\csname ccval@lineup.sonnet-5-sub.S2.plants\endcsname{7}
\expandafter\def\csname ccval@lineup.sonnet-5-sub.S2.plants_measured\endcsname{7}
\expandafter\def\csname ccval@lineup.sonnet-5-sub.S2.unmeasured\endcsname{0}
\expandafter\def\csname ccval@lineup.sonnet-5-sub.S3.bip39.draws\endcsname{3}
\expandafter\def\csname ccval@lineup.sonnet-5-sub.S3.bip39.max\endcsname{12}
\expandafter\def\csname ccval@lineup.sonnet-5-sub.S3.bip39.median\endcsname{12}
\expandafter\def\csname ccval@lineup.sonnet-5-sub.S3.bip39.min\endcsname{12}
\expandafter\def\csname ccval@lineup.sonnet-5-sub.S3.bip39.total\endcsname{14}
\expandafter\def\csname ccval@lineup.sonnet-5-sub.S3.fixed_total\endcsname{32}
\expandafter\def\csname ccval@lineup.sonnet-5-sub.S3.mahjongg.draws\endcsname{1}
\expandafter\def\csname ccval@lineup.sonnet-5-sub.S3.mahjongg.false\endcsname{0}
\expandafter\def\csname ccval@lineup.sonnet-5-sub.S3.order\endcsname{2=}
\expandafter\def\csname ccval@lineup.sonnet-5-sub.S3.voyager.draws\endcsname{3}
\expandafter\def\csname ccval@lineup.sonnet-5-sub.S3.voyager.max\endcsname{23}
\expandafter\def\csname ccval@lineup.sonnet-5-sub.S3.voyager.median\endcsname{20}
\expandafter\def\csname ccval@lineup.sonnet-5-sub.S3.voyager.min\endcsname{18}
\expandafter\def\csname ccval@lineup.sonnet-5-sub.S3.voyager.total\endcsname{44}
\expandafter\def\csname ccval@lineup.sonnet-5-sub.audit.all.cases\endcsname{5}
\expandafter\def\csname ccval@lineup.sonnet-5-sub.audit.all.gone\endcsname{0}
\expandafter\def\csname ccval@lineup.sonnet-5-sub.audit.all.lost\endcsname{0}
\expandafter\def\csname ccval@lineup.sonnet-5-sub.audit.all.not_content\endcsname{0}
\expandafter\def\csname ccval@lineup.sonnet-5-sub.audit.all.partly\endcsname{0}
\expandafter\def\csname ccval@lineup.sonnet-5-sub.audit.all.present\endcsname{5}
\expandafter\def\csname ccval@lineup.sonnet-5-sub.audit.all.runs\endcsname{2}
\expandafter\def\csname ccval@lineup.sonnet-5-sub.audit.borderline_only\endcsname{no}
\expandafter\def\csname ccval@lineup.sonnet-5-sub.audit.cases\endcsname{9}
\expandafter\def\csname ccval@lineup.sonnet-5-sub.audit.counted\endcsname{5}
\expandafter\def\csname ccval@lineup.sonnet-5-sub.audit.counted.lost\endcsname{0}
\expandafter\def\csname ccval@lineup.sonnet-5-sub.audit.disqualified\endcsname{no}
\expandafter\def\csname ccval@lineup.sonnet-5-sub.audit.found\endcsname{0 of 9}
\expandafter\def\csname ccval@lineup.sonnet-5-sub.audit.gone\endcsname{0}
\expandafter\def\csname ccval@lineup.sonnet-5-sub.audit.lost\endcsname{0}
\expandafter\def\csname ccval@lineup.sonnet-5-sub.audit.lost_borderline\endcsname{0}
\expandafter\def\csname ccval@lineup.sonnet-5-sub.audit.not_content\endcsname{0}
\expandafter\def\csname ccval@lineup.sonnet-5-sub.audit.partly\endcsname{0}
\expandafter\def\csname ccval@lineup.sonnet-5-sub.audit.present\endcsname{9}
\expandafter\def\csname ccval@lineup.sonnet-5-sub.audit.reworded\endcsname{4}
\expandafter\def\csname ccval@lineup.sonnet-5-sub.audit.reworded.lost\endcsname{0}
\expandafter\def\csname ccval@lineup.sonnet-5-sub.audit.second.lost\endcsname{0}
\expandafter\def\csname ccval@lineup.sonnet-5-sub.dev_per_pair\endcsname{5.67}
\expandafter\def\csname ccval@lineup.sonnet-5-sub.deviations\endcsname{51}
\expandafter\def\csname ccval@lineup.sonnet-5-sub.effort\endcsname{medium}
\expandafter\def\csname ccval@lineup.sonnet-5-sub.model\endcsname{claude-sonnet-5}
\expandafter\def\csname ccval@lineup.sonnet-5-sub.model_failures\endcsname{0}
\expandafter\def\csname ccval@lineup.sonnet-5-sub.name\endcsname{sonnet-5-sub}
\expandafter\def\csname ccval@lineup.sonnet-5-sub.no_better_than_union\endcsname{no}
\expandafter\def\csname ccval@lineup.sonnet-5-sub.pairs\endcsname{9}
\expandafter\def\csname ccval@lineup.sonnet-5-sub.pairs_completed\endcsname{9}
\expandafter\def\csname ccval@lineup.sonnet-5-sub.place\endcsname{2=}
\expandafter\def\csname ccval@lineup.sonnet-5-sub.rank\endcsname{2=}
\expandafter\def\csname ccval@lineup.sonnet-5-sub.seconds\endcsname{186.5}
\expandafter\def\csname ccval@lineup.sonnet-5-sub.silent\endcsname{2}
\expandafter\def\csname ccval@lineup.sonnet-5-sub.silent_per_pair\endcsname{0.22}
\expandafter\def\csname ccval@lineup.sonnet-5-sub.silent_runs\endcsname{1}
\expandafter\def\csname ccval@lineup.sonnet-5-sub.spread.badge_access.deviations.max\endcsname{0}
\expandafter\def\csname ccval@lineup.sonnet-5-sub.spread.badge_access.deviations.median\endcsname{0}
\expandafter\def\csname ccval@lineup.sonnet-5-sub.spread.badge_access.deviations.min\endcsname{0}
\expandafter\def\csname ccval@lineup.sonnet-5-sub.spread.badge_access.draws\endcsname{3}
\expandafter\def\csname ccval@lineup.sonnet-5-sub.spread.badge_access.silent.max\endcsname{0}
\expandafter\def\csname ccval@lineup.sonnet-5-sub.spread.badge_access.silent.median\endcsname{0}
\expandafter\def\csname ccval@lineup.sonnet-5-sub.spread.badge_access.silent.min\endcsname{0}
\expandafter\def\csname ccval@lineup.sonnet-5-sub.spread.rate_limits.deviations.max\endcsname{19}
\expandafter\def\csname ccval@lineup.sonnet-5-sub.spread.rate_limits.deviations.median\endcsname{17}
\expandafter\def\csname ccval@lineup.sonnet-5-sub.spread.rate_limits.deviations.min\endcsname{16}
\expandafter\def\csname ccval@lineup.sonnet-5-sub.spread.rate_limits.draws\endcsname{3}
\expandafter\def\csname ccval@lineup.sonnet-5-sub.spread.rate_limits.silent.max\endcsname{2}
\expandafter\def\csname ccval@lineup.sonnet-5-sub.spread.rate_limits.silent.median\endcsname{0}
\expandafter\def\csname ccval@lineup.sonnet-5-sub.spread.rate_limits.silent.min\endcsname{0}
\expandafter\def\csname ccval@lineup.sonnet-5-sub.spread.trace_names.deviations.max\endcsname{7}
\expandafter\def\csname ccval@lineup.sonnet-5-sub.spread.trace_names.deviations.median\endcsname{5}
\expandafter\def\csname ccval@lineup.sonnet-5-sub.spread.trace_names.deviations.min\endcsname{4}
\expandafter\def\csname ccval@lineup.sonnet-5-sub.spread.trace_names.draws\endcsname{3}
\expandafter\def\csname ccval@lineup.sonnet-5-sub.spread.trace_names.silent.max\endcsname{0}
\expandafter\def\csname ccval@lineup.sonnet-5-sub.spread.trace_names.silent.median\endcsname{0}
\expandafter\def\csname ccval@lineup.sonnet-5-sub.spread.trace_names.silent.min\endcsname{0}
\expandafter\def\csname ccval@lineup.sonnet-5-sub.spread_pairs\endcsname{3}
\expandafter\def\csname ccval@lineup.sonnet-5-sub.story.exit\endcsname{1}
\expandafter\def\csname ccval@lineup.sonnet-5-sub.story.flagged\endcsname{yes}
\expandafter\def\csname ccval@lineup.sonnet-5-sub.story.kept_a\endcsname{yes}
\expandafter\def\csname ccval@lineup.sonnet-5-sub.story.kept_b\endcsname{no}
\expandafter\def\csname ccval@lineup.sonnet-5-sub.tokens\endcsname{73188}
\expandafter\def\csname ccval@lineup.sonnet-5-sub.union_dev\endcsname{70}
\expandafter\def\csname ccval@lineup.sonnet-5-sub.unit\endcsname{API-equivalent}
\expandafter\def\csname ccval@lineup.sonnet-5-sub.untaken\endcsname{0}
\expandafter\def\csname ccval@lineup.sonnet-5-sub.untaken_all\endcsname{0}
\expandafter\def\csname ccval@lineup.sonnet-5-sub.untaken_runs\endcsname{0}
\expandafter\def\csname ccval@lineup.sonnet-5-sub.usd\endcsname{0.309}
\expandafter\def\csname ccval@lineup.spend.anthropic\endcsname{34.29}
\expandafter\def\csname ccval@lineup.spend.fable\endcsname{7.33}
\expandafter\def\csname ccval@lineup.spend.openai\endcsname{7.26}
\expandafter\def\csname ccval@lineup.spend.subscription\endcsname{49.47}
\expandafter\def\csname ccval@lineup.spread.max_width\endcsname{3}
\expandafter\def\csname ccval@lineup.spread.max_width.ranges\endcsname{4}
\expandafter\def\csname ccval@lineup.spread.max_width.rows\endcsname{3}
\expandafter\def\csname ccval@lineup.spread.repeating_rows\endcsname{4}
\expandafter\def\csname ccval@lineup.story.a_count\endcsname{2}
\expandafter\def\csname ccval@lineup.story.b_count\endcsname{3}
\expandafter\def\csname ccval@lineup.story.b_line\endcsname{15}
\expandafter\def\csname ccval@lineup.story.code\endcsname{E2}
\expandafter\def\csname ccval@lineup.story.disposition\endcsname{superseded}
\expandafter\def\csname ccval@lineup.story.dropped\endcsname{7}
\expandafter\def\csname ccval@lineup.story.exit\endcsname{1}
\expandafter\def\csname ccval@lineup.story.flagged\endcsname{5}
\expandafter\def\csname ccval@lineup.story.merged_line\endcsname{17}
\expandafter\def\csname ccval@lineup.story.reason\endcsname{Threshold conflict resolved in favour of the base value.}
\expandafter\def\csname ccval@lineup.story.reference_keeps\endcsname{a}
\expandafter\def\csname ccval@lineup.story.row_label\endcsname{Opus 5.5}
\expandafter\def\csname ccval@lineup.story.rows\endcsname{8}
\expandafter\def\csname ccval@lineup.story.silent\endcsname{0}
\expandafter\def\csname ccval@lineup.story.untaken\endcsname{0}
\expandafter\def\csname ccval@lineup.survivors\endcsname{5}
\expandafter\def\csname ccval@lineup.to\endcsname{2026-09-28}
\expandafter\def\csname ccval@lineup.unmeasured.name\endcsname{Gemini 3.8 Flash}
\expandafter\def\csname ccval@lineup.untaken_segments\endcsname{12}
\expandafter\def\csname ccval@matrix.block1.arm1\endcsname{qwen3:14b}
\expandafter\def\csname ccval@matrix.block1.arm2\endcsname{qwen3:4b}
\expandafter\def\csname ccval@matrix.block1.arm3\endcsname{qwen3:8b}
\expandafter\def\csname ccval@matrix.block1.arm4\endcsname{qwen3.8:27b}
\expandafter\def\csname ccval@matrix.block1.arms\endcsname{qwen3:14b, qwen3:4b, qwen3:8b, qwen3.8:27b}
\expandafter\def\csname ccval@matrix.block1.registered\endcsname{2026-09-06}
\expandafter\def\csname ccval@matrix.block2.arm1\endcsname{llama3.3:70b}
\expandafter\def\csname ccval@matrix.block2.arm2\endcsname{deepseek-r1:70b}
\expandafter\def\csname ccval@matrix.block2.arms\endcsname{llama3.3:70b, deepseek-r1:70b}
\expandafter\def\csname ccval@matrix.block2.registered\endcsname{2026-09-06}
\expandafter\def\csname ccval@matrix.block3.arm1\endcsname{gemma3:27b}
\expandafter\def\csname ccval@matrix.block3.arm2\endcsname{mistral-small}
\expandafter\def\csname ccval@matrix.block3.arm3\endcsname{qwen3.8:27b}
\expandafter\def\csname ccval@matrix.block3.arms\endcsname{gemma3:27b, mistral-small, qwen3.8:27b}
\expandafter\def\csname ccval@matrix.block3.registered\endcsname{2026-09-06}
\expandafter\def\csname ccval@matrix.block4.arm1\endcsname{gpt-oss:20b}
\expandafter\def\csname ccval@matrix.block4.arm2\endcsname{gpt-oss:120b}
\expandafter\def\csname ccval@matrix.block4.arms\endcsname{gpt-oss:20b, gpt-oss:120b}
\expandafter\def\csname ccval@matrix.block4.registered\endcsname{2026-09-06}
\expandafter\def\csname ccval@matrix.block5.registered\endcsname{2026-08-31}
\expandafter\def\csname ccval@merge.A-27b.forward_declared\endcsname{93}
\expandafter\def\csname ccval@merge.A-27b.forward_measured\endcsname{93}
\expandafter\def\csname ccval@merge.A-27b.forward_supported\endcsname{93}
\expandafter\def\csname ccval@merge.A-27b.model\endcsname{qwen3.8:27b}
\expandafter\def\csname ccval@merge.A-27b.samples\endcsname{1}
\expandafter\def\csname ccval@merge.B-70b.forward_declared\endcsname{93}
\expandafter\def\csname ccval@merge.B-70b.forward_measured\endcsname{93}
\expandafter\def\csname ccval@merge.B-70b.forward_supported\endcsname{91}
\expandafter\def\csname ccval@merge.B-70b.model\endcsname{deepseek-r1:70b}
\expandafter\def\csname ccval@merge.B-70b.samples\endcsname{1}
\expandafter\def\csname ccval@merge.C-120b.forward_declared\endcsname{93}
\expandafter\def\csname ccval@merge.C-120b.forward_measured\endcsname{93}
\expandafter\def\csname ccval@merge.C-120b.forward_supported\endcsname{93}
\expandafter\def\csname ccval@merge.C-120b.model\endcsname{gpt-oss:120b}
\expandafter\def\csname ccval@merge.C-120b.samples\endcsname{1}
\expandafter\def\csname ccval@opus5.id\endcsname{claude-opus-5}
\expandafter\def\csname ccval@opus5.replaced_by\endcsname{claude-opus-5-5}
\expandafter\def\csname ccval@opus5.retired_on\endcsname{2026-09-25}
\expandafter\def\csname ccval@phase4.27b-fp8.claims.max\endcsname{75}
\expandafter\def\csname ccval@phase4.27b-fp8.claims.min\endcsname{59}
\expandafter\def\csname ccval@phase4.27b-fp8.model\endcsname{Qwen/Qwen3.8-27B-FP8}
\expandafter\def\csname ccval@phase4.8b.claims.max\endcsname{20}
\expandafter\def\csname ccval@phase4.8b.claims.min\endcsname{15}
\expandafter\def\csname ccval@phase4.8b.model\endcsname{qwen/qwen3-8b}
\expandafter\def\csname ccval@phase4.date\endcsname{2026-09-16}
\expandafter\def\csname ccval@planted.bip39.errors\endcsname{14}
\expandafter\def\csname ccval@planted.voyager.errors\endcsname{44}
\expandafter\def\csname ccval@reg.2026-09-24-depth.amendments\endcsname{0}
\expandafter\def\csname ccval@reg.2026-09-24-depth.date\endcsname{2026-09-24}
\expandafter\def\csname ccval@reg.2026-09-24-depth.now\endcsname{4f68c3271576}
\expandafter\def\csname ccval@reg.2026-09-24-depth.path\endcsname{arms/2026-09-24/depth}
\expandafter\def\csname ccval@reg.2026-09-24-depth.recorded\endcsname{none kept}
\expandafter\def\csname ccval@reg.2026-09-24-subscription.amendments\endcsname{0}
\expandafter\def\csname ccval@reg.2026-09-24-subscription.date\endcsname{2026-09-24}
\expandafter\def\csname ccval@reg.2026-09-24-subscription.now\endcsname{3a627876c69c}
\expandafter\def\csname ccval@reg.2026-09-24-subscription.path\endcsname{arms/2026-09-24/subscription}
\expandafter\def\csname ccval@reg.2026-09-24-subscription.recorded\endcsname{none kept}
\expandafter\def\csname ccval@reg.2026-09-25-google.amendments\endcsname{4}
\expandafter\def\csname ccval@reg.2026-09-25-google.date\endcsname{2026-09-25}
\expandafter\def\csname ccval@reg.2026-09-25-google.now\endcsname{7abb032866f6}
\expandafter\def\csname ccval@reg.2026-09-25-google.path\endcsname{arms/2026-09-25/google}
\expandafter\def\csname ccval@reg.2026-09-25-google.recorded\endcsname{22a7b91b4484}
\expandafter\def\csname ccval@reg.2026-09-25-subscription-comparison.amendments\endcsname{1}
\expandafter\def\csname ccval@reg.2026-09-25-subscription-comparison.date\endcsname{2026-09-25}
\expandafter\def\csname ccval@reg.2026-09-25-subscription-comparison.now\endcsname{07af34914ab9}
\expandafter\def\csname ccval@reg.2026-09-25-subscription-comparison.path\endcsname{arms/2026-09-25/subscription-comparison}
\expandafter\def\csname ccval@reg.2026-09-25-subscription-comparison.recorded\endcsname{08a82717e20b}
\expandafter\def\csname ccval@reg.2026-09-25-vendor.amendments\endcsname{3}
\expandafter\def\csname ccval@reg.2026-09-25-vendor.date\endcsname{2026-09-25}
\expandafter\def\csname ccval@reg.2026-09-25-vendor.now\endcsname{6794b479e5fc}
\expandafter\def\csname ccval@reg.2026-09-25-vendor.path\endcsname{arms/2026-09-25/vendor}
\expandafter\def\csname ccval@reg.2026-09-25-vendor.recorded\endcsname{21d4125562f4}
\expandafter\def\csname ccval@reg.2026-09-26-opus-max.amendments\endcsname{1}
\expandafter\def\csname ccval@reg.2026-09-26-opus-max.date\endcsname{2026-09-26}
\expandafter\def\csname ccval@reg.2026-09-26-opus-max.now\endcsname{da0e5f664b66}
\expandafter\def\csname ccval@reg.2026-09-26-opus-max.path\endcsname{arms/2026-09-26/opus-max}
\expandafter\def\csname ccval@reg.2026-09-26-opus-max.recorded\endcsname{da0e5f664b66}
\expandafter\def\csname ccval@reg.2026-09-26-safemode-spot.amendments\endcsname{2}
\expandafter\def\csname ccval@reg.2026-09-26-safemode-spot.date\endcsname{2026-09-26}
\expandafter\def\csname ccval@reg.2026-09-26-safemode-spot.now\endcsname{3d24e2d04b57}
\expandafter\def\csname ccval@reg.2026-09-26-safemode-spot.path\endcsname{arms/2026-09-26/safemode-spot}
\expandafter\def\csname ccval@reg.2026-09-26-safemode-spot.recorded\endcsname{3d24e2d04b57}
\expandafter\def\csname ccval@reg.2026-09-27-lineup-pilot.amendments\endcsname{0}
\expandafter\def\csname ccval@reg.2026-09-27-lineup-pilot.date\endcsname{2026-09-27}
\expandafter\def\csname ccval@reg.2026-09-27-lineup-pilot.now\endcsname{047562a2afc7}
\expandafter\def\csname ccval@reg.2026-09-27-lineup-pilot.path\endcsname{arms/2026-09-27/lineup-pilot}
\expandafter\def\csname ccval@reg.2026-09-27-lineup-pilot.recorded\endcsname{none kept}
\expandafter\def\csname ccval@reg.2026-09-27-lineup.amendments\endcsname{6}
\expandafter\def\csname ccval@reg.2026-09-27-lineup.date\endcsname{2026-09-27}
\expandafter\def\csname ccval@reg.2026-09-27-lineup.now\endcsname{1d0d5eea65ee}
\expandafter\def\csname ccval@reg.2026-09-27-lineup.path\endcsname{arms/2026-09-27/lineup}
\expandafter\def\csname ccval@reg.2026-09-27-lineup.recorded\endcsname{none kept}
\expandafter\def\csname ccval@reg.count\endcsname{9}
\expandafter\def\csname ccval@regrade.detect.A-27b.fixtures\endcsname{13}
\expandafter\def\csname ccval@regrade.detect.A-27b.fixtures_measured\endcsname{12}
\expandafter\def\csname ccval@regrade.detect.A-27b.invented_on_unmeasured\endcsname{0}
\expandafter\def\csname ccval@regrade.detect.A-27b.invented_total\endcsname{2}
\expandafter\def\csname ccval@regrade.detect.A-27b.matched\endcsname{11}
\expandafter\def\csname ccval@regrade.detect.A-27b.missed\endcsname{1}
\expandafter\def\csname ccval@regrade.detect.A-27b.model\endcsname{qwen3.8:27b}
\expandafter\def\csname ccval@regrade.detect.A-27b.plants_detected\endcsname{6}
\expandafter\def\csname ccval@regrade.detect.A-27b.plants_missed\endcsname{0}
\expandafter\def\csname ccval@regrade.detect.A-27b.plants_total\endcsname{6}
\expandafter\def\csname ccval@regrade.detect.A-27b.plants_unmeasured\endcsname{0}
\expandafter\def\csname ccval@regrade.detect.A-27b.unmeasured\endcsname{1}
\expandafter\def\csname ccval@regrade.detect.B-70b.fixtures\endcsname{13}
\expandafter\def\csname ccval@regrade.detect.B-70b.fixtures_measured\endcsname{11}
\expandafter\def\csname ccval@regrade.detect.B-70b.invented_on_unmeasured\endcsname{0}
\expandafter\def\csname ccval@regrade.detect.B-70b.invented_total\endcsname{1}
\expandafter\def\csname ccval@regrade.detect.B-70b.matched\endcsname{11}
\expandafter\def\csname ccval@regrade.detect.B-70b.missed\endcsname{0}
\expandafter\def\csname ccval@regrade.detect.B-70b.model\endcsname{deepseek-r1:70b}
\expandafter\def\csname ccval@regrade.detect.B-70b.plants_detected\endcsname{5}
\expandafter\def\csname ccval@regrade.detect.B-70b.plants_missed\endcsname{0}
\expandafter\def\csname ccval@regrade.detect.B-70b.plants_total\endcsname{6}
\expandafter\def\csname ccval@regrade.detect.B-70b.plants_unmeasured\endcsname{1}
\expandafter\def\csname ccval@regrade.detect.B-70b.unmeasured\endcsname{2}
\expandafter\def\csname ccval@regrade.detect.C-120b.fixtures\endcsname{13}
\expandafter\def\csname ccval@regrade.detect.C-120b.fixtures_measured\endcsname{12}
\expandafter\def\csname ccval@regrade.detect.C-120b.invented_on_unmeasured\endcsname{1}
\expandafter\def\csname ccval@regrade.detect.C-120b.invented_total\endcsname{3}
\expandafter\def\csname ccval@regrade.detect.C-120b.matched\endcsname{12}
\expandafter\def\csname ccval@regrade.detect.C-120b.missed\endcsname{0}
\expandafter\def\csname ccval@regrade.detect.C-120b.model\endcsname{gpt-oss:120b}
\expandafter\def\csname ccval@regrade.detect.C-120b.plants_detected\endcsname{6}
\expandafter\def\csname ccval@regrade.detect.C-120b.plants_missed\endcsname{0}
\expandafter\def\csname ccval@regrade.detect.C-120b.plants_total\endcsname{6}
\expandafter\def\csname ccval@regrade.detect.C-120b.plants_unmeasured\endcsname{0}
\expandafter\def\csname ccval@regrade.detect.C-120b.unmeasured\endcsname{1}
\expandafter\def\csname ccval@regrade.detect.amended\endcsname{2026-09-03}
\expandafter\def\csname ccval@regrade.guard-probes.A-27b.fixtures\endcsname{[{'fixture': 'attribution_invented', 'guard_probes_total': 6, 'guard_confirmed': ['a-listen-port', 'a-read-timeout', 'a-health-path', 'b-tls-min', 'b-max-connections', 'm-read-timeout'], 'guard_wrong': [], 'guard_unclaimed': [], 'detail': [{'probe_id': 'a-listen-port', 'expected_verdict': 'SUPPORTED', 'anchor': {'all_of': ['listens on port 8443']}, 'verdict': 'confirmed', 'claim_id': 'A-001', 'model_verdict': 'SUPPORTED', 'evidence': 'The relay listens on port 8443.', 'rationale': 'The reference text explicitly states that the relay listens on port 8443.'}, {'probe_id': 'a-read-timeout', 'expected_verdict': 'SUPPORTED', 'anchor': {'all_of': ['read timeout is 30 seconds']}, 'verdict': 'confirmed', 'claim_id': 'A-002', 'model_verdict': 'SUPPORTED', 'evidence': 'The read timeout is 30 seconds.', 'rationale': 'The reference text explicitly states that the read timeout is 30 seconds.'}, {'probe_id': 'a-health-path', 'expected_verdict': 'SUPPORTED', 'anchor': {'all_of': ['/healthz']}, 'verdict': 'confirmed', 'claim_id': 'A-003', 'model_verdict': 'SUPPORTED', 'evidence': 'The health check path is /healthz.', 'rationale': 'The reference text explicitly states that the health check path is /healthz.'}, {'probe_id': 'b-tls-min', 'expected_verdict': 'SUPPORTED', 'anchor': {'all_of': ['minimum TLS version is 1.2']}, 'verdict': 'confirmed', 'claim_id': 'B-002', 'model_verdict': 'SUPPORTED', 'evidence': 'According to the Operator Guide, the minimum TLS version is 1.2.', 'rationale': 'The reference text states that the minimum TLS version is 1.2, attributed to the Operator Guide.'}, {'probe_id': 'b-max-connections', 'expected_verdict': 'SUPPORTED', 'anchor': {'all_of': ['512']}, 'verdict': 'confirmed', 'claim_id': 'B-003', 'model_verdict': 'SUPPORTED', 'evidence': 'The maximum concurrent connections is 512.', 'rationale': 'The reference text explicitly states that the maximum concurrent connections is 512.'}, {'probe_id': 'm-read-timeout', 'expected_verdict': 'SUPPORTED', 'anchor': {'all_of': ['read timeout is 30 seconds']}, 'verdict': 'confirmed', 'claim_id': 'M-002', 'model_verdict': 'SUPPORTED', 'evidence': 'The read timeout is 30 seconds.', 'rationale': 'The claim is explicitly stated in source_a.md.'}]}, {'fixture': 'attribution_swapped', 'guard_probes_total': 2, 'guard_confirmed': ['a-read-timeout', 'b-read-timeout'], 'guard_wrong': [], 'guard_unclaimed': [], 'detail': [{'probe_id': 'a-read-timeout', 'expected_verdict': 'SUPPORTED', 'anchor': {'all_of': ['read timeout is 30 seconds']}, 'verdict': 'confirmed', 'claim_id': 'A-002', 'model_verdict': 'SUPPORTED', 'evidence': 'According to the Deployment Notes, the read timeout is 30 seconds.', 'rationale': 'The reference text attributes a read timeout of 30 seconds to the Deployment Notes, which supports the claim.'}, {'probe_id': 'b-read-timeout', 'expected_verdict': 'SUPPORTED', 'anchor': {'all_of': ['read timeout is 60 seconds']}, 'verdict': 'confirmed', 'claim_id': 'B-002', 'model_verdict': 'SUPPORTED', 'evidence': 'According to the Operator Guide, the read timeout is 60 seconds.', 'rationale': 'The reference text attributes a read timeout of 60 seconds to the Operator Guide, which supports the claim.'}]}, {'fixture': 'conflict_surfaced', 'guard_probes_total': 11, 'guard_confirmed': ['a-connect-timeout', 'b-connect-timeout', 'a-listen-port', 'a-read-timeout', 'a-max-connections', 'a-tls-version', 'b-log-format', 'b-health-path', 'merged-attributed-timeout-a', 'merged-attributed-timeout-b', 'merged-max-connections'], 'guard_wrong': [], 'guard_unclaimed': [], 'detail': [{'probe_id': 'a-connect-timeout', 'expected_verdict': 'SUPPORTED', 'anchor': {'all_of': ['connect timeout', '30 seconds']}, 'verdict': 'confirmed', 'claim_id': 'A-003', 'model_verdict': 'SUPPORTED', 'evidence': 'According to the Operator Guide, the default connect timeout is 30 seconds.', 'rationale': 'The reference text attributes the value of 30 seconds to the Operator Guide, which supports the claim as one of the stated values.'}, {'probe_id': 'b-connect-timeout', 'expected_verdict': 'SUPPORTED', 'anchor': {'all_of': ['connect timeout', '60 seconds']}, 'verdict': 'confirmed', 'claim_id': 'B-003', 'model_verdict': 'SUPPORTED', 'evidence': 'According to the Deployment Notes, the default connect timeout is 60 seconds.', 'rationale': 'The reference text attributes the value of 60 seconds to the Deployment Notes, which supports the claim as one of the stated values.'}, {'probe_id': 'a-listen-port', 'expected_verdict': 'SUPPORTED', 'anchor': {'all_of': ['listens on port', '8443']}, 'verdict': 'confirmed', 'claim_id': 'A-002', 'model_verdict': 'SUPPORTED', 'evidence': 'The relay listens on port 8443 by default.', 'rationale': 'The reference text explicitly states that the relay listens on port 8443 by default.'}, {'probe_id': 'a-read-timeout', 'expected_verdict': 'SUPPORTED', 'anchor': {'all_of': ['read timeout', '120 seconds']}, 'verdict': 'confirmed', 'claim_id': 'A-004', 'model_verdict': 'SUPPORTED', 'evidence': 'The default read timeout is 120 seconds.', 'rationale': 'The reference text explicitly states that the default read timeout is 120 seconds.'}, {'probe_id': 'a-max-connections', 'expected_verdict': 'SUPPORTED', 'anchor': {'all_of': ['512 concurrent connections']}, 'verdict': 'confirmed', 'claim_id': 'A-005', 'model_verdict': 'SUPPORTED', 'evidence': 'The relay accepts at most 512 concurrent connections.', 'rationale': 'The reference text explicitly states that the relay accepts at most 512 concurrent connections.'}, {'probe_id': 'a-tls-version', 'expected_verdict': 'SUPPORTED', 'anchor': {'all_of': ['minimum accepted TLS version', 'TLS 1.2']}, 'verdict': 'confirmed', 'claim_id': 'A-007', 'model_verdict': 'SUPPORTED', 'evidence': 'The minimum accepted TLS version is TLS 1.2.', 'rationale': 'The reference text explicitly states that the minimum accepted TLS version is TLS 1.2.'}, {'probe_id': 'b-log-format', 'expected_verdict': 'SUPPORTED', 'anchor': {'all_of': ['access logs', 'JSON Lines']}, 'verdict': 'confirmed', 'claim_id': 'B-005', 'model_verdict': 'SUPPORTED', 'evidence': 'Access logs are written in JSON Lines format.', 'rationale': 'The reference text explicitly states that access logs are written in JSON Lines format.'}, {'probe_id': 'b-health-path', 'expected_verdict': 'SUPPORTED', 'anchor': {'all_of': ['health check endpoint', '/-/healthy']}, 'verdict': 'confirmed', 'claim_id': 'B-006', 'model_verdict': 'SUPPORTED', 'evidence': 'The health check endpoint is /-/healthy.', 'rationale': 'The reference text explicitly states that the health check endpoint is /-/healthy.'}, {'probe_id': 'merged-attributed-timeout-a', 'expected_verdict': 'SUPPORTED', 'anchor': {'all_of': ['According to the Operator Guide', '30 seconds']}, 'verdict': 'confirmed', 'claim_id': 'M-004', 'model_verdict': 'SUPPORTED', 'evidence': 'The default connect timeout is 30 seconds.', 'rationale': "Source A, titled 'Vandrell Relay — Operator Guide', states the default connect timeout is 30 seconds."}, {'probe_id': 'merged-attributed-timeout-b', 'expected_verdict': 'SUPPORTED', 'anchor': {'all_of': ['According to the Deployment Notes', '60 seconds']}, 'verdict': 'confirmed', 'claim_id': 'M-005', 'model_verdict': 'SUPPORTED', 'evidence': 'The default connect timeout is 60 seconds.', 'rationale': "Source B, titled 'Vandrell Relay — Deployment Notes', states the default connect timeout is 60 seconds."}, {'probe_id': 'merged-max-connections', 'expected_verdict': 'SUPPORTED', 'anchor': {'all_of': ['512 concurrent connections']}, 'verdict': 'confirmed', 'claim_id': 'M-006', 'model_verdict': 'SUPPORTED', 'evidence': 'The relay accepts at most 512 concurrent connections.', 'rationale': 'Source A states that the relay accepts at most 512 concurrent connections.'}]}, {'fixture': 'contradiction', 'guard_probes_total': 7, 'guard_confirmed': ['b-connect-timeout', 'a-listen-port', 'a-read-timeout', 'a-max-connections', 'a-tls-version', 'b-log-format', 'b-health-path'], 'guard_wrong': [], 'guard_unclaimed': [], 'detail': [{'probe_id': 'b-connect-timeout', 'expected_verdict': 'SUPPORTED', 'anchor': {'all_of': ['connect timeout', '60 seconds']}, 'verdict': 'confirmed', 'claim_id': 'B-003', 'model_verdict': 'SUPPORTED', 'evidence': 'The default connect timeout is 60 seconds.', 'rationale': 'The reference text explicitly states that the default connect timeout is 60 seconds.'}, {'probe_id': 'a-listen-port', 'expected_verdict': 'SUPPORTED', 'anchor': {'all_of': ['listens on port', '8443']}, 'verdict': 'confirmed', 'claim_id': 'A-002', 'model_verdict': 'SUPPORTED', 'evidence': 'The relay listens on port 8443 by default.', 'rationale': 'The reference text explicitly states that the relay listens on port 8443 by default.'}, {'probe_id': 'a-read-timeout', 'expected_verdict': 'SUPPORTED', 'anchor': {'all_of': ['read timeout', '120 seconds']}, 'verdict': 'confirmed', 'claim_id': 'A-004', 'model_verdict': 'SUPPORTED', 'evidence': 'The default read timeout is 120 seconds.', 'rationale': 'The reference text explicitly states that the default read timeout is 120 seconds.'}, {'probe_id': 'a-max-connections', 'expected_verdict': 'SUPPORTED', 'anchor': {'all_of': ['512 concurrent connections']}, 'verdict': 'confirmed', 'claim_id': 'A-005', 'model_verdict': 'SUPPORTED', 'evidence': 'The relay accepts at most 512 concurrent connections.', 'rationale': 'The reference text explicitly states that the relay accepts at most 512 concurrent connections.'}, {'probe_id': 'a-tls-version', 'expected_verdict': 'SUPPORTED', 'anchor': {'all_of': ['minimum accepted TLS version', 'TLS 1.2']}, 'verdict': 'confirmed', 'claim_id': 'A-007', 'model_verdict': 'SUPPORTED', 'evidence': 'The minimum accepted TLS version is TLS 1.2.', 'rationale': 'The reference text explicitly states that the minimum accepted TLS version is TLS 1.2.'}, {'probe_id': 'b-log-format', 'expected_verdict': 'SUPPORTED', 'anchor': {'all_of': ['access logs', 'JSON Lines']}, 'verdict': 'confirmed', 'claim_id': 'B-005', 'model_verdict': 'SUPPORTED', 'evidence': 'Access logs are written in JSON Lines format.', 'rationale': 'The reference text explicitly states that access logs are written in JSON Lines format.'}, {'probe_id': 'b-health-path', 'expected_verdict': 'SUPPORTED', 'anchor': {'all_of': ['health check endpoint', '/-/healthy']}, 'verdict': 'confirmed', 'claim_id': 'B-006', 'model_verdict': 'SUPPORTED', 'evidence': 'The health check endpoint is /-/healthy.', 'rationale': 'The reference text explicitly states that the health check endpoint is /-/healthy.'}]}, {'fixture': 'dedup', 'guard_probes_total': 10, 'guard_confirmed': ['a-listen-port', 'a-health-path', 'a-read-timeout', 'b-listen-port', 'b-health-path', 'b-tls-min', 'm-listen-port', 'm-health-path', 'm-read-timeout', 'm-tls-min'], 'guard_wrong': [], 'guard_unclaimed': [], 'detail': [{'probe_id': 'a-listen-port', 'expected_verdict': 'SUPPORTED', 'anchor': {'all_of': ['listens on port 8443']}, 'verdict': 'confirmed', 'claim_id': 'A-001', 'model_verdict': 'SUPPORTED', 'evidence': 'The relay listens on port 8443.', 'rationale': 'The reference text explicitly states that the relay listens on port 8443.'}, {'probe_id': 'a-health-path', 'expected_verdict': 'SUPPORTED', 'anchor': {'all_of': ['/healthz']}, 'verdict': 'confirmed', 'claim_id': 'A-002', 'model_verdict': 'SUPPORTED', 'evidence': 'The health check path is /healthz.', 'rationale': 'The reference text explicitly states that the health check path is /healthz.'}, {'probe_id': 'a-read-timeout', 'expected_verdict': 'SUPPORTED', 'anchor': {'all_of': ['read timeout is 30 seconds']}, 'verdict': 'confirmed', 'claim_id': 'A-003', 'model_verdict': 'SUPPORTED', 'evidence': 'The read timeout is 30 seconds.', 'rationale': 'The reference text explicitly states that the read timeout is 30 seconds.'}, {'probe_id': 'b-listen-port', 'expected_verdict': 'SUPPORTED', 'anchor': {'all_of': ['listens on port 8443']}, 'verdict': 'confirmed', 'claim_id': 'A-001', 'model_verdict': 'SUPPORTED', 'evidence': 'The relay listens on port 8443.', 'rationale': 'The reference text explicitly states that the relay listens on port 8443.'}, {'probe_id': 'b-health-path', 'expected_verdict': 'SUPPORTED', 'anchor': {'all_of': ['/healthz']}, 'verdict': 'confirmed', 'claim_id': 'A-002', 'model_verdict': 'SUPPORTED', 'evidence': 'The health check path is /healthz.', 'rationale': 'The reference text explicitly states that the health check path is /healthz.'}, {'probe_id': 'b-tls-min', 'expected_verdict': 'SUPPORTED', 'anchor': {'all_of': ['minimum TLS version is 1.2']}, 'verdict': 'confirmed', 'claim_id': 'B-003', 'model_verdict': 'SUPPORTED', 'evidence': 'The minimum TLS version is 1.2.', 'rationale': 'The reference text explicitly states that the minimum TLS version is 1.2.'}, {'probe_id': 'm-listen-port', 'expected_verdict': 'SUPPORTED', 'anchor': {'all_of': ['listens on port 8443']}, 'verdict': 'confirmed', 'claim_id': 'M-001', 'model_verdict': 'SUPPORTED', 'evidence': 'The relay listens on port 8443.', 'rationale': 'The claim is explicitly stated in source_a.md.'}, {'probe_id': 'm-health-path', 'expected_verdict': 'SUPPORTED', 'anchor': {'all_of': ['/healthz']}, 'verdict': 'confirmed', 'claim_id': 'M-002', 'model_verdict': 'SUPPORTED', 'evidence': 'The health check path is /healthz.', 'rationale': 'The claim is explicitly stated in source_a.md.'}, {'probe_id': 'm-read-timeout', 'expected_verdict': 'SUPPORTED', 'anchor': {'all_of': ['read timeout is 30 seconds']}, 'verdict': 'confirmed', 'claim_id': 'M-003', 'model_verdict': 'SUPPORTED', 'evidence': 'The read timeout is 30 seconds.', 'rationale': 'The claim is explicitly stated in source_a.md.'}, {'probe_id': 'm-tls-min', 'expected_verdict': 'SUPPORTED', 'anchor': {'all_of': ['minimum TLS version is 1.2']}, 'verdict': 'confirmed', 'claim_id': 'M-004', 'model_verdict': 'SUPPORTED', 'evidence': 'The minimum TLS version is 1.2.', 'rationale': 'The claim is explicitly stated in source_b.md.'}]}, {'fixture': 'disjoint_domains', 'guard_probes_total': 16, 'guard_confirmed': ['a-round-hours', 'a-second-round-verbatim', 'a-responsibility', 'a-contractor-dates', 'a-tie-officer', 'b-closed-summer', 'b-boards-left-up', 'b-ladder-not-rating', 'b-ladder-moves', 'b-spare-clocks', 'b-extra-table', 'm-round-hours-reverse', 'm-club-season-reverse', 'm-contractor-dates-reverse', 'm-ladder-reverse'], 'guard_wrong': [], 'guard_unclaimed': ['a-frost-latch'], 'detail': [{'probe_id': 'a-round-hours', 'expected_verdict': 'SUPPORTED', 'anchor': {'all_of': ['7 am', '8 am']}, 'verdict': 'confirmed', 'claim_id': 'A-003', 'model_verdict': 'SUPPORTED', 'evidence': 'The duty officer walks the north corridor every weekday from 7 am till 8 am', 'rationale': 'The reference text confirms the duty officer walks the north corridor every weekday from 7 am till 8 am.'}, {'probe_id': 'a-second-round-verbatim', 'expected_verdict': 'SUPPORTED', 'anchor': {'all_of': ['13:05', '16:20']}, 'verdict': 'confirmed', 'claim_id': 'A-005', 'model_verdict': 'SUPPORTED', 'evidence': 'A second round covers the workshop and the paint store from 13:05 till 16:20 hours on Wednesdays only', 'rationale': 'The reference text explicitly describes the second round covering the workshop and paint store on Wednesdays from 13:05 to 16:20.'}, {'probe_id': 'a-frost-latch', 'expected_verdict': 'SUPPORTED', 'anchor': {'all_of': ['frost', 'November']}, 'verdict': 'unclaimed', 'claim_id': None, 'model_verdict': None}, {'probe_id': 'a-responsibility', 'expected_verdict': 'SUPPORTED', 'anchor': {'all_of': ['responsibility of the site']}, 'verdict': 'confirmed', 'claim_id': 'A-011', 'model_verdict': 'SUPPORTED', 'evidence': 'the doors are the responsibility of the site', 'rationale': 'The reference text states that the doors are the responsibility of the site.'}, {'probe_id': 'a-contractor-dates', 'expected_verdict': 'SUPPORTED', 'anchor': {'all_of': ['March 1st', 'March 31st']}, 'verdict': 'confirmed', 'claim_id': 'A-013', 'model_verdict': 'SUPPORTED', 'evidence': 'attends every March 1st to March 31st', 'rationale': 'The reference text states the contractor attends every March 1st to March 31st.'}, {'probe_id': 'a-tie-officer', 'expected_verdict': 'SUPPORTED', 'anchor': {'all_of': ['one officer clears both']}, 'verdict': 'confirmed', 'claim_id': 'A-022', 'model_verdict': 'SUPPORTED', 'evidence': 'one officer clears both at the end of the week', 'rationale': 'The reference text states that one officer clears both the card and the written record at the end of the week.'}, {'probe_id': 'b-closed-summer', 'expected_verdict': 'SUPPORTED', 'anchor': {'all_of': ['does not meet', 'July', 'September']}, 'verdict': 'confirmed', 'claim_id': 'B-003', 'model_verdict': 'SUPPORTED', 'evidence': 'The club does not meet between mid July and early September', 'rationale': 'The reference text states the club does not meet between mid July and early September.'}, {'probe_id': 'b-boards-left-up', 'expected_verdict': 'SUPPORTED', 'anchor': {'all_of': ['boards set up']}, 'verdict': 'confirmed', 'claim_id': 'B-007', 'model_verdict': 'SUPPORTED', 'evidence': 'the members have learned to leave the boards set up rather than pack them away', 'rationale': 'The text states that members have learned to leave the boards set up rather than pack them away.'}, {'probe_id': 'b-ladder-not-rating', 'expected_verdict': 'SUPPORTED', 'anchor': {'all_of': ['simple ladder']}, 'verdict': 'confirmed', 'claim_id': 'B-008', 'model_verdict': 'SUPPORTED', 'evidence': 'the club runs a simple ladder that anyone may join at any point', 'rationale': 'The text states the club runs a simple ladder that anyone may join at any point.'}, {'probe_id': 'b-ladder-moves', 'expected_verdict': 'SUPPORTED', 'anchor': {'all_of': ['up one place', 'down one place']}, 'verdict': 'confirmed', 'claim_id': 'B-009', 'model_verdict': 'SUPPORTED', 'evidence': 'Members who win a game move up one place', 'rationale': 'The text states that members who win a game move up one place.'}, {'probe_id': 'b-spare-clocks', 'expected_verdict': 'SUPPORTED', 'anchor': {'all_of': ['spare clocks']}, 'verdict': 'confirmed', 'claim_id': 'B-025', 'model_verdict': 'SUPPORTED', 'evidence': 'the county side gave the club enough spare clocks to run every board on the ladder at the same time', 'rationale': 'The text states the county side gave the club enough spare clocks to run every board on the ladder at the same time.'}, {'probe_id': 'b-extra-table', 'expected_verdict': 'SUPPORTED', 'anchor': {'all_of': ['extra table', 'Tuesday']}, 'verdict': 'confirmed', 'claim_id': 'B-029', 'model_verdict': 'SUPPORTED', 'evidence': 'Eventually the club grew to thirty members', 'rationale': 'The reference text explicitly states that the club grew to thirty members.'}, {'probe_id': 'm-round-hours-reverse', 'expected_verdict': 'SUPPORTED', 'anchor': {'all_of': ['7 am', '8 am']}, 'verdict': 'confirmed', 'claim_id': 'M-003', 'model_verdict': 'SUPPORTED', 'evidence': 'The duty officer walks the north corridor every weekday from 7 am till 8 am', 'rationale': 'The source confirms the duty officer walks the north corridor every weekday from 7 am till 8 am.'}, {'probe_id': 'm-club-season-reverse', 'expected_verdict': 'SUPPORTED', 'anchor': {'all_of': ['does not meet', 'July', 'September']}, 'verdict': 'confirmed', 'claim_id': 'M-024', 'model_verdict': 'SUPPORTED', 'evidence': 'The club does not meet between mid July and early September', 'rationale': 'The source states the club does not meet between mid July and early September.'}, {'probe_id': 'm-contractor-dates-reverse', 'expected_verdict': 'SUPPORTED', 'anchor': {'all_of': ['March 1st', 'March 31st']}, 'verdict': 'confirmed', 'claim_id': 'M-013', 'model_verdict': 'SUPPORTED', 'evidence': 'attends every March 1st to March 31st', 'rationale': 'The source states the contractor attends every March 1st to March 31st.'}, {'probe_id': 'm-ladder-reverse', 'expected_verdict': 'SUPPORTED', 'anchor': {'all_of': ['simple ladder']}, 'verdict': 'confirmed', 'claim_id': 'M-029', 'model_verdict': 'SUPPORTED', 'evidence': 'Instead of a rating system, the club runs a simple ladder that anyone may join at any point.', 'rationale': 'The source explicitly states that the club runs a simple ladder open to anyone at any point.'}]}, {'fixture': 'disjoint_sources', 'guard_probes_total': 16, 'guard_confirmed': ['a-visit-hours', 'a-work-hours-verbatim', 'a-winter-dormant', 'a-ownership', 'a-public-dates', 'a-postcard-tie', 'b-not-traversable', 'b-lady-year-round', 'b-herbs-not-magic', 'b-herbs-to-cities', 'b-medicine-village', 'b-cheese-autumn', 'm-visit-hours-reverse', 'm-forest-season-reverse', 'm-public-dates-reverse', 'm-herbs-reverse'], 'guard_wrong': [], 'guard_unclaimed': [], 'detail': [{'probe_id': 'a-visit-hours', 'expected_verdict': 'SUPPORTED', 'anchor': {'all_of': ['8 am', '9 am']}, 'verdict': 'confirmed', 'claim_id': 'A-001', 'model_verdict': 'SUPPORTED', 'evidence': 'She visited the rosegarden every morning from 8 am till 9 am', 'rationale': 'The reference text explicitly states that she visited the rosegarden every morning from 8 am till 9 am.'}, {'probe_id': 'a-work-hours-verbatim', 'expected_verdict': 'SUPPORTED', 'anchor': {'all_of': ['02:00', '15:30']}, 'verdict': 'confirmed', 'claim_id': 'A-003', 'model_verdict': 'SUPPORTED', 'evidence': "From 02:00 till 15:30 o'clock, she worked in the rosegarden to tidy it up", 'rationale': 'The reference text confirms she worked in the rosegarden to tidy it up during the specified time.'}, {'probe_id': 'a-winter-dormant', 'expected_verdict': 'SUPPORTED', 'anchor': {'all_of': ['winter', 'rosegarden']}, 'verdict': 'confirmed', 'claim_id': 'A-004', 'model_verdict': 'SUPPORTED', 'evidence': 'The rosegarden lost all his life in winter', 'rationale': 'The reference text states that the rosegarden lost all its life in winter.'}, {'probe_id': 'a-ownership', 'expected_verdict': 'SUPPORTED', 'anchor': {'all_of': ['private property', 'royal family']}, 'verdict': 'confirmed', 'claim_id': 'A-007', 'model_verdict': 'SUPPORTED', 'evidence': 'the rose garden was the private property of the royal family', 'rationale': 'The reference text explicitly states that the rose garden was the private property of the royal family.'}, {'probe_id': 'a-public-dates', 'expected_verdict': 'SUPPORTED', 'anchor': {'all_of': ['March 1st', 'March 31st']}, 'verdict': 'confirmed', 'claim_id': 'A-007', 'model_verdict': 'SUPPORTED', 'evidence': 'the rose garden was the private property of the royal family', 'rationale': 'The reference text explicitly states that the rose garden was the private property of the royal family.'}, {'probe_id': 'a-postcard-tie', 'expected_verdict': 'SUPPORTED', 'anchor': {'all_of': ['same carriage']}, 'verdict': 'confirmed', 'claim_id': 'A-013', 'model_verdict': 'SUPPORTED', 'evidence': 'In many cases, it was a tie because the people and the postcards were usually transported in the same carriage', 'rationale': 'The reference text states that it was often a tie because they were transported in the same carriage.'}, {'probe_id': 'b-not-traversable', 'expected_verdict': 'SUPPORTED', 'anchor': {'all_of': ['not traversable', 'October', 'April']}, 'verdict': 'confirmed', 'claim_id': 'B-001', 'model_verdict': 'SUPPORTED', 'evidence': 'The forest was not traversable between October and April', 'rationale': 'The reference text states that the forest was not traversable between October and April.'}, {'probe_id': 'b-lady-year-round', 'expected_verdict': 'SUPPORTED', 'anchor': {'all_of': ['entire year', 'forest']}, 'verdict': 'confirmed', 'claim_id': 'B-003', 'model_verdict': 'SUPPORTED', 'evidence': 'The old lady nevertheless lived the entire year in the forest', 'rationale': 'The reference text states that the old lady lived in the forest the entire year.'}, {'probe_id': 'b-herbs-not-magic', 'expected_verdict': 'SUPPORTED', 'anchor': {'all_of': ['wild herbs']}, 'verdict': 'confirmed', 'claim_id': 'B-007', 'model_verdict': 'SUPPORTED', 'evidence': 'she collected wild herbs', 'rationale': 'The reference text states that she collected wild herbs.'}, {'probe_id': 'b-herbs-to-cities', 'expected_verdict': 'SUPPORTED', 'anchor': {'all_of': ['travelers', 'cities']}, 'verdict': 'confirmed', 'claim_id': 'B-008', 'model_verdict': 'SUPPORTED', 'evidence': 'Those wild herbs were tried and sold to wandering travelers', 'rationale': 'The reference text states that the herbs were tried and sold to wandering travelers.'}, {'probe_id': 'b-medicine-village', 'expected_verdict': 'SUPPORTED', 'anchor': {'all_of': ['medicine', 'entire village']}, 'verdict': 'confirmed', 'claim_id': 'B-020', 'model_verdict': 'SUPPORTED', 'evidence': 'When the man recovered again, she gave him enough medicine made out of the herbs she collected in spring for the entire village.', 'rationale': 'The reference text explicitly states that she gave him enough medicine made from herbs collected in spring for the entire village.'}, {'probe_id': 'b-cheese-autumn', 'expected_verdict': 'SUPPORTED', 'anchor': {'all_of': ['cheese', 'autumn']}, 'verdict': 'confirmed', 'claim_id': 'B-024', 'model_verdict': 'SUPPORTED', 'evidence': 'Eventually, everyone got healthy again', 'rationale': 'The reference text explicitly states that everyone got healthy again.'}, {'probe_id': 'm-visit-hours-reverse', 'expected_verdict': 'SUPPORTED', 'anchor': {'all_of': ['8 am', '9 am']}, 'verdict': 'confirmed', 'claim_id': 'M-001', 'model_verdict': 'SUPPORTED', 'evidence': 'She visited the rosegarden every morning from 8 am till 9 am', 'rationale': 'The source explicitly states that the princess visited the rosegarden every morning from 8 am till 9 am.'}, {'probe_id': 'm-forest-season-reverse', 'expected_verdict': 'SUPPORTED', 'anchor': {'all_of': ['not traversable', 'October', 'April']}, 'verdict': 'confirmed', 'claim_id': 'M-014', 'model_verdict': 'SUPPORTED', 'evidence': 'The forest was not traversable between October and April', 'rationale': 'The source states the forest was not traversable between October and April.'}, {'probe_id': 'm-public-dates-reverse', 'expected_verdict': 'SUPPORTED', 'anchor': {'all_of': ['March 1st', 'March 31st']}, 'verdict': 'confirmed', 'claim_id': 'M-007', 'model_verdict': 'SUPPORTED', 'evidence': 'While the rose garden was the private property of the royal family', 'rationale': 'The source explicitly states the rose garden was the private property of the royal family.'}, {'probe_id': 'm-herbs-reverse', 'expected_verdict': 'SUPPORTED', 'anchor': {'all_of': ['wild herbs']}, 'verdict': 'confirmed', 'claim_id': 'M-019', 'model_verdict': 'SUPPORTED', 'evidence': 'Instead of harvesting skulls and bones, she collected wild herbs.', 'rationale': 'The source states the old lady collected wild herbs.'}]}, {'fixture': 'dropped_claim', 'guard_probes_total': 9, 'guard_confirmed': ['b-log-rotation', 'b-worker-per-core', 'b-health-path', 'b-health-failures', 'a-listen-port', 'a-connect-timeout', 'a-read-timeout', 'a-retry-limit', 'a-tls-version'], 'guard_wrong': [], 'guard_unclaimed': [], 'detail': [{'probe_id': 'b-log-rotation', 'expected_verdict': 'SUPPORTED', 'anchor': {'all_of': ['rotated', '100 megabytes']}, 'verdict': 'confirmed', 'claim_id': 'B-006', 'model_verdict': 'SUPPORTED', 'evidence': 'Log files are rotated when they reach 100 megabytes.', 'rationale': 'The reference text states that log files are rotated when they reach 100 megabytes.'}, {'probe_id': 'b-worker-per-core', 'expected_verdict': 'SUPPORTED', 'anchor': {'all_of': ['one worker process', 'per CPU core']}, 'verdict': 'confirmed', 'claim_id': 'B-004', 'model_verdict': 'SUPPORTED', 'evidence': 'One worker process is started per CPU core.', 'rationale': 'The reference text states that one worker process is started per CPU core.'}, {'probe_id': 'b-health-path', 'expected_verdict': 'SUPPORTED', 'anchor': {'all_of': ['health check endpoint', '/-/healthy']}, 'verdict': 'confirmed', 'claim_id': 'B-007', 'model_verdict': 'SUPPORTED', 'evidence': 'The health check endpoint is /-/healthy.', 'rationale': 'The reference text states that the health check endpoint is /-/healthy.'}, {'probe_id': 'b-health-failures', 'expected_verdict': 'SUPPORTED', 'anchor': {'all_of': ['3 consecutive health checks', 'removed from rotation']}, 'verdict': 'confirmed', 'claim_id': 'B-008', 'model_verdict': 'SUPPORTED', 'evidence': 'A relay that fails 3 consecutive health checks is removed from rotation.', 'rationale': 'The reference text states that a relay failing 3 consecutive health checks is removed from rotation.'}, {'probe_id': 'a-listen-port', 'expected_verdict': 'SUPPORTED', 'anchor': {'all_of': ['listens on port', '8443']}, 'verdict': 'confirmed', 'claim_id': 'A-002', 'model_verdict': 'SUPPORTED', 'evidence': 'The relay listens on port 8443 by default.', 'rationale': 'The reference text states that the relay listens on port 8443 by default.'}, {'probe_id': 'a-connect-timeout', 'expected_verdict': 'SUPPORTED', 'anchor': {'all_of': ['connect timeout', '30 seconds']}, 'verdict': 'confirmed', 'claim_id': 'A-003', 'model_verdict': 'SUPPORTED', 'evidence': 'The default connect timeout is 30 seconds.', 'rationale': 'The reference text states that the default connect timeout is 30 seconds.'}, {'probe_id': 'a-read-timeout', 'expected_verdict': 'SUPPORTED', 'anchor': {'all_of': ['read timeout', '120 seconds']}, 'verdict': 'confirmed', 'claim_id': 'A-004', 'model_verdict': 'SUPPORTED', 'evidence': 'The default read timeout is 120 seconds.', 'rationale': 'The reference text states that the default read timeout is 120 seconds.'}, {'probe_id': 'a-retry-limit', 'expected_verdict': 'SUPPORTED', 'anchor': {'all_of': ['retried at most 3 times']}, 'verdict': 'confirmed', 'claim_id': 'A-006', 'model_verdict': 'SUPPORTED', 'evidence': 'Each request is retried at most 3 times.', 'rationale': 'The reference text states that each request is retried at most 3 times.'}, {'probe_id': 'a-tls-version', 'expected_verdict': 'SUPPORTED', 'anchor': {'all_of': ['minimum accepted TLS version', 'TLS 1.2']}, 'verdict': 'confirmed', 'claim_id': 'A-007', 'model_verdict': 'SUPPORTED', 'evidence': 'The minimum accepted TLS version is TLS 1.2.', 'rationale': 'The reference text states that the minimum accepted TLS version is TLS 1.2.'}]}, {'fixture': 'hallucination', 'guard_probes_total': 7, 'guard_confirmed': ['merged-connect-timeout', 'merged-max-connections', 'a-listen-port', 'a-tls-version', 'a-client-cert', 'b-log-format', 'b-health-path'], 'guard_wrong': [], 'guard_unclaimed': [], 'detail': [{'probe_id': 'merged-connect-timeout', 'expected_verdict': 'SUPPORTED', 'anchor': {'all_of': ['connect timeout', '30 seconds']}, 'verdict': 'confirmed', 'claim_id': 'M-003', 'model_verdict': 'SUPPORTED', 'evidence': 'The default connect timeout is 30 seconds.', 'rationale': 'The claim is stated verbatim in source_a.md.'}, {'probe_id': 'merged-max-connections', 'expected_verdict': 'SUPPORTED', 'anchor': {'all_of': ['512 concurrent connections']}, 'verdict': 'confirmed', 'claim_id': 'M-004', 'model_verdict': 'SUPPORTED', 'evidence': 'The relay accepts at most 512 concurrent connections.', 'rationale': 'The claim is stated verbatim in source_b.md.'}, {'probe_id': 'a-listen-port', 'expected_verdict': 'SUPPORTED', 'anchor': {'all_of': ['listens on port', '8443']}, 'verdict': 'confirmed', 'claim_id': 'A-002', 'model_verdict': 'SUPPORTED', 'evidence': 'The relay listens on port 8443 by default.', 'rationale': 'The reference text explicitly states that the relay listens on port 8443 by default.'}, {'probe_id': 'a-tls-version', 'expected_verdict': 'SUPPORTED', 'anchor': {'all_of': ['minimum accepted TLS version', 'TLS 1.2']}, 'verdict': 'confirmed', 'claim_id': 'A-004', 'model_verdict': 'SUPPORTED', 'evidence': 'The minimum accepted TLS version is TLS 1.2.', 'rationale': 'The reference text explicitly states that the minimum accepted TLS version is TLS 1.2.'}, {'probe_id': 'a-client-cert', 'expected_verdict': 'SUPPORTED', 'anchor': {'all_of': ['client certificate verification', 'disabled by default']}, 'verdict': 'confirmed', 'claim_id': 'A-005', 'model_verdict': 'SUPPORTED', 'evidence': 'Client certificate verification is disabled by default.', 'rationale': 'The reference text explicitly states that client certificate verification is disabled by default.'}, {'probe_id': 'b-log-format', 'expected_verdict': 'SUPPORTED', 'anchor': {'all_of': ['access logs', 'JSON Lines']}, 'verdict': 'confirmed', 'claim_id': 'B-004', 'model_verdict': 'SUPPORTED', 'evidence': 'Access logs are written in JSON Lines format.', 'rationale': 'The reference text explicitly states that access logs are written in JSON Lines format.'}, {'probe_id': 'b-health-path', 'expected_verdict': 'SUPPORTED', 'anchor': {'all_of': ['health check endpoint', '/-/healthy']}, 'verdict': 'confirmed', 'claim_id': 'B-005', 'model_verdict': 'SUPPORTED', 'evidence': 'The health check endpoint is /-/healthy.', 'rationale': 'The reference text explicitly states that the health check endpoint is /-/healthy.'}]}, {'fixture': 'numeric_drift', 'guard_probes_total': 4, 'guard_confirmed': ['b-max-connections', 'a-retry-limit', 'b-listen-port', 'a-connect-timeout'], 'guard_wrong': [], 'guard_unclaimed': [], 'detail': [{'probe_id': 'b-max-connections', 'expected_verdict': 'SUPPORTED', 'anchor': {'all_of': ['connection limit', '512']}, 'verdict': 'confirmed', 'claim_id': 'B-001', 'model_verdict': 'SUPPORTED', 'evidence': 'The connection limit is 512 concurrent connections.', 'rationale': 'The reference text explicitly states the connection limit is 512 concurrent connections.'}, {'probe_id': 'a-retry-limit', 'expected_verdict': 'SUPPORTED', 'anchor': {'all_of': ['retried at most 3 times']}, 'verdict': 'confirmed', 'claim_id': 'A-003', 'model_verdict': 'SUPPORTED', 'evidence': 'Each request is retried at most 3 times.', 'rationale': 'The reference text states that each request is retried at most 3 times.'}, {'probe_id': 'b-listen-port', 'expected_verdict': 'SUPPORTED', 'anchor': {'all_of': ['listens on port', '8443']}, 'verdict': 'confirmed', 'claim_id': 'A-004', 'model_verdict': 'SUPPORTED', 'evidence': 'The relay listens on port 8443 by default.', 'rationale': 'The reference text states that the relay listens on port 8443 by default.'}, {'probe_id': 'a-connect-timeout', 'expected_verdict': 'SUPPORTED', 'anchor': {'all_of': ['connect timeout', '30 seconds']}, 'verdict': 'confirmed', 'claim_id': 'A-005', 'model_verdict': 'SUPPORTED', 'evidence': 'The default connect timeout is 30 seconds.', 'rationale': 'The reference text states that the default connect timeout is 30 seconds.'}]}, {'fixture': 'ordering_only', 'guard_probes_total': 11, 'guard_confirmed': ['a-listen-port', 'a-connect-timeout', 'a-max-connections', 'a-tls-version', 'b-log-format', 'b-retry-limit', 'b-read-timeout', 'b-health-path', 'merged-backoff', 'merged-client-cert', 'merged-log-rotation'], 'guard_wrong': [], 'guard_unclaimed': [], 'detail': [{'probe_id': 'a-listen-port', 'expected_verdict': 'SUPPORTED', 'anchor': {'all_of': ['listens on port', '8443']}, 'verdict': 'confirmed', 'claim_id': 'A-002', 'model_verdict': 'SUPPORTED', 'evidence': 'The relay listens on port 8443 by default.', 'rationale': 'The reference text states that the relay listens on port 8443 by default.'}, {'probe_id': 'a-connect-timeout', 'expected_verdict': 'SUPPORTED', 'anchor': {'all_of': ['connect timeout', '30 seconds']}, 'verdict': 'confirmed', 'claim_id': 'A-003', 'model_verdict': 'SUPPORTED', 'evidence': 'The default connect timeout is 30 seconds.', 'rationale': 'The reference text states that the default connect timeout is 30 seconds.'}, {'probe_id': 'a-max-connections', 'expected_verdict': 'SUPPORTED', 'anchor': {'all_of': ['512 concurrent connections']}, 'verdict': 'confirmed', 'claim_id': 'A-006', 'model_verdict': 'SUPPORTED', 'evidence': 'The relay accepts at most 512 concurrent connections.', 'rationale': 'The reference text states that the relay accepts at most 512 concurrent connections.'}, {'probe_id': 'a-tls-version', 'expected_verdict': 'SUPPORTED', 'anchor': {'all_of': ['minimum accepted TLS version', 'TLS 1.2']}, 'verdict': 'confirmed', 'claim_id': 'A-010', 'model_verdict': 'SUPPORTED', 'evidence': 'The minimum accepted TLS version is TLS 1.2.', 'rationale': 'The reference text states that the minimum accepted TLS version is TLS 1.2.'}, {'probe_id': 'b-log-format', 'expected_verdict': 'SUPPORTED', 'anchor': {'all_of': ['access logs', 'JSON Lines']}, 'verdict': 'confirmed', 'claim_id': 'A-012', 'model_verdict': 'SUPPORTED', 'evidence': 'Access logs are written in JSON Lines format.', 'rationale': 'The reference text states that access logs are written in JSON Lines format.'}, {'probe_id': 'b-retry-limit', 'expected_verdict': 'SUPPORTED', 'anchor': {'all_of': ['retried at most 3 times']}, 'verdict': 'confirmed', 'claim_id': 'A-007', 'model_verdict': 'SUPPORTED', 'evidence': 'Each request is retried at most 3 times.', 'rationale': 'The reference text states that each request is retried at most 3 times.'}, {'probe_id': 'b-read-timeout', 'expected_verdict': 'SUPPORTED', 'anchor': {'all_of': ['read timeout', '120 seconds']}, 'verdict': 'confirmed', 'claim_id': 'A-004', 'model_verdict': 'SUPPORTED', 'evidence': 'The default read timeout is 120 seconds.', 'rationale': 'The reference text states that the default read timeout is 120 seconds.'}, {'probe_id': 'b-health-path', 'expected_verdict': 'SUPPORTED', 'anchor': {'all_of': ['health check endpoint', '/-/healthy']}, 'verdict': 'confirmed', 'claim_id': 'A-005', 'model_verdict': 'SUPPORTED', 'evidence': 'The health check endpoint is /-/healthy.', 'rationale': 'The reference text states that the health check endpoint is /-/healthy.'}, {'probe_id': 'merged-backoff', 'expected_verdict': 'SUPPORTED', 'anchor': {'all_of': ['250 milliseconds', 'doubles']}, 'verdict': 'confirmed', 'claim_id': 'M-004', 'model_verdict': 'SUPPORTED', 'evidence': 'Retry backoff starts at 250 milliseconds', 'rationale': 'The claim is stated verbatim in the source document.'}, {'probe_id': 'merged-client-cert', 'expected_verdict': 'SUPPORTED', 'anchor': {'all_of': ['client certificate verification', 'disabled by default']}, 'verdict': 'confirmed', 'claim_id': 'M-013', 'model_verdict': 'SUPPORTED', 'evidence': 'Client certificate verification is disabled by default.', 'rationale': 'The claim is stated verbatim in the source document.'}, {'probe_id': 'merged-log-rotation', 'expected_verdict': 'SUPPORTED', 'anchor': {'all_of': ['rotated', '100 megabytes']}, 'verdict': 'confirmed', 'claim_id': 'M-007', 'model_verdict': 'SUPPORTED', 'evidence': 'Log files are rotated when they reach 100 megabytes.', 'rationale': 'The claim is stated verbatim in the source document.'}]}, {'fixture': 'paraphrase', 'guard_probes_total': 11, 'guard_confirmed': ['a-connect-timeout', 'a-read-timeout', 'a-max-connections', 'a-tls-version', 'b-retry-limit', 'b-log-format', 'b-health-path', 'merged-listen-port', 'merged-tls-version', 'merged-read-timeout'], 'guard_wrong': ['a-listen-port'], 'guard_unclaimed': [], 'detail': [{'probe_id': 'a-listen-port', 'expected_verdict': 'SUPPORTED', 'anchor': {'all_of': ['binds to port', '8443']}, 'verdict': 'false_alarm', 'claim_id': 'A-002', 'finding': 'partially_dropped', 'evidence': "The relay's default listen port is 8443.", 'rationale': "The reference text states the default port is 8443, but does not mention the condition 'unless overridden'."}, {'probe_id': 'a-connect-timeout', 'expected_verdict': 'SUPPORTED', 'anchor': {'all_of': ['30 seconds', 'abandoned']}, 'verdict': 'confirmed', 'claim_id': 'A-003', 'model_verdict': 'SUPPORTED', 'evidence': 'Connections that do not complete a handshake inside 30 seconds are dropped.', 'rationale': 'The reference text states that connections not completing a handshake within 30 seconds are dropped, which is semantically equivalent to being abandoned.'}, {'probe_id': 'a-read-timeout', 'expected_verdict': 'SUPPORTED', 'anchor': {'all_of': ['waits up to 120 seconds']}, 'verdict': 'confirmed', 'claim_id': 'A-004', 'model_verdict': 'SUPPORTED', 'evidence': 'The relay allows a response 120 seconds to arrive before it times out.', 'rationale': 'The reference text states the relay allows a response to arrive within 120 seconds before timing out, matching the claim.'}, {'probe_id': 'a-max-connections', 'expected_verdict': 'SUPPORTED', 'anchor': {'all_of': ['512 connections', 'same time']}, 'verdict': 'confirmed', 'claim_id': 'A-005', 'model_verdict': 'SUPPORTED', 'evidence': 'At most 512 connections may be held open concurrently.', 'rationale': 'The reference text states that at most 512 connections may be held open concurrently, which matches the claim.'}, {'probe_id': 'a-tls-version', 'expected_verdict': 'SUPPORTED', 'anchor': {'all_of': ['older than TLS 1.2', 'rejected']}, 'verdict': 'confirmed', 'claim_id': 'A-007', 'model_verdict': 'SUPPORTED', 'evidence': 'The relay will not negotiate any protocol version below TLS 1.2.', 'rationale': 'The reference text states the relay will not negotiate versions below TLS 1.2, meaning anything older is rejected.'}, {'probe_id': 'b-retry-limit', 'expected_verdict': 'SUPPORTED', 'anchor': {'all_of': ['Retry attempts per request', '3']}, 'verdict': 'confirmed', 'claim_id': 'B-006', 'model_verdict': 'SUPPORTED', 'evidence': 'The relay makes up to 3 further attempts after a request fails.', 'rationale': 'The reference text states the relay makes up to 3 further attempts, supporting the claim of 3 retry attempts.'}, {'probe_id': 'b-log-format', 'expected_verdict': 'SUPPORTED', 'anchor': {'all_of': ['JSON Lines', 'one record per line']}, 'verdict': 'confirmed', 'claim_id': 'B-008', 'model_verdict': 'SUPPORTED', 'evidence': 'Access log records are emitted one JSON object per line.', 'rationale': 'The reference text states access log records are one JSON object per line, which is the definition of JSON Lines format.'}, {'probe_id': 'b-health-path', 'expected_verdict': 'SUPPORTED', 'anchor': {'all_of': ['health checks', '/-/healthy']}, 'verdict': 'confirmed', 'claim_id': 'B-010', 'model_verdict': 'SUPPORTED', 'evidence': 'The endpoint /-/healthy answers health checks.', 'rationale': 'The reference text states the endpoint /-/healthy answers health checks, supporting the claim that health checks should target it.'}, {'probe_id': 'merged-listen-port', 'expected_verdict': 'SUPPORTED', 'anchor': {'all_of': ['default listen port', '8443']}, 'verdict': 'confirmed', 'claim_id': 'M-002', 'model_verdict': 'SUPPORTED', 'evidence': 'Unless you override it, the relay binds to port 8443.', 'rationale': 'Source A states the relay binds to port 8443 by default, and Source B lists the listen port as 8443.'}, {'probe_id': 'merged-tls-version', 'expected_verdict': 'SUPPORTED', 'anchor': {'all_of': ['will not negotiate', 'TLS 1.2']}, 'verdict': 'confirmed', 'claim_id': 'M-007', 'model_verdict': 'SUPPORTED', 'evidence': 'Anything older than TLS 1.2 is rejected.', 'rationale': 'Source A states that anything older than TLS 1.2 is rejected, which means the relay will not negotiate versions below TLS 1.2.'}, {'probe_id': 'merged-read-timeout', 'expected_verdict': 'SUPPORTED', 'anchor': {'all_of': ['120 seconds', 'times out']}, 'verdict': 'confirmed', 'claim_id': 'M-004', 'model_verdict': 'SUPPORTED', 'evidence': 'Once connected, the relay waits up to 120 seconds for a response.', 'rationale': 'Source A states the relay waits up to 120 seconds for a response, matching the claim about the timeout duration.'}]}, {'fixture': 'structure_added', 'guard_probes_total': 8, 'guard_confirmed': ['a-listen-port', 'a-read-timeout', 'b-health-path', 'b-tls-min', 'm-listen-port', 'm-read-timeout', 'm-health-path', 'm-tls-min'], 'guard_wrong': [], 'guard_unclaimed': [], 'detail': [{'probe_id': 'a-listen-port', 'expected_verdict': 'SUPPORTED', 'anchor': {'all_of': ['listens on port 8443']}, 'verdict': 'confirmed', 'claim_id': 'A-001', 'model_verdict': 'SUPPORTED', 'evidence': 'The relay listens on port 8443.', 'rationale': 'The reference text explicitly states that the relay listens on port 8443.'}, {'probe_id': 'a-read-timeout', 'expected_verdict': 'SUPPORTED', 'anchor': {'all_of': ['read timeout is 30 seconds']}, 'verdict': 'confirmed', 'claim_id': 'A-002', 'model_verdict': 'SUPPORTED', 'evidence': 'The read timeout is 30 seconds.', 'rationale': 'The reference text explicitly states that the read timeout is 30 seconds.'}, {'probe_id': 'b-health-path', 'expected_verdict': 'SUPPORTED', 'anchor': {'all_of': ['/healthz']}, 'verdict': 'confirmed', 'claim_id': 'B-002', 'model_verdict': 'SUPPORTED', 'evidence': 'The health check path is /healthz.', 'rationale': 'The reference text explicitly states that the health check path is /healthz.'}, {'probe_id': 'b-tls-min', 'expected_verdict': 'SUPPORTED', 'anchor': {'all_of': ['minimum TLS version is 1.2']}, 'verdict': 'confirmed', 'claim_id': 'B-003', 'model_verdict': 'SUPPORTED', 'evidence': 'The minimum TLS version is 1.2.', 'rationale': 'The reference text explicitly states that the minimum TLS version is 1.2.'}, {'probe_id': 'm-listen-port', 'expected_verdict': 'SUPPORTED', 'anchor': {'all_of': ['listens on port 8443']}, 'verdict': 'confirmed', 'claim_id': 'M-001', 'model_verdict': 'SUPPORTED', 'evidence': 'The relay listens on port 8443.', 'rationale': 'The claim is explicitly stated in source_a.md.'}, {'probe_id': 'm-read-timeout', 'expected_verdict': 'SUPPORTED', 'anchor': {'all_of': ['read timeout is 30 seconds']}, 'verdict': 'confirmed', 'claim_id': 'M-002', 'model_verdict': 'SUPPORTED', 'evidence': 'The read timeout is 30 seconds.', 'rationale': 'The claim is explicitly stated in source_a.md.'}, {'probe_id': 'm-health-path', 'expected_verdict': 'SUPPORTED', 'anchor': {'all_of': ['/healthz']}, 'verdict': 'confirmed', 'claim_id': 'M-003', 'model_verdict': 'SUPPORTED', 'evidence': 'The health check path is /healthz.', 'rationale': 'The claim is explicitly stated in source_b.md.'}, {'probe_id': 'm-tls-min', 'expected_verdict': 'SUPPORTED', 'anchor': {'all_of': ['minimum TLS version is 1.2']}, 'verdict': 'confirmed', 'claim_id': 'M-004', 'model_verdict': 'SUPPORTED', 'evidence': 'The minimum TLS version is 1.2.', 'rationale': 'The claim is explicitly stated in source_b.md.'}]}]}
\expandafter\def\csname ccval@regrade.guard-probes.B-70b.fixtures\endcsname{[{'fixture': 'attribution_invented', 'guard_probes_total': 6, 'guard_confirmed': ['a-listen-port', 'a-read-timeout', 'a-health-path', 'b-tls-min', 'b-max-connections', 'm-read-timeout'], 'guard_wrong': [], 'guard_unclaimed': [], 'detail': [{'probe_id': 'a-listen-port', 'expected_verdict': 'SUPPORTED', 'anchor': {'all_of': ['listens on port 8443']}, 'verdict': 'confirmed', 'claim_id': 'A-001', 'model_verdict': 'SUPPORTED', 'evidence': 'The relay listens on port 8443.', 'rationale': 'The reference text explicitly states that the relay listens on port 8443, matching the claim exactly. The wording is identical, so the claim is supported. The evidence is the exact sentence from the reference text, preserving the wording and any potential errors. The rationale is clear and concise, fitting within the character limit without being cut off mid-sentence. The verdict is correctly set,'}, {'probe_id': 'a-read-timeout', 'expected_verdict': 'SUPPORTED', 'anchor': {'all_of': ['read timeout is 30 seconds']}, 'verdict': 'confirmed', 'claim_id': 'A-002', 'model_verdict': 'SUPPORTED', 'evidence': 'The read timeout is 30 seconds.', 'rationale': 'The reference text explicitly states that the read timeout is 30 seconds, matching the claim exactly. The wording is identical, so the claim is supported. The evidence is the exact sentence from the reference text, preserving the wording and any potential errors. The rationale is clear and concise, fitting within the character limit without being cut off mid-sentence. The verdict is correctly set,'}, {'probe_id': 'a-health-path', 'expected_verdict': 'SUPPORTED', 'anchor': {'all_of': ['/healthz']}, 'verdict': 'confirmed', 'claim_id': 'A-003', 'model_verdict': 'SUPPORTED', 'evidence': 'The health check path is /healthz.', 'rationale': 'The reference text explicitly states that the health check path is /healthz, matching the claim exactly. The wording is identical, so the claim is supported. The evidence is the exact sentence from the reference text, preserving the wording and any potential errors. The rationale is clear and concise, fitting within the character limit without being cut off mid-sentence. The verdict is correctly'}, {'probe_id': 'b-tls-min', 'expected_verdict': 'SUPPORTED', 'anchor': {'all_of': ['minimum TLS version is 1.2']}, 'verdict': 'confirmed', 'claim_id': 'B-002', 'model_verdict': 'SUPPORTED', 'evidence': 'According to the Operator Guide, the minimum TLS version is 1.2.', 'rationale': 'The reference text explicitly states that the minimum TLS version is 1.2, according to the Operator Guide, matching the claim exactly. The wording is identical, so the claim is supported. The evidence is the exact sentence from the reference text, preserving the wording and any potential errors. The rationale is clear and concise, fitting within the character limit without being cut off mid-sent'}, {'probe_id': 'b-max-connections', 'expected_verdict': 'SUPPORTED', 'anchor': {'all_of': ['512']}, 'verdict': 'confirmed', 'claim_id': 'B-003', 'model_verdict': 'SUPPORTED', 'evidence': 'The maximum concurrent connections is 512.', 'rationale': 'The reference text explicitly states that the maximum concurrent connections is 512, matching the claim exactly. The wording is identical, so the claim is supported. The evidence is the exact sentence from the reference text, preserving the wording and any potential errors. The rationale is clear and concise, fitting within the character limit without being cut off mid-sentence. The verdict is cor'}, {'probe_id': 'm-read-timeout', 'expected_verdict': 'SUPPORTED', 'anchor': {'all_of': ['read timeout is 30 seconds']}, 'verdict': 'confirmed', 'claim_id': 'M-002', 'model_verdict': 'SUPPORTED', 'evidence': 'The read timeout is 30 seconds.', 'rationale': 'The claim is directly stated in source_a.md with the exact wording and meaning. Source_b.md does not mention the read timeout, so the support comes solely from source_a.md. The evidence is taken from source_a.md as it matches the claim exactly. The rationale is that the claim is directly stated in source_a.md, making it fully supported without any contradictions or missing parts. A'}]}, {'fixture': 'attribution_swapped', 'guard_probes_total': 2, 'guard_confirmed': ['a-read-timeout', 'b-read-timeout'], 'guard_wrong': [], 'guard_unclaimed': [], 'detail': [{'probe_id': 'a-read-timeout', 'expected_verdict': 'SUPPORTED', 'anchor': {'all_of': ['read timeout is 30 seconds']}, 'verdict': 'confirmed', 'claim_id': 'A-002', 'model_verdict': 'SUPPORTED', 'evidence': 'According to the Deployment Notes, the read timeout is 30 seconds.', 'rationale': "The reference text states that the Deployment Notes specify a read timeout of 30 seconds, which matches the claim. The evidence is the exact sentence from the reference text, including the attribution to the Deployment Notes. The rationale explains that the claim is supported by the specific source mentioned in the reference text. The verdict is 'SUPPORTED'."}, {'probe_id': 'b-read-timeout', 'expected_verdict': 'SUPPORTED', 'anchor': {'all_of': ['read timeout is 60 seconds']}, 'verdict': 'confirmed', 'claim_id': 'B-002', 'model_verdict': 'SUPPORTED', 'evidence': 'According to the Operator Guide, the read timeout is 60 seconds.', 'rationale': "The reference text states that the Operator Guide specifies a read timeout of 60 seconds, which matches the claim. The evidence is the exact sentence from the reference text, including the attribution to the Operator Guide. The rationale explains that the claim is supported by the specific source mentioned in the reference text. The verdict is 'SUPPORTED'."}]}, {'fixture': 'conflict_surfaced', 'guard_probes_total': 11, 'guard_confirmed': ['a-connect-timeout', 'b-connect-timeout', 'a-listen-port', 'a-read-timeout', 'a-max-connections', 'a-tls-version', 'b-log-format', 'b-health-path', 'merged-attributed-timeout-a', 'merged-attributed-timeout-b', 'merged-max-connections'], 'guard_wrong': [], 'guard_unclaimed': [], 'detail': [{'probe_id': 'a-connect-timeout', 'expected_verdict': 'SUPPORTED', 'anchor': {'all_of': ['connect timeout', '30 seconds']}, 'verdict': 'confirmed', 'claim_id': 'A-002', 'model_verdict': 'SUPPORTED', 'evidence': 'According to the Operator Guide, the default connect timeout is 30 seconds.', 'rationale': "The reference text states that the Operator Guide specifies the default connect timeout as 30 seconds, which directly supports the claim. The evidence is the exact span mentioning the 30-second timeout, and the rationale explains that the claim is supported by the Operator Guide's specification."}, {'probe_id': 'b-connect-timeout', 'expected_verdict': 'SUPPORTED', 'anchor': {'all_of': ['connect timeout', '60 seconds']}, 'verdict': 'confirmed', 'claim_id': 'B-002', 'model_verdict': 'SUPPORTED', 'evidence': 'According to the Deployment Notes, the default connect timeout is 60 seconds.', 'rationale': "The reference text states that the Deployment Notes specify the default connect timeout as 60 seconds, which directly supports the claim. The evidence is the exact span mentioning the 60-second timeout, and the rationale explains that the claim is supported by the Deployment Notes' specification."}, {'probe_id': 'a-listen-port', 'expected_verdict': 'SUPPORTED', 'anchor': {'all_of': ['listens on port', '8443']}, 'verdict': 'confirmed', 'claim_id': 'A-001', 'model_verdict': 'SUPPORTED', 'evidence': 'The relay listens on port 8443 by default.', 'rationale': 'The reference text explicitly states that the relay listens on port 8443 by default, matching the claim exactly. The wording is identical, so the claim is supported. The evidence is the exact sentence from the reference text, preserving the wording and any potential errors. The rationale is concise and fits within the character limit, clearly explaining why the verdict is SUPPORTED without being a'}, {'probe_id': 'a-read-timeout', 'expected_verdict': 'SUPPORTED', 'anchor': {'all_of': ['read timeout', '120 seconds']}, 'verdict': 'confirmed', 'claim_id': 'A-003', 'model_verdict': 'SUPPORTED', 'evidence': 'The default read timeout is 120 seconds.', 'rationale': 'The reference text explicitly states that the default read timeout is 120 seconds, matching the claim exactly. The evidence is the exact sentence from the reference text, and the rationale is concise, explaining that the claim is supported by the direct statement in the reference text.'}, {'probe_id': 'a-max-connections', 'expected_verdict': 'SUPPORTED', 'anchor': {'all_of': ['512 concurrent connections']}, 'verdict': 'confirmed', 'claim_id': 'A-004', 'model_verdict': 'SUPPORTED', 'evidence': 'The relay accepts at most 512 concurrent connections.', 'rationale': 'The reference text explicitly states that the relay accepts at most 512 concurrent connections, matching the claim exactly. The evidence is the exact sentence from the reference text, and the rationale is concise, explaining that the claim is supported by the direct statement in the reference text.'}, {'probe_id': 'a-tls-version', 'expected_verdict': 'SUPPORTED', 'anchor': {'all_of': ['minimum accepted TLS version', 'TLS 1.2']}, 'verdict': 'confirmed', 'claim_id': 'A-006', 'model_verdict': 'SUPPORTED', 'evidence': 'The minimum accepted TLS version is TLS 1.2.', 'rationale': 'The reference text explicitly states that the minimum accepted TLS version is TLS 1.2, matching the claim exactly. The evidence is the exact sentence from the reference text, and the rationale is concise, explaining that the claim is supported by the direct statement in the reference text.'}, {'probe_id': 'b-log-format', 'expected_verdict': 'SUPPORTED', 'anchor': {'all_of': ['access logs', 'JSON Lines']}, 'verdict': 'confirmed', 'claim_id': 'B-004', 'model_verdict': 'SUPPORTED', 'evidence': 'Access logs are written in JSON Lines format.', 'rationale': 'The reference text explicitly states that access logs are written in JSON Lines format, matching the claim exactly. The evidence is the exact sentence from the reference text, and the rationale is concise, explaining that the claim is supported by the direct statement in the reference text.'}, {'probe_id': 'b-health-path', 'expected_verdict': 'SUPPORTED', 'anchor': {'all_of': ['health check endpoint', '/-/healthy']}, 'verdict': 'confirmed', 'claim_id': 'B-005', 'model_verdict': 'SUPPORTED', 'evidence': 'The health check endpoint is /-/healthy.', 'rationale': 'The reference text explicitly states that the health check endpoint is /-/healthy, matching the claim exactly. The evidence is the exact sentence from the reference text, and the rationale is concise, explaining that the claim is supported by the direct statement in the reference text.'}, {'probe_id': 'merged-attributed-timeout-a', 'expected_verdict': 'SUPPORTED', 'anchor': {'all_of': ['According to the Operator Guide', '30 seconds']}, 'verdict': 'confirmed', 'claim_id': 'M-003', 'model_verdict': 'SUPPORTED', 'evidence': 'The default connect timeout is 30 seconds.', 'rationale': 'Source_a.md, the Operator Guide, states the default connect timeout is 30 seconds, directly supporting the claim as it is specifically attributed to the correct source document. The evidence is accurately quoted, and the source is correctly identified, ensuring the claim is supported without any discrepancies or additional information that would contradict it. The'}, {'probe_id': 'merged-attributed-timeout-b', 'expected_verdict': 'SUPPORTED', 'anchor': {'all_of': ['According to the Deployment Notes', '60 seconds']}, 'verdict': 'confirmed', 'claim_id': 'M-004', 'model_verdict': 'SUPPORTED', 'evidence': 'The default connect timeout is 60 seconds.', 'rationale': 'Source_b.md, the Deployment Notes, states the default connect timeout is 60 seconds, directly supporting the claim as it is specifically attributed to the correct source document. The evidence is accurately quoted, and the source is correctly identified, ensuring the claim is supported without any discrepancies or additional information that would contradict it. The'}, {'probe_id': 'merged-max-connections', 'expected_verdict': 'SUPPORTED', 'anchor': {'all_of': ['512 concurrent connections']}, 'verdict': 'confirmed', 'claim_id': 'M-005', 'model_verdict': 'SUPPORTED', 'evidence': 'The relay accepts at most 512 concurrent connections.', 'rationale': 'Source_a.md states the relay accepts at most 512 concurrent connections, providing direct support for the claim with the exact wording preserved, ensuring the evidence is accurate and the source is correctly identified. There are no contradictions or missing parts, so the claim is fully supported as it appears in the source document. The'}]}, {'fixture': 'contradiction', 'guard_probes_total': 7, 'guard_confirmed': ['b-connect-timeout', 'a-listen-port', 'a-read-timeout', 'a-max-connections', 'a-tls-version', 'b-log-format', 'b-health-path'], 'guard_wrong': [], 'guard_unclaimed': [], 'detail': [{'probe_id': 'b-connect-timeout', 'expected_verdict': 'SUPPORTED', 'anchor': {'all_of': ['connect timeout', '60 seconds']}, 'verdict': 'confirmed', 'claim_id': 'B-002', 'model_verdict': 'SUPPORTED', 'evidence': 'The default connect timeout is 60 seconds.', 'rationale': 'The reference text explicitly states that the default connect timeout is 60 seconds, matching the claim exactly. SUPPORTED is correct because the wording matches exactly with the reference text. The evidence is the exact sentence from the reference text. The rationale explains that the claim is directly supported by the reference text. The verdict is correct because the claim is fully supported by'}, {'probe_id': 'a-listen-port', 'expected_verdict': 'SUPPORTED', 'anchor': {'all_of': ['listens on port', '8443']}, 'verdict': 'confirmed', 'claim_id': 'A-001', 'model_verdict': 'SUPPORTED', 'evidence': 'The relay listens on port 8443 by default.', 'rationale': 'The reference text explicitly states that the relay listens on port 8443 by default, matching the claim exactly. SUPPORTED is correct because the wording matches exactly with the reference text. The evidence is the exact sentence from the reference text. The rationale explains that the claim is directly supported by the reference text. The verdict is correct because the claim is fully supported by'}, {'probe_id': 'a-read-timeout', 'expected_verdict': 'SUPPORTED', 'anchor': {'all_of': ['read timeout', '120 seconds']}, 'verdict': 'confirmed', 'claim_id': 'A-003', 'model_verdict': 'SUPPORTED', 'evidence': 'The default read timeout is 120 seconds.', 'rationale': 'The reference text explicitly states that the default read timeout is 120 seconds, matching the claim exactly. SUPPORTED is correct because the wording matches exactly with the reference text. The evidence is the exact sentence from the reference text. The rationale explains that the claim is directly supported by the reference text. The verdict is correct because the claim is fully supported by'}, {'probe_id': 'a-max-connections', 'expected_verdict': 'SUPPORTED', 'anchor': {'all_of': ['512 concurrent connections']}, 'verdict': 'confirmed', 'claim_id': 'A-004', 'model_verdict': 'SUPPORTED', 'evidence': 'The relay accepts at most 512 concurrent connections.', 'rationale': 'The reference text explicitly states that the relay accepts at most 512 concurrent connections, matching the claim exactly. SUPPORTED is correct because the wording matches exactly with the reference text. The evidence is the exact sentence from the reference text. The rationale explains that the claim is directly supported by the reference text. The verdict is correct because the claim is fully'}, {'probe_id': 'a-tls-version', 'expected_verdict': 'SUPPORTED', 'anchor': {'all_of': ['minimum accepted TLS version', 'TLS 1.2']}, 'verdict': 'confirmed', 'claim_id': 'A-006', 'model_verdict': 'SUPPORTED', 'evidence': 'The minimum accepted TLS version is TLS 1.2.', 'rationale': 'The reference text explicitly states that the minimum accepted TLS version is TLS 1.2, matching the claim exactly. SUPPORTED is correct because the wording matches exactly with the reference text. The evidence is the exact sentence from the reference text. The rationale explains that the claim is directly supported by the reference text. The verdict is correct because the claim is fully supported'}, {'probe_id': 'b-log-format', 'expected_verdict': 'SUPPORTED', 'anchor': {'all_of': ['access logs', 'JSON Lines']}, 'verdict': 'confirmed', 'claim_id': 'B-004', 'model_verdict': 'SUPPORTED', 'evidence': 'Access logs are written in JSON Lines format.', 'rationale': 'The reference text explicitly states that access logs are written in JSON Lines format, matching the claim exactly. SUPPORTED is correct because the wording matches exactly with the reference text. The evidence is the exact sentence from the reference text. The rationale explains that the claim is directly supported by the reference text. The verdict is correct because the claim is fully supported'}, {'probe_id': 'b-health-path', 'expected_verdict': 'SUPPORTED', 'anchor': {'all_of': ['health check endpoint', '/-/healthy']}, 'verdict': 'confirmed', 'claim_id': 'B-005', 'model_verdict': 'SUPPORTED', 'evidence': 'The health check endpoint is /-/healthy.', 'rationale': 'The reference text explicitly states that the health check endpoint is /-/healthy, matching the claim exactly. SUPPORTED is correct because the wording matches exactly with the reference text. The evidence is the exact sentence from the reference text. The rationale explains that the claim is directly supported by the reference text. The verdict is correct because the claim is fully supported by'}]}, {'fixture': 'dedup', 'guard_probes_total': 10, 'guard_confirmed': ['a-listen-port', 'a-health-path', 'a-read-timeout', 'b-listen-port', 'b-health-path', 'b-tls-min', 'm-listen-port', 'm-health-path', 'm-read-timeout', 'm-tls-min'], 'guard_wrong': [], 'guard_unclaimed': [], 'detail': [{'probe_id': 'a-listen-port', 'expected_verdict': 'SUPPORTED', 'anchor': {'all_of': ['listens on port 8443']}, 'verdict': 'confirmed', 'claim_id': 'A-001', 'model_verdict': 'SUPPORTED', 'evidence': 'The relay listens on port 8443.', 'rationale': "The reference text states that the relay listens on port 8443, matching the claim exactly. Verdict: SUPPORTED, evidence: 'The relay listens on port 8443.', source: merged.md, rationale: The claim is directly stated in the reference text. The evidence is the exact sentence from the reference text that matches the claim. The rationale explains that the claim is supported because it is directly found"}, {'probe_id': 'a-health-path', 'expected_verdict': 'SUPPORTED', 'anchor': {'all_of': ['/healthz']}, 'verdict': 'confirmed', 'claim_id': 'A-002', 'model_verdict': 'SUPPORTED', 'evidence': 'The health check path is /healthz.', 'rationale': "The reference text states that the health check path is /healthz, matching the claim exactly. Verdict: SUPPORTED, evidence: 'The health check path is /healthz.', source: merged.md, rationale: The claim is directly stated in the reference text. The evidence is the exact sentence from the reference text that matches the claim. The rationale explains that the claim is supported because it is directly"}, {'probe_id': 'a-read-timeout', 'expected_verdict': 'SUPPORTED', 'anchor': {'all_of': ['read timeout is 30 seconds']}, 'verdict': 'confirmed', 'claim_id': 'A-003', 'model_verdict': 'SUPPORTED', 'evidence': 'The read timeout is 30 seconds.', 'rationale': "The reference text states that the read timeout is 30 seconds, matching the claim exactly. Verdict: SUPPORTED, evidence: 'The read timeout is 30 seconds.', source: merged.md, rationale: The claim is directly stated in the reference text. The evidence is the exact sentence from the reference text that matches the claim. The rationale explains that the claim is supported because it is directly"}, {'probe_id': 'b-listen-port', 'expected_verdict': 'SUPPORTED', 'anchor': {'all_of': ['listens on port 8443']}, 'verdict': 'confirmed', 'claim_id': 'A-001', 'model_verdict': 'SUPPORTED', 'evidence': 'The relay listens on port 8443.', 'rationale': "The reference text states that the relay listens on port 8443, matching the claim exactly. Verdict: SUPPORTED, evidence: 'The relay listens on port 8443.', source: merged.md, rationale: The claim is directly stated in the reference text. The evidence is the exact sentence from the reference text that matches the claim. The rationale explains that the claim is supported because it is directly found"}, {'probe_id': 'b-health-path', 'expected_verdict': 'SUPPORTED', 'anchor': {'all_of': ['/healthz']}, 'verdict': 'confirmed', 'claim_id': 'A-002', 'model_verdict': 'SUPPORTED', 'evidence': 'The health check path is /healthz.', 'rationale': "The reference text states that the health check path is /healthz, matching the claim exactly. Verdict: SUPPORTED, evidence: 'The health check path is /healthz.', source: merged.md, rationale: The claim is directly stated in the reference text. The evidence is the exact sentence from the reference text that matches the claim. The rationale explains that the claim is supported because it is directly"}, {'probe_id': 'b-tls-min', 'expected_verdict': 'SUPPORTED', 'anchor': {'all_of': ['minimum TLS version is 1.2']}, 'verdict': 'confirmed', 'claim_id': 'B-003', 'model_verdict': 'SUPPORTED', 'evidence': 'The minimum TLS version is 1.2.', 'rationale': "The reference text states that the minimum TLS version is 1.2, matching the claim exactly. Verdict: SUPPORTED, evidence: 'The minimum TLS version is 1.2.', source: merged.md, rationale: The claim is directly stated in the reference text. The evidence is the exact sentence from the reference text that matches the claim. The rationale explains that the claim is supported because it is directly"}, {'probe_id': 'm-listen-port', 'expected_verdict': 'SUPPORTED', 'anchor': {'all_of': ['listens on port 8443']}, 'verdict': 'confirmed', 'claim_id': 'M-001', 'model_verdict': 'SUPPORTED', 'evidence': 'The relay listens on port 8443.', 'rationale': 'Both source_a.md and source_b.md state that the relay listens on port 8443, so the claim is supported by at least one source document. The exact wording matches source_a.md, so the evidence is taken from there. The claim is SUPPORTED because it is stated in both source documents, but evidence is only required from one. The rationale is that the claim is present in both sources, so it is supported.'}, {'probe_id': 'm-health-path', 'expected_verdict': 'SUPPORTED', 'anchor': {'all_of': ['/healthz']}, 'verdict': 'confirmed', 'claim_id': 'M-002', 'model_verdict': 'SUPPORTED', 'evidence': 'The health check path is /healthz.', 'rationale': 'Both source_a.md and source_b.md state that the health check path is /healthz, so the claim is supported by at least one source document. The exact wording matches source_a.md, so the evidence is taken from there. The claim is SUPPORTED because it is stated in both source documents, but evidence is only required from one. The rationale is that the claim is present in both sources, so it is fully'}, {'probe_id': 'm-read-timeout', 'expected_verdict': 'SUPPORTED', 'anchor': {'all_of': ['read timeout is 30 seconds']}, 'verdict': 'confirmed', 'claim_id': 'M-003', 'model_verdict': 'SUPPORTED', 'evidence': 'The read timeout is 30 seconds.', 'rationale': 'Source_a.md states that the read timeout is 30 seconds, so the claim is supported by at least one source document. Source_b.md does not mention the read timeout, but the claim is still supported because it appears in source_a.md. The evidence is taken from source_a.md because it explicitly states the claim. The rationale is that the claim is present in source_a.md, so it is supported.'}, {'probe_id': 'm-tls-min', 'expected_verdict': 'SUPPORTED', 'anchor': {'all_of': ['minimum TLS version is 1.2']}, 'verdict': 'confirmed', 'claim_id': 'M-004', 'model_verdict': 'SUPPORTED', 'evidence': 'The minimum TLS version is 1.2.', 'rationale': 'Source_b.md states that the minimum TLS version is 1.2, so the claim is supported by at least one source document. Source_a.md does not mention the TLS version, but the claim is still supported because it appears in source_b.md. The evidence is taken from source_b.md because it explicitly states the claim. The rationale is that the claim is present in source_b.md, so it is supported.'}]}, {'fixture': 'disjoint_domains', 'guard_probes_total': 16, 'guard_confirmed': ['a-round-hours', 'a-second-round-verbatim', 'a-frost-latch', 'a-responsibility', 'a-contractor-dates', 'a-tie-officer', 'b-closed-summer', 'b-boards-left-up', 'b-ladder-not-rating', 'b-ladder-moves', 'b-spare-clocks', 'b-extra-table', 'm-round-hours-reverse', 'm-club-season-reverse', 'm-contractor-dates-reverse', 'm-ladder-reverse'], 'guard_wrong': [], 'guard_unclaimed': [], 'detail': [{'probe_id': 'a-round-hours', 'expected_verdict': 'SUPPORTED', 'anchor': {'all_of': ['7 am', '8 am']}, 'verdict': 'confirmed', 'claim_id': 'A-002', 'model_verdict': 'SUPPORTED', 'evidence': 'The duty officer walks the north corridor every weekday from 7 am till 8 am, right after the shift handover.', 'rationale': 'The reference text explicitly states that the duty officer walks the north corridor every weekday from 7 am till 8 am, right after the shift handover, which directly supports the claim A-002. The details provided in the text match the claim precisely, confirming the support without any ambiguity or contradiction. The evidence is clear and directly addresses the claim, leaving no room for doubt or'}, {'probe_id': 'a-second-round-verbatim', 'expected_verdict': 'SUPPORTED', 'anchor': {'all_of': ['13:05', '16:20']}, 'verdict': 'confirmed', 'claim_id': 'A-003', 'model_verdict': 'SUPPORTED', 'evidence': 'A second round covers the workshop and the paint store from 13:05 till 16:20 hours on Wednesdays only.', 'rationale': 'The reference text explicitly states that a second round covers the workshop and the paint store from 13:05 till 16:20 hours on Wednesdays only, which directly supports the claim A-003. The details provided in the text match the claim precisely, confirming the support without any ambiguity or contradiction. The evidence is clear and directly addresses the claim, leaving no room for doubt or'}, {'probe_id': 'a-frost-latch', 'expected_verdict': 'SUPPORTED', 'anchor': {'all_of': ['frost', 'November']}, 'verdict': 'confirmed', 'claim_id': 'A-005', 'model_verdict': 'SUPPORTED', 'evidence': 'Frost swells the timber frames on the older doors, so between November and February several of them will not latch unless pulled hard.', 'rationale': 'The reference text explicitly states that frost swells the timber frames on the older doors, so between November and February several of them will not latch unless pulled hard, which directly supports the claim A-005. The details provided in the text match the claim precisely, confirming the support without any ambiguity or contradiction. The evidence is clear and directly addresses the claim,'}, {'probe_id': 'a-responsibility', 'expected_verdict': 'SUPPORTED', 'anchor': {'all_of': ['responsibility of the site']}, 'verdict': 'confirmed', 'claim_id': 'A-007', 'model_verdict': 'SUPPORTED', 'evidence': 'While the doors are the responsibility of the site, the certified inspection is carried out by an outside contractor who attends every March 1st to March 31st.', 'rationale': 'The reference text explicitly states that while the doors are the responsibility of the site, the certified inspection is carried out by an outside contractor who attends every March 1st to March 31st, which directly supports the claim A-007. The details provided in the text match the claim precisely, confirming the support without any ambiguity or contradiction. The evidence is clear and directly'}, {'probe_id': 'a-contractor-dates', 'expected_verdict': 'SUPPORTED', 'anchor': {'all_of': ['March 1st', 'March 31st']}, 'verdict': 'confirmed', 'claim_id': 'A-007', 'model_verdict': 'SUPPORTED', 'evidence': 'While the doors are the responsibility of the site, the certified inspection is carried out by an outside contractor who attends every March 1st to March 31st.', 'rationale': 'The reference text explicitly states that while the doors are the responsibility of the site, the certified inspection is carried out by an outside contractor who attends every March 1st to March 31st, which directly supports the claim A-007. The details provided in the text match the claim precisely, confirming the support without any ambiguity or contradiction. The evidence is clear and directly'}, {'probe_id': 'a-tie-officer', 'expected_verdict': 'SUPPORTED', 'anchor': {'all_of': ['one officer clears both']}, 'verdict': 'confirmed', 'claim_id': 'A-012', 'model_verdict': 'SUPPORTED', 'evidence': 'In most cases it is a tie, because one officer clears both at the end of the week.', 'rationale': 'The reference text explicitly states that in most cases it is a tie because one officer clears both at the end of the week, which directly supports the claim A-012. The details provided in the text match the claim precisely, confirming the support without any ambiguity or contradiction. The evidence is clear and directly addresses the claim, leaving no room for doubt or partial support.'}, {'probe_id': 'b-closed-summer', 'expected_verdict': 'SUPPORTED', 'anchor': {'all_of': ['does not meet', 'July', 'September']}, 'verdict': 'confirmed', 'claim_id': 'B-002', 'model_verdict': 'SUPPORTED', 'evidence': 'The club does not meet between mid July and early September, because most members take leave and the canteen is closed for cleaning.', 'rationale': 'The reference text explicitly states that the club does not meet between mid-July and early September because most members take leave and the canteen is closed for cleaning, which directly supports the claim B-002. The details provided in the text match the claim precisely, confirming the support without any ambiguity or contradiction. The evidence is clear and directly addresses the claim,'}, {'probe_id': 'b-boards-left-up', 'expected_verdict': 'SUPPORTED', 'anchor': {'all_of': ['boards set up']}, 'verdict': 'confirmed', 'claim_id': 'B-004', 'model_verdict': 'SUPPORTED', 'evidence': 'Over the years the members have learned to leave the boards set up rather than pack them away.', 'rationale': 'The reference text explicitly states that over the years the members have learned to leave the boards set up rather than pack them away, which directly supports the claim B-004. The details provided in the text match the claim precisely, confirming the support without any ambiguity or contradiction. The evidence is clear and directly addresses the claim, leaving no room for doubt or partial'}, {'probe_id': 'b-ladder-not-rating', 'expected_verdict': 'SUPPORTED', 'anchor': {'all_of': ['simple ladder']}, 'verdict': 'confirmed', 'claim_id': 'B-005', 'model_verdict': 'SUPPORTED', 'evidence': 'Instead of a rating system, the club runs a simple ladder that anyone may join at any point.', 'rationale': 'The reference text explicitly states that instead of a rating system, the club runs a simple ladder that anyone may join at any point, which directly supports the claim B-005. The details provided in the text match the claim precisely, confirming the support without any ambiguity or contradiction. The evidence is clear and directly addresses the claim, leaving no room for doubt or partial'}, {'probe_id': 'b-ladder-moves', 'expected_verdict': 'SUPPORTED', 'anchor': {'all_of': ['up one place', 'down one place']}, 'verdict': 'confirmed', 'claim_id': 'B-006', 'model_verdict': 'SUPPORTED', 'evidence': 'Members who win a game move up one place, and the person they beat moves down one place.', 'rationale': 'The reference text explicitly states that members who win a game move up one place, and the person they beat moves down one place, which directly supports the claim B-006. The details provided in the text match the claim precisely, confirming the support without any ambiguity or contradiction. The evidence is clear and directly addresses the claim, leaving no room for doubt or partial'}, {'probe_id': 'b-spare-clocks', 'expected_verdict': 'SUPPORTED', 'anchor': {'all_of': ['spare clocks']}, 'verdict': 'confirmed', 'claim_id': 'B-014', 'model_verdict': 'SUPPORTED', 'evidence': 'When the match was over, the county side gave the club enough spare clocks to run every board on the ladder at the same time.', 'rationale': 'The reference text directly states that the county side provided clocks to the club after the match ended, matching the claim exactly. The wording is identical, so the claim is supported without any contradiction or partial support issues. The evidence is taken verbatim from the text, preserving any original formatting or spelling, as required by the instructions. The rationale is concise and fits'}, {'probe_id': 'b-extra-table', 'expected_verdict': 'SUPPORTED', 'anchor': {'all_of': ['extra table', 'Tuesday']}, 'verdict': 'confirmed', 'claim_id': 'B-016', 'model_verdict': 'SUPPORTED', 'evidence': 'Eventually the club grew to thirty members, and from then on the canteen staff put out an extra table every Tuesday so that nobody had to stand.', 'rationale': "The reference text clearly states that the club's membership reached thirty, leading to the addition of an extra table each Tuesday. This directly aligns with the claim, providing full support without any contradictions or partial elements. The evidence is quoted exactly as it appears, adhering to the specified rules. The rationale is succinct and fits within the required"}, {'probe_id': 'm-round-hours-reverse', 'expected_verdict': 'SUPPORTED', 'anchor': {'all_of': ['7 am', '8 am']}, 'verdict': 'confirmed', 'claim_id': 'M-002', 'model_verdict': 'SUPPORTED', 'evidence': 'The duty officer walks the north corridor every weekday from 7 am till 8 am, right after the shift handover.', 'rationale': "The claim is exactly stated in source_a.md's Rounds section."}, {'probe_id': 'm-club-season-reverse', 'expected_verdict': 'SUPPORTED', 'anchor': {'all_of': ['does not meet', 'July', 'September']}, 'verdict': 'confirmed', 'claim_id': 'M-014', 'model_verdict': 'SUPPORTED', 'evidence': 'The club does not meet between mid July and early September, because most members take leave and the canteen is closed for cleaning.', 'rationale': "The claim is exactly stated in source_b.md's Season section."}, {'probe_id': 'm-contractor-dates-reverse', 'expected_verdict': 'SUPPORTED', 'anchor': {'all_of': ['March 1st', 'March 31st']}, 'verdict': 'confirmed', 'claim_id': 'M-007', 'model_verdict': 'SUPPORTED', 'evidence': 'While the doors are the responsibility of the site, the certified inspection is carried out by an outside contractor who attends every March 1st to March 31st.', 'rationale': "The claim is exactly stated in source_a.md's Contractor visits section."}, {'probe_id': 'm-ladder-reverse', 'expected_verdict': 'SUPPORTED', 'anchor': {'all_of': ['simple ladder']}, 'verdict': 'confirmed', 'claim_id': 'M-018', 'model_verdict': 'SUPPORTED', 'evidence': 'Instead of a rating system, the club runs a simple ladder that anyone may join at any point.', 'rationale': "The claim is exactly stated in source_b.md's Ladder section."}]}, {'fixture': 'disjoint_sources', 'guard_probes_total': 16, 'guard_confirmed': ['a-visit-hours', 'a-work-hours-verbatim', 'a-winter-dormant', 'a-ownership', 'a-public-dates', 'a-postcard-tie', 'b-not-traversable', 'b-lady-year-round', 'b-herbs-not-magic', 'b-herbs-to-cities', 'b-medicine-village', 'b-cheese-autumn', 'm-visit-hours-reverse', 'm-forest-season-reverse', 'm-public-dates-reverse', 'm-herbs-reverse'], 'guard_wrong': [], 'guard_unclaimed': [], 'detail': [{'probe_id': 'a-visit-hours', 'expected_verdict': 'SUPPORTED', 'anchor': {'all_of': ['8 am', '9 am']}, 'verdict': 'confirmed', 'claim_id': 'A-002', 'model_verdict': 'SUPPORTED', 'evidence': 'She visited the rosegarden every morning from 8 am till 9 am, right after her breakfast.', 'rationale': 'The reference text directly states the claim about the princess visiting the rosegarden every morning from 8 am till 9 am after breakfast, matching the wording exactly. Therefore, the claim is SUPPORTED with the provided evidence from the text. The evidence is a direct quote from the reference text, confirming the claim without any discrepancies or contradictions. The rationale is clear and'}, {'probe_id': 'a-work-hours-verbatim', 'expected_verdict': 'SUPPORTED', 'anchor': {'all_of': ['02:00', '15:30']}, 'verdict': 'confirmed', 'claim_id': 'A-003', 'model_verdict': 'SUPPORTED', 'evidence': "From 02:00 till 15:30 o'clock, she worked in the rosegarden to tidy it up.", 'rationale': "The reference text directly states the claim about the princess working in the rosegarden from 02:00 till 15:30 o'clock, matching the wording exactly. Therefore, the claim is SUPPORTED with the provided evidence from the text. The evidence is a direct quote from the reference text, confirming the claim without any discrepancies or contradictions. The rationale is clear and concise, explaining why"}, {'probe_id': 'a-winter-dormant', 'expected_verdict': 'SUPPORTED', 'anchor': {'all_of': ['winter', 'rosegarden']}, 'verdict': 'confirmed', 'claim_id': 'A-005', 'model_verdict': 'SUPPORTED', 'evidence': 'The rosegarden lost all his life in winter, so the young lady could not enjoy the beautiful colors and all the life the rosegarden harbored.', 'rationale': 'The reference text directly states the claim about the rosegarden losing life in winter, matching the wording exactly. Therefore, the claim is SUPPORTED with the provided evidence from the text. The evidence is a direct quote from the reference text, confirming the claim without any discrepancies or contradictions. The rationale is clear and concise, explaining why the verdict is'}, {'probe_id': 'a-ownership', 'expected_verdict': 'SUPPORTED', 'anchor': {'all_of': ['private property', 'royal family']}, 'verdict': 'confirmed', 'claim_id': 'A-007', 'model_verdict': 'SUPPORTED', 'evidence': 'While the rose garden was the private property of the royal family, the young princess insisted on making the garden accessible to the public every March 1st to March 31st.', 'rationale': 'The reference text directly states the claim about the princess making the garden public from March 1st to 31st, matching the wording exactly. Therefore, the claim is SUPPORTED with the provided evidence from the text. The evidence is a direct quote from the reference text, confirming the claim without any discrepancies or contradictions. The rationale is clear and concise, explaining why'}, {'probe_id': 'a-public-dates', 'expected_verdict': 'SUPPORTED', 'anchor': {'all_of': ['March 1st', 'March 31st']}, 'verdict': 'confirmed', 'claim_id': 'A-007', 'model_verdict': 'SUPPORTED', 'evidence': 'While the rose garden was the private property of the royal family, the young princess insisted on making the garden accessible to the public every March 1st to March 31st.', 'rationale': 'The reference text directly states the claim about the princess making the garden public from March 1st to 31st, matching the wording exactly. Therefore, the claim is SUPPORTED with the provided evidence from the text. The evidence is a direct quote from the reference text, confirming the claim without any discrepancies or contradictions. The rationale is clear and concise, explaining why'}, {'probe_id': 'a-postcard-tie', 'expected_verdict': 'SUPPORTED', 'anchor': {'all_of': ['same carriage']}, 'verdict': 'confirmed', 'claim_id': 'A-012', 'model_verdict': 'SUPPORTED', 'evidence': 'In many cases, it was a tie because the people and the postcards were usually transported in the same carriage.', 'rationale': 'The reference text directly states the claim about ties due to same carriage transport, matching the wording exactly. Therefore, the claim is SUPPORTED with the provided evidence from the text. The evidence is a direct quote from the reference text, confirming the claim without any discrepancies or contradictions. The rationale is clear and concise, explaining why the verdict is'}, {'probe_id': 'b-not-traversable', 'expected_verdict': 'SUPPORTED', 'anchor': {'all_of': ['not traversable', 'October', 'April']}, 'verdict': 'confirmed', 'claim_id': 'B-002', 'model_verdict': 'SUPPORTED', 'evidence': 'The forest was not traversable between October and April as it was too wet and too cold.', 'rationale': 'The reference text directly states the claim about the forest being non-traversable from October to April, matching the wording exactly. Therefore, the claim is SUPPORTED with the provided evidence from the text. The evidence is a direct quote from the reference text, confirming the claim without any discrepancies or contradictions. The rationale is clear and concise, explaining why the'}, {'probe_id': 'b-lady-year-round', 'expected_verdict': 'SUPPORTED', 'anchor': {'all_of': ['entire year', 'forest']}, 'verdict': 'confirmed', 'claim_id': 'B-003', 'model_verdict': 'SUPPORTED', 'evidence': 'The old lady nevertheless lived the entire year in the forest.', 'rationale': 'The reference text directly states the claim about the old lady living in the forest all year, matching the wording exactly. Therefore, the claim is SUPPORTED with the provided evidence from the text. The evidence is a direct quote from the reference text, confirming the claim without any discrepancies or contradictions. The rationale is clear and concise, explaining why the verdict'}, {'probe_id': 'b-herbs-not-magic', 'expected_verdict': 'SUPPORTED', 'anchor': {'all_of': ['wild herbs']}, 'verdict': 'confirmed', 'claim_id': 'B-006', 'model_verdict': 'SUPPORTED', 'evidence': 'she collected wild herbs.', 'rationale': 'The reference text directly states the claim about the old lady collecting wild herbs, matching the wording exactly. Therefore, the claim is SUPPORTED with the provided evidence from the text. The evidence is a direct quote from the reference text, confirming the claim without any discrepancies or contradictions. The rationale is clear and concise, explaining why the verdict is'}, {'probe_id': 'b-herbs-to-cities', 'expected_verdict': 'SUPPORTED', 'anchor': {'all_of': ['travelers', 'cities']}, 'verdict': 'confirmed', 'claim_id': 'B-007', 'model_verdict': 'SUPPORTED', 'evidence': 'Those wild herbs were sold to wandering travelers who transported and sold the herbs into the cities.', 'rationale': 'The reference text directly states the claim about selling herbs to travelers, matching the wording exactly. Therefore, the claim is SUPPORTED with the provided evidence from the text. The evidence is a direct quote from the reference text, confirming the claim without any discrepancies or contradictions. The rationale is clear and concise, explaining why the verdict is'}, {'probe_id': 'b-medicine-village', 'expected_verdict': 'SUPPORTED', 'anchor': {'all_of': ['medicine', 'entire village']}, 'verdict': 'confirmed', 'claim_id': 'B-015', 'model_verdict': 'SUPPORTED', 'evidence': 'she gave him enough medicine made out of the herbs she collected in spring for the entire village.', 'rationale': 'The reference text directly states the claim in the same words and meaning, so it is supported by the text provided in merged.md. The evidence is an exact match for the claim, confirming that the claim is accurate based on the reference text. The claim is fully supported without any contradictions or missing information, making the verdict clear and straightforward. The evidence is taken verbatim,'}, {'probe_id': 'b-cheese-autumn', 'expected_verdict': 'SUPPORTED', 'anchor': {'all_of': ['cheese', 'autumn']}, 'verdict': 'confirmed', 'claim_id': 'B-017', 'model_verdict': 'SUPPORTED', 'evidence': 'Eventually, everyone got healthy again and from then on the villagers sent a huge basket of cheese and dried meat to the old lady every autumn so that she could feast in winter.', 'rationale': 'The reference text directly states the claim in the same words and meaning, so it is supported by the text provided in merged.md. The evidence is an exact match for the claim, confirming that the claim is accurate based on the reference text. The claim is fully supported without any contradictions or missing information, making the verdict clear and straightforward. The evidence is taken verbatim,'}, {'probe_id': 'm-visit-hours-reverse', 'expected_verdict': 'SUPPORTED', 'anchor': {'all_of': ['8 am', '9 am']}, 'verdict': 'confirmed', 'claim_id': 'M-002', 'model_verdict': 'SUPPORTED', 'evidence': 'She visited the rosegarden every morning from 8 am till 9 am, right after her breakfast.', 'rationale': 'The claim is directly stated in source_a.md with the exact wording and meaning.'}, {'probe_id': 'm-forest-season-reverse', 'expected_verdict': 'SUPPORTED', 'anchor': {'all_of': ['not traversable', 'October', 'April']}, 'verdict': 'confirmed', 'claim_id': 'M-014', 'model_verdict': 'SUPPORTED', 'evidence': 'The forest was not traversable between October and April as it was too wet and too cold.', 'rationale': 'The claim is directly stated in source_b.md with the exact wording and meaning.'}, {'probe_id': 'm-public-dates-reverse', 'expected_verdict': 'SUPPORTED', 'anchor': {'all_of': ['March 1st', 'March 31st']}, 'verdict': 'confirmed', 'claim_id': 'M-007', 'model_verdict': 'SUPPORTED', 'evidence': 'While the rose garden was the private property of the royal family, the young princess insisted on making the garden accessible to the public every March 1st to March 31st.', 'rationale': 'The claim is directly stated in source_a.md with the exact wording and meaning.'}, {'probe_id': 'm-herbs-reverse', 'expected_verdict': 'SUPPORTED', 'anchor': {'all_of': ['wild herbs']}, 'verdict': 'confirmed', 'claim_id': 'M-018', 'model_verdict': 'SUPPORTED', 'evidence': 'Instead of harvesting skulls and bones, she collected wild herbs.', 'rationale': 'The claim is directly stated in source_b.md with the exact wording and meaning.'}]}, {'fixture': 'dropped_claim', 'guard_probes_total': 9, 'guard_confirmed': ['b-log-rotation', 'b-worker-per-core', 'b-health-path', 'b-health-failures', 'a-listen-port', 'a-connect-timeout', 'a-read-timeout', 'a-retry-limit', 'a-tls-version'], 'guard_wrong': [], 'guard_unclaimed': [], 'detail': [{'probe_id': 'b-log-rotation', 'expected_verdict': 'SUPPORTED', 'anchor': {'all_of': ['rotated', '100 megabytes']}, 'verdict': 'confirmed', 'claim_id': 'B-005', 'model_verdict': 'SUPPORTED', 'evidence': 'Log files are rotated when they reach 100 megabytes.', 'rationale': 'The reference text explicitly states that log files are rotated when they reach 100 megabytes.'}, {'probe_id': 'b-worker-per-core', 'expected_verdict': 'SUPPORTED', 'anchor': {'all_of': ['one worker process', 'per CPU core']}, 'verdict': 'confirmed', 'claim_id': 'B-003', 'model_verdict': 'SUPPORTED', 'evidence': 'One worker process is started per CPU core.', 'rationale': 'The reference text explicitly states that one worker process is started per CPU core.'}, {'probe_id': 'b-health-path', 'expected_verdict': 'SUPPORTED', 'anchor': {'all_of': ['health check endpoint', '/-/healthy']}, 'verdict': 'confirmed', 'claim_id': 'B-006', 'model_verdict': 'SUPPORTED', 'evidence': 'The health check endpoint is /-/healthy.', 'rationale': 'The reference text explicitly states that the health check endpoint is /-/healthy.'}, {'probe_id': 'b-health-failures', 'expected_verdict': 'SUPPORTED', 'anchor': {'all_of': ['3 consecutive health checks', 'removed from rotation']}, 'verdict': 'confirmed', 'claim_id': 'B-007', 'model_verdict': 'SUPPORTED', 'evidence': 'A relay that fails 3 consecutive health checks is removed from rotation.', 'rationale': 'The reference text explicitly states that a relay that fails 3 consecutive health checks is removed from rotation.'}, {'probe_id': 'a-listen-port', 'expected_verdict': 'SUPPORTED', 'anchor': {'all_of': ['listens on port', '8443']}, 'verdict': 'confirmed', 'claim_id': 'A-001', 'model_verdict': 'SUPPORTED', 'evidence': 'The relay listens on port 8443 by default.', 'rationale': 'The reference text explicitly states that the relay listens on port 8443 by default.'}, {'probe_id': 'a-connect-timeout', 'expected_verdict': 'SUPPORTED', 'anchor': {'all_of': ['connect timeout', '30 seconds']}, 'verdict': 'confirmed', 'claim_id': 'A-002', 'model_verdict': 'SUPPORTED', 'evidence': 'The default connect timeout is 30 seconds.', 'rationale': 'The reference text explicitly states that the default connect timeout is 30 seconds.'}, {'probe_id': 'a-read-timeout', 'expected_verdict': 'SUPPORTED', 'anchor': {'all_of': ['read timeout', '120 seconds']}, 'verdict': 'confirmed', 'claim_id': 'A-003', 'model_verdict': 'SUPPORTED', 'evidence': 'The default read timeout is 120 seconds.', 'rationale': 'The reference text explicitly states that the default read timeout is 120 seconds.'}, {'probe_id': 'a-retry-limit', 'expected_verdict': 'SUPPORTED', 'anchor': {'all_of': ['retried at most 3 times']}, 'verdict': 'confirmed', 'claim_id': 'A-005', 'model_verdict': 'SUPPORTED', 'evidence': 'Each request is retried at most 3 times.', 'rationale': 'The reference text explicitly states that each request is retried at most 3 times.'}, {'probe_id': 'a-tls-version', 'expected_verdict': 'SUPPORTED', 'anchor': {'all_of': ['minimum accepted TLS version', 'TLS 1.2']}, 'verdict': 'confirmed', 'claim_id': 'A-006', 'model_verdict': 'SUPPORTED', 'evidence': 'The minimum accepted TLS version is TLS 1.2.', 'rationale': 'The reference text explicitly states that the minimum accepted TLS version is TLS 1.2.'}]}, {'fixture': 'hallucination', 'guard_probes_total': 7, 'guard_confirmed': ['merged-connect-timeout', 'merged-max-connections', 'a-listen-port', 'a-tls-version', 'a-client-cert', 'b-log-format', 'b-health-path'], 'guard_wrong': [], 'guard_unclaimed': [], 'detail': [{'probe_id': 'merged-connect-timeout', 'expected_verdict': 'SUPPORTED', 'anchor': {'all_of': ['connect timeout', '30 seconds']}, 'verdict': 'confirmed', 'claim_id': 'M-002', 'model_verdict': 'SUPPORTED', 'evidence': 'The default connect timeout is 30 seconds.', 'rationale': 'The claim is exactly stated in source_a.md under the Networking section. The evidence is correctly quoted from source_a.md, preserving the exact wording and formatting. The rationale is concise and explains why the claim is supported by the evidence from source_a.md.'}, {'probe_id': 'merged-max-connections', 'expected_verdict': 'SUPPORTED', 'anchor': {'all_of': ['512 concurrent connections']}, 'verdict': 'confirmed', 'claim_id': 'M-003', 'model_verdict': 'SUPPORTED', 'evidence': 'The relay accepts at most 512 concurrent connections.', 'rationale': 'The claim is exactly stated in source_b.md under the Networking section. The evidence is correctly quoted from source_b.md, preserving the exact wording and formatting. The rationale is concise and explains why the claim is supported by the evidence from source_b.md.'}, {'probe_id': 'a-listen-port', 'expected_verdict': 'SUPPORTED', 'anchor': {'all_of': ['listens on port', '8443']}, 'verdict': 'confirmed', 'claim_id': 'A-001', 'model_verdict': 'SUPPORTED', 'evidence': 'The relay listens on port 8443 by default.', 'rationale': 'The reference text explicitly states that the relay listens on port 8443 by default, matching the claim exactly. The wording is identical, so the claim is supported. The evidence is the exact sentence from the reference text, preserving the wording and any potential errors. The rationale is concise and explains why the verdict is SUPPORTED without exceeding the character limit. The evidence_source'}, {'probe_id': 'a-tls-version', 'expected_verdict': 'SUPPORTED', 'anchor': {'all_of': ['minimum accepted TLS version', 'TLS 1.2']}, 'verdict': 'confirmed', 'claim_id': 'A-003', 'model_verdict': 'SUPPORTED', 'evidence': 'The minimum accepted TLS version is TLS 1.2.', 'rationale': 'The reference text explicitly states that the minimum accepted TLS version is TLS 1.2, matching the claim exactly. The wording is identical, so the claim is supported. The evidence is the exact sentence from the reference text, preserving the wording and any potential errors. The rationale is concise and explains why the verdict is SUPPORTED without exceeding the character limit. The evidence_sour'}, {'probe_id': 'a-client-cert', 'expected_verdict': 'SUPPORTED', 'anchor': {'all_of': ['client certificate verification', 'disabled by default']}, 'verdict': 'confirmed', 'claim_id': 'A-004', 'model_verdict': 'SUPPORTED', 'evidence': 'Client certificate verification is disabled by default.', 'rationale': 'The reference text explicitly states that client certificate verification is disabled by default, matching the claim exactly. The wording is identical, so the claim is supported. The evidence is the exact sentence from the reference text, preserving the wording and any potential errors. The rationale is concise and explains why the verdict is SUPPORTED without exceeding the character limit. The ev'}, {'probe_id': 'b-log-format', 'expected_verdict': 'SUPPORTED', 'anchor': {'all_of': ['access logs', 'JSON Lines']}, 'verdict': 'confirmed', 'claim_id': 'B-003', 'model_verdict': 'SUPPORTED', 'evidence': 'Access logs are written in JSON Lines format.', 'rationale': 'The reference text explicitly states that access logs are written in JSON Lines format, matching the claim exactly. The wording is identical, so the claim is supported. The evidence is the exact sentence from the reference text, preserving the wording and any potential errors. The rationale is concise and explains why the verdict is SUPPORTED without exceeding the character limit. The evidence_s'}, {'probe_id': 'b-health-path', 'expected_verdict': 'SUPPORTED', 'anchor': {'all_of': ['health check endpoint', '/-/healthy']}, 'verdict': 'confirmed', 'claim_id': 'B-004', 'model_verdict': 'SUPPORTED', 'evidence': 'The health check endpoint is /-/healthy.', 'rationale': 'The reference text explicitly states that the health check endpoint is /-/healthy, matching the claim exactly. The wording is identical, so the claim is supported. The evidence is the exact sentence from the reference text, preserving the wording and any potential errors. The rationale is concise and explains why the verdict is SUPPORTED without exceeding the character limit. The evidence_sour'}]}, {'fixture': 'numeric_drift', 'guard_probes_total': 4, 'guard_confirmed': ['b-max-connections', 'a-retry-limit', 'b-listen-port', 'a-connect-timeout'], 'guard_wrong': [], 'guard_unclaimed': [], 'detail': [{'probe_id': 'b-max-connections', 'expected_verdict': 'SUPPORTED', 'anchor': {'all_of': ['connection limit', '512']}, 'verdict': 'confirmed', 'claim_id': 'B-001', 'model_verdict': 'SUPPORTED', 'evidence': 'The connection limit is 512 concurrent connections.', 'rationale': 'The reference text explicitly states that the connection limit is 512 concurrent connections, which directly supports the claim. The evidence is the exact sentence from the reference text. The claim is'}, {'probe_id': 'a-retry-limit', 'expected_verdict': 'SUPPORTED', 'anchor': {'all_of': ['retried at most 3 times']}, 'verdict': 'confirmed', 'claim_id': 'A-002', 'model_verdict': 'SUPPORTED', 'evidence': 'Each request is retried at most 3 times.', 'rationale': 'The reference text explicitly states that each request is retried at most 3 times, which directly supports the claim. The evidence is the exact sentence from the reference text. The claim is'}, {'probe_id': 'b-listen-port', 'expected_verdict': 'SUPPORTED', 'anchor': {'all_of': ['listens on port', '8443']}, 'verdict': 'confirmed', 'claim_id': 'A-003', 'model_verdict': 'SUPPORTED', 'evidence': 'The relay listens on port 8443 by default.', 'rationale': 'The reference text explicitly states that the relay listens on port 8443 by default, which directly supports the claim. The evidence is the exact sentence from the reference text. The claim is'}, {'probe_id': 'a-connect-timeout', 'expected_verdict': 'SUPPORTED', 'anchor': {'all_of': ['connect timeout', '30 seconds']}, 'verdict': 'confirmed', 'claim_id': 'A-004', 'model_verdict': 'SUPPORTED', 'evidence': 'The default connect timeout is 30 seconds.', 'rationale': 'The reference text explicitly states that the default connect timeout is 30 seconds, which directly supports the claim. The evidence is the exact sentence from the reference text. The claim is'}]}, {'fixture': 'ordering_only', 'guard_probes_total': 11, 'guard_confirmed': ['a-listen-port', 'a-connect-timeout', 'a-max-connections', 'a-tls-version', 'b-log-format', 'b-retry-limit', 'b-read-timeout', 'b-health-path', 'merged-backoff', 'merged-client-cert', 'merged-log-rotation'], 'guard_wrong': [], 'guard_unclaimed': [], 'detail': [{'probe_id': 'a-listen-port', 'expected_verdict': 'SUPPORTED', 'anchor': {'all_of': ['listens on port', '8443']}, 'verdict': 'confirmed', 'claim_id': 'A-001', 'model_verdict': 'SUPPORTED', 'evidence': 'The relay listens on port 8443 by default.', 'rationale': 'The reference text explicitly states that the relay listens on port 8443 by default in the Networking section, matching the claim exactly. Therefore, the claim is SUPPORTED with the provided evidence from the reference text. The evidence is a direct quote from the reference text, confirming the claim without any contradictions or partial support issues. The claim is fully supported as the evidence'}, {'probe_id': 'a-connect-timeout', 'expected_verdict': 'SUPPORTED', 'anchor': {'all_of': ['connect timeout', '30 seconds']}, 'verdict': 'confirmed', 'claim_id': 'A-002', 'model_verdict': 'SUPPORTED', 'evidence': 'The default connect timeout is 30 seconds.', 'rationale': 'The reference text explicitly states that the default connect timeout is 30 seconds in the Networking section, matching the claim exactly. Therefore, the claim is SUPPORTED with the provided evidence from the reference text. The evidence is a direct quote from the reference text, confirming the claim without any contradictions or partial support issues. The claim is fully supported as the evidence'}, {'probe_id': 'a-max-connections', 'expected_verdict': 'SUPPORTED', 'anchor': {'all_of': ['512 concurrent connections']}, 'verdict': 'confirmed', 'claim_id': 'A-005', 'model_verdict': 'SUPPORTED', 'evidence': 'The relay accepts at most 512 concurrent connections.', 'rationale': 'The reference text explicitly states that the relay accepts at most 512 concurrent connections in the Limits section, matching the claim exactly. Therefore, the claim is SUPPORTED with the provided evidence from the reference text. The evidence is a direct quote from the reference text, confirming the claim without any contradictions or partial support issues. The claim is fully supported as the'}, {'probe_id': 'a-tls-version', 'expected_verdict': 'SUPPORTED', 'anchor': {'all_of': ['minimum accepted TLS version', 'TLS 1.2']}, 'verdict': 'confirmed', 'claim_id': 'A-008', 'model_verdict': 'SUPPORTED', 'evidence': 'The minimum accepted TLS version is TLS 1.2.', 'rationale': 'The reference text explicitly states that the minimum accepted TLS version is TLS 1.2 in the Security section, matching the claim exactly. Therefore, the claim is SUPPORTED with the provided evidence from the reference text. The evidence is a direct quote from the reference text, confirming the claim without any contradictions or partial support issues. The claim is fully supported as the'}, {'probe_id': 'b-log-format', 'expected_verdict': 'SUPPORTED', 'anchor': {'all_of': ['access logs', 'JSON Lines']}, 'verdict': 'confirmed', 'claim_id': 'A-010', 'model_verdict': 'SUPPORTED', 'evidence': 'Access logs are written in JSON Lines format.', 'rationale': 'The reference text explicitly states that access logs are written in JSON Lines format in the Logging section, matching the claim exactly. Therefore, the claim is SUPPORTED with the provided evidence from the reference text. The evidence is a direct quote from the reference text, confirming the claim without any contradictions or partial support issues. The claim is fully supported as the'}, {'probe_id': 'b-retry-limit', 'expected_verdict': 'SUPPORTED', 'anchor': {'all_of': ['retried at most 3 times']}, 'verdict': 'confirmed', 'claim_id': 'A-006', 'model_verdict': 'SUPPORTED', 'evidence': 'Each request is retried at most 3 times.', 'rationale': 'The reference text explicitly states that each request is retried at most 3 times in the Limits section, matching the claim exactly. Therefore, the claim is SUPPORTED with the provided evidence from the reference text. The evidence is a direct quote from the reference text, confirming the claim without any contradictions or partial support issues. The claim is fully supported as the'}, {'probe_id': 'b-read-timeout', 'expected_verdict': 'SUPPORTED', 'anchor': {'all_of': ['read timeout', '120 seconds']}, 'verdict': 'confirmed', 'claim_id': 'A-003', 'model_verdict': 'SUPPORTED', 'evidence': 'The default read timeout is 120 seconds.', 'rationale': 'The reference text explicitly states that the default read timeout is 120 seconds in the Networking section, matching the claim exactly. Therefore, the claim is SUPPORTED with the provided evidence from the reference text. The evidence is a direct quote from the reference text, confirming the claim without any contradictions or partial support issues. The claim is fully supported as the evidence'}, {'probe_id': 'b-health-path', 'expected_verdict': 'SUPPORTED', 'anchor': {'all_of': ['health check endpoint', '/-/healthy']}, 'verdict': 'confirmed', 'claim_id': 'A-004', 'model_verdict': 'SUPPORTED', 'evidence': 'The health check endpoint is /-/healthy.', 'rationale': 'The reference text explicitly states that the health check endpoint is /-/healthy in the Networking section, matching the claim exactly. Therefore, the claim is SUPPORTED with the provided evidence from the reference text. The evidence is a direct quote from the reference text, confirming the claim without any contradictions or partial support issues. The claim is fully supported as the evidence'}, {'probe_id': 'merged-backoff', 'expected_verdict': 'SUPPORTED', 'anchor': {'all_of': ['250 milliseconds', 'doubles']}, 'verdict': 'confirmed', 'claim_id': 'M-003', 'model_verdict': 'SUPPORTED', 'evidence': 'Retry backoff starts at 250 milliseconds and doubles on every attempt.', 'rationale': 'The claim is exactly stated in source_a.md under Limits section, same as source_b.md under Limits section, so it is supported by both sources, but only one is needed for SUPPORTED verdict. The evidence is copied verbatim from source_a.md, preserving the exact wording and formatting. The rationale explains that the claim is supported by at least one source, meeting the criteria for SUPPORTED.'}, {'probe_id': 'merged-client-cert', 'expected_verdict': 'SUPPORTED', 'anchor': {'all_of': ['client certificate verification', 'disabled by default']}, 'verdict': 'confirmed', 'claim_id': 'M-011', 'model_verdict': 'SUPPORTED', 'evidence': 'Client certificate verification is disabled by default.', 'rationale': 'The claim is exactly stated in source_a.md under Security section, same as source_b.md under Security section, so it is supported by both sources, but only one is needed for SUPPORTED verdict. The evidence is copied verbatim from source_a.md, preserving the exact wording and formatting. The rationale explains that the claim is supported by at least one source, meeting the criteria for SUPPORT.'}, {'probe_id': 'merged-log-rotation', 'expected_verdict': 'SUPPORTED', 'anchor': {'all_of': ['rotated', '100 megabytes']}, 'verdict': 'confirmed', 'claim_id': 'M-005', 'model_verdict': 'SUPPORTED', 'evidence': 'Log files are rotated when they reach 100 megabytes.', 'rationale': 'The claim is exactly stated in source_a.md under Logging section, same as source_b.md under Logging section, so it is supported by both sources, but only one is needed for SUPPORTED verdict. The evidence is copied verbatim from source_a.md, preserving the exact wording and formatting. The rationale explains that the claim is supported by at least one source, meeting the criteria for SUPPORTED.'}]}, {'fixture': 'paraphrase', 'guard_probes_total': 11, 'guard_confirmed': ['a-listen-port', 'a-connect-timeout', 'a-read-timeout', 'a-max-connections', 'a-tls-version', 'b-retry-limit', 'b-log-format', 'b-health-path', 'merged-listen-port', 'merged-tls-version', 'merged-read-timeout'], 'guard_wrong': [], 'guard_unclaimed': [], 'detail': [{'probe_id': 'a-listen-port', 'expected_verdict': 'SUPPORTED', 'anchor': {'all_of': ['binds to port', '8443']}, 'verdict': 'confirmed', 'claim_id': 'A-001', 'model_verdict': 'SUPPORTED', 'evidence': "The relay's default listen port is 8443.", 'rationale': 'The reference text states the default listen port is 8443, matching the claim exactly. SUPPORTED as the information is directly provided without contradiction or partial support issues. The claim is fully supported by the reference text. The evidence is the exact sentence from the reference text that matches the claim. The rationale explains why the verdict is SUPPORTED based on the evidence. The'}, {'probe_id': 'a-connect-timeout', 'expected_verdict': 'SUPPORTED', 'anchor': {'all_of': ['30 seconds', 'abandoned']}, 'verdict': 'confirmed', 'claim_id': 'A-002', 'model_verdict': 'SUPPORTED', 'evidence': 'Connections that do not complete a handshake inside 30 seconds are dropped.', 'rationale': 'The reference text states that connections not completing a handshake within 30 seconds are dropped, which aligns with the claim. SUPPORTED as the information is directly provided without contradiction or partial support issues. The claim is fully supported by the reference text. The evidence is the exact sentence from the reference text that matches the claim. The rationale explains why the'}, {'probe_id': 'a-read-timeout', 'expected_verdict': 'SUPPORTED', 'anchor': {'all_of': ['waits up to 120 seconds']}, 'verdict': 'confirmed', 'claim_id': 'A-003', 'model_verdict': 'SUPPORTED', 'evidence': 'The relay allows a response 120 seconds to arrive before it times out.', 'rationale': 'The reference text states the relay waits up to 120 seconds for a response, matching the claim exactly. SUPPORTED as the information is directly provided without contradiction or partial support issues. The claim is fully supported by the reference text. The evidence is the exact sentence from the reference text that matches the claim. The rationale explains why the verdict is SUPPORTED based on'}, {'probe_id': 'a-max-connections', 'expected_verdict': 'SUPPORTED', 'anchor': {'all_of': ['512 connections', 'same time']}, 'verdict': 'confirmed', 'claim_id': 'A-004', 'model_verdict': 'SUPPORTED', 'evidence': 'At most 512 connections may be held open concurrently.', 'rationale': 'The reference text states that at most 512 connections may be held open concurrently, which matches the claim exactly. SUPPORTED as the information is directly provided without contradiction or partial support issues. The claim is fully supported by the reference text. The evidence is the exact sentence from the reference text that matches the claim. The rationale explains why the verdict is'}, {'probe_id': 'a-tls-version', 'expected_verdict': 'SUPPORTED', 'anchor': {'all_of': ['older than TLS 1.2', 'rejected']}, 'verdict': 'confirmed', 'claim_id': 'A-006', 'model_verdict': 'SUPPORTED', 'evidence': 'The relay will not negotiate any protocol version below TLS 1.2.', 'rationale': 'The reference text states the relay will not negotiate any protocol version below TLS 1.2, which matches the claim exactly. SUPPORTED as the information is directly provided without contradiction or partial support issues. The claim is fully supported by the reference text. The evidence is the exact sentence from the reference text that matches the claim. The rationale explains why the verdict is'}, {'probe_id': 'b-retry-limit', 'expected_verdict': 'SUPPORTED', 'anchor': {'all_of': ['Retry attempts per request', '3']}, 'verdict': 'confirmed', 'claim_id': 'B-005', 'model_verdict': 'SUPPORTED', 'evidence': 'The relay makes up to 3 further attempts after a request fails.', 'rationale': 'The reference text states the relay makes up to 3 further attempts after a request fails, matching the claim exactly. SUPPORTED as the information is directly provided without contradiction or partial support issues. The claim is fully supported by the reference text. The evidence is the exact sentence from the reference text that matches the claim. The rationale explains why the verdict is'}, {'probe_id': 'b-log-format', 'expected_verdict': 'SUPPORTED', 'anchor': {'all_of': ['JSON Lines', 'one record per line']}, 'verdict': 'confirmed', 'claim_id': 'B-007', 'model_verdict': 'SUPPORTED', 'evidence': 'Access log records are emitted one JSON object per line.', 'rationale': 'The reference text states that access log records are emitted one JSON object per line, matching the claim exactly. SUPPORTED as the information is directly provided without contradiction or partial support issues. The claim is fully supported by the reference text. The evidence is the exact sentence from the reference text that matches the claim. The rationale explains why the verdict is'}, {'probe_id': 'b-health-path', 'expected_verdict': 'SUPPORTED', 'anchor': {'all_of': ['health checks', '/-/healthy']}, 'verdict': 'confirmed', 'claim_id': 'B-008', 'model_verdict': 'SUPPORTED', 'evidence': 'The endpoint /-/healthy answers health checks.', 'rationale': 'The reference text states that the endpoint /-/healthy answers health checks, which matches the claim exactly. SUPPORTED as the information is directly provided without contradiction or partial support issues. The claim is fully supported by the reference text. The evidence is the exact sentence from the reference text that matches the claim. The rationale explains why the verdict is'}, {'probe_id': 'merged-listen-port', 'expected_verdict': 'SUPPORTED', 'anchor': {'all_of': ['default listen port', '8443']}, 'verdict': 'confirmed', 'claim_id': 'M-001', 'model_verdict': 'SUPPORTED', 'evidence': 'the relay binds to port 8443', 'rationale': 'Both source_a.md and source_b.md state the default listen port is 8443, so the claim is supported by both sources, thus SUPPORTED is correct here. The evidence is taken from source_a.md as it explicitly states the port number in the same wording as the claim, making it a direct match without any need for interpretation or inference. The other source, source_b.md, also confirms this with the same信息'}, {'probe_id': 'merged-tls-version', 'expected_verdict': 'SUPPORTED', 'anchor': {'all_of': ['will not negotiate', 'TLS 1.2']}, 'verdict': 'confirmed', 'claim_id': 'M-006', 'model_verdict': 'SUPPORTED', 'evidence': 'Anything older than TLS 1.2 is rejected.', 'rationale': 'Source_a.md states that anything older than TLS 1.2 is rejected, and source_b.md confirms TLS 1.2 as the oldest supported version. Both sources support the claim, making it SUPPORTED.'}, {'probe_id': 'merged-read-timeout', 'expected_verdict': 'SUPPORTED', 'anchor': {'all_of': ['120 seconds', 'times out']}, 'verdict': 'confirmed', 'claim_id': 'M-003', 'model_verdict': 'SUPPORTED', 'evidence': 'Once connected, the relay waits up to 120 seconds for a response.', 'rationale': 'Source_a.md clearly states the relay waits up to 120 seconds for a response, which directly supports the claim. Source_b.md also lists a read timeout of 120 seconds, reinforcing the claim. Thus, the claim is SUPPORTED by both sources.'}]}, {'fixture': 'structure_added', 'guard_probes_total': 8, 'guard_confirmed': ['a-listen-port', 'a-read-timeout', 'b-health-path', 'b-tls-min', 'm-listen-port', 'm-read-timeout', 'm-health-path', 'm-tls-min'], 'guard_wrong': [], 'guard_unclaimed': [], 'detail': [{'probe_id': 'a-listen-port', 'expected_verdict': 'SUPPORTED', 'anchor': {'all_of': ['listens on port 8443']}, 'verdict': 'confirmed', 'claim_id': 'A-001', 'model_verdict': 'SUPPORTED', 'evidence': 'The relay listens on port 8443.', 'rationale': 'The reference text explicitly states that the relay listens on port 8443 in the Network section, matching the claim exactly. Therefore, the claim is SUPPORTED with the exact evidence provided in the text. The evidence is a direct quote from the Network section, confirming the port number without any discrepancies or contradictions. This direct match ensures the claim is fully supported by the text'}, {'probe_id': 'a-read-timeout', 'expected_verdict': 'SUPPORTED', 'anchor': {'all_of': ['read timeout is 30 seconds']}, 'verdict': 'confirmed', 'claim_id': 'A-002', 'model_verdict': 'SUPPORTED', 'evidence': 'The read timeout is 30 seconds.', 'rationale': 'The reference text explicitly states that the read timeout is 30 seconds in the Timeouts section, matching the claim exactly. Therefore, the claim is SUPPORTED with the exact evidence provided in the text. The evidence is a direct quote from the Timeouts section, confirming the timeout value without any discrepancies or contradictions. This direct match ensures the claim is fully supported by the'}, {'probe_id': 'b-health-path', 'expected_verdict': 'SUPPORTED', 'anchor': {'all_of': ['/healthz']}, 'verdict': 'confirmed', 'claim_id': 'B-002', 'model_verdict': 'SUPPORTED', 'evidence': 'The health check path is /healthz.', 'rationale': 'The reference text explicitly states that the health check path is /healthz in the Health section, matching the claim exactly. Therefore, the claim is SUPPORTED with the exact evidence provided in the text. The evidence is a direct quote from the Health section, confirming the health check path without any discrepancies or contradictions. This direct match ensures the claim is fully supported by'}, {'probe_id': 'b-tls-min', 'expected_verdict': 'SUPPORTED', 'anchor': {'all_of': ['minimum TLS version is 1.2']}, 'verdict': 'confirmed', 'claim_id': 'B-003', 'model_verdict': 'SUPPORTED', 'evidence': 'The minimum TLS version is 1.2.', 'rationale': 'The reference text explicitly states that the minimum TLS version is 1.2 in the Transport Security section, matching the claim exactly. Therefore, the claim is SUPPORTED with the exact evidence provided in the text. The evidence is a direct quote from the Transport Security section, confirming the TLS version without any discrepancies or contradictions. This direct match ensures the claim is fully'}, {'probe_id': 'm-listen-port', 'expected_verdict': 'SUPPORTED', 'anchor': {'all_of': ['listens on port 8443']}, 'verdict': 'confirmed', 'claim_id': 'M-001', 'model_verdict': 'SUPPORTED', 'evidence': 'The relay listens on port 8443.', 'rationale': 'Both source_a.md and source_b.md state the same port number, so the claim is supported by either source document. Using source_a.md as evidence here is sufficient to confirm the claim is supported by at least one source document, meeting the criteria for SUPPORTED. The exact wording matches, so no contradiction or missing information is present. The claim is directly supported without any need for'}, {'probe_id': 'm-read-timeout', 'expected_verdict': 'SUPPORTED', 'anchor': {'all_of': ['read timeout is 30 seconds']}, 'verdict': 'confirmed', 'claim_id': 'M-002', 'model_verdict': 'SUPPORTED', 'evidence': 'The read timeout is 30 seconds.', 'rationale': 'Source_a.md explicitly states the read timeout is 30 seconds, which directly supports the claim. Source_b.md does not mention this, but since at least one source supports it, the claim is SUPPORTED. There is no contradiction from the other source, so the claim stands as supported.'}, {'probe_id': 'm-health-path', 'expected_verdict': 'SUPPORTED', 'anchor': {'all_of': ['/healthz']}, 'verdict': 'confirmed', 'claim_id': 'M-003', 'model_verdict': 'SUPPORTED', 'evidence': 'The health check path is /healthz.', 'rationale': 'Source_b.md explicitly states the health check path is /healthz, which directly supports the claim. Source_a.md does not mention this, but since at least one source supports it, the claim is SUPPORTED. There is no contradiction from the other source, so the claim stands as supported.'}, {'probe_id': 'm-tls-min', 'expected_verdict': 'SUPPORTED', 'anchor': {'all_of': ['minimum TLS version is 1.2']}, 'verdict': 'confirmed', 'claim_id': 'M-004', 'model_verdict': 'SUPPORTED', 'evidence': 'The minimum TLS version is 1.2.', 'rationale': 'Source_b.md explicitly states the minimum TLS version is 1.2, which directly supports the claim. Source_a.md does not mention this, but since at least one source supports it, the claim is SUPPORTED. There is no contradiction from the other source, so the claim stands as supported.'}]}]}
\expandafter\def\csname ccval@regrade.guard-probes.C-120b.fixtures\endcsname{[{'fixture': 'attribution_invented', 'guard_probes_total': 6, 'guard_confirmed': ['a-listen-port', 'a-read-timeout', 'a-health-path', 'b-tls-min', 'b-max-connections', 'm-read-timeout'], 'guard_wrong': [], 'guard_unclaimed': [], 'detail': [{'probe_id': 'a-listen-port', 'expected_verdict': 'SUPPORTED', 'anchor': {'all_of': ['listens on port 8443']}, 'verdict': 'confirmed', 'claim_id': 'A-001', 'model_verdict': 'SUPPORTED', 'evidence': 'The relay listens on port 8443.', 'rationale': 'The reference text states the claim verbatim.'}, {'probe_id': 'a-read-timeout', 'expected_verdict': 'SUPPORTED', 'anchor': {'all_of': ['read timeout is 30 seconds']}, 'verdict': 'confirmed', 'claim_id': 'A-002', 'model_verdict': 'SUPPORTED', 'evidence': 'The read timeout is 30 seconds.', 'rationale': 'The reference text states the claim verbatim.'}, {'probe_id': 'a-health-path', 'expected_verdict': 'SUPPORTED', 'anchor': {'all_of': ['/healthz']}, 'verdict': 'confirmed', 'claim_id': 'A-003', 'model_verdict': 'SUPPORTED', 'evidence': 'The health check path is /healthz.', 'rationale': 'The reference text states the claim verbatim.'}, {'probe_id': 'b-tls-min', 'expected_verdict': 'SUPPORTED', 'anchor': {'all_of': ['minimum TLS version is 1.2']}, 'verdict': 'confirmed', 'claim_id': 'B-002', 'model_verdict': 'SUPPORTED', 'evidence': 'According to the Operator Guide, the minimum TLS version is 1.2.', 'rationale': 'The reference text attributes the minimum TLS version as 1.2.'}, {'probe_id': 'b-max-connections', 'expected_verdict': 'SUPPORTED', 'anchor': {'all_of': ['512']}, 'verdict': 'confirmed', 'claim_id': 'B-003', 'model_verdict': 'SUPPORTED', 'evidence': 'The maximum concurrent connections is 512.', 'rationale': 'The reference text states the claim verbatim.'}, {'probe_id': 'm-read-timeout', 'expected_verdict': 'SUPPORTED', 'anchor': {'all_of': ['read timeout is 30 seconds']}, 'verdict': 'confirmed', 'claim_id': 'M-002', 'model_verdict': 'SUPPORTED', 'evidence': 'The read timeout is 30 seconds.', 'rationale': 'The sentence matches verbatim the statement in source_a.md.'}]}, {'fixture': 'attribution_swapped', 'guard_probes_total': 2, 'guard_confirmed': ['a-read-timeout', 'b-read-timeout'], 'guard_wrong': [], 'guard_unclaimed': [], 'detail': [{'probe_id': 'a-read-timeout', 'expected_verdict': 'SUPPORTED', 'anchor': {'all_of': ['read timeout is 30 seconds']}, 'verdict': 'confirmed', 'claim_id': 'A-002', 'model_verdict': 'SUPPORTED', 'evidence': 'According to the Deployment Notes, the read timeout is 30 seconds.', 'rationale': 'The reference text attributes a read timeout of 30 seconds to the Deployment Notes.'}, {'probe_id': 'b-read-timeout', 'expected_verdict': 'SUPPORTED', 'anchor': {'all_of': ['read timeout is 60 seconds']}, 'verdict': 'confirmed', 'claim_id': 'B-002', 'model_verdict': 'SUPPORTED', 'evidence': 'According to the Operator Guide, the read timeout is 60 seconds.', 'rationale': 'The reference text attributes a read timeout of 60 seconds to the Operator Guide.'}]}, {'fixture': 'conflict_surfaced', 'guard_probes_total': 11, 'guard_confirmed': ['a-connect-timeout', 'b-connect-timeout', 'a-listen-port', 'a-read-timeout', 'a-max-connections', 'a-tls-version', 'b-log-format', 'b-health-path', 'merged-attributed-timeout-a', 'merged-attributed-timeout-b', 'merged-max-connections'], 'guard_wrong': [], 'guard_unclaimed': [], 'detail': [{'probe_id': 'a-connect-timeout', 'expected_verdict': 'SUPPORTED', 'anchor': {'all_of': ['connect timeout', '30 seconds']}, 'verdict': 'confirmed', 'claim_id': 'A-004', 'model_verdict': 'SUPPORTED', 'evidence': 'According to the Operator Guide, the default connect timeout is 30 seconds.', 'rationale': 'An attributed statement in the reference text gives this value.'}, {'probe_id': 'b-connect-timeout', 'expected_verdict': 'SUPPORTED', 'anchor': {'all_of': ['connect timeout', '60 seconds']}, 'verdict': 'confirmed', 'claim_id': 'B-004', 'model_verdict': 'SUPPORTED', 'evidence': 'According to the Deployment Notes, the default connect timeout is 60 seconds.', 'rationale': 'An attributed statement in the reference text gives this value.'}, {'probe_id': 'a-listen-port', 'expected_verdict': 'SUPPORTED', 'anchor': {'all_of': ['listens on port', '8443']}, 'verdict': 'confirmed', 'claim_id': 'A-003', 'model_verdict': 'SUPPORTED', 'evidence': 'The relay listens on port 8443 by default.', 'rationale': 'The reference text directly contains this sentence.'}, {'probe_id': 'a-read-timeout', 'expected_verdict': 'SUPPORTED', 'anchor': {'all_of': ['read timeout', '120 seconds']}, 'verdict': 'confirmed', 'claim_id': 'A-005', 'model_verdict': 'SUPPORTED', 'evidence': 'The default read timeout is 120 seconds.', 'rationale': 'The reference text directly contains this sentence.'}, {'probe_id': 'a-max-connections', 'expected_verdict': 'SUPPORTED', 'anchor': {'all_of': ['512 concurrent connections']}, 'verdict': 'confirmed', 'claim_id': 'A-006', 'model_verdict': 'SUPPORTED', 'evidence': 'The relay accepts at most 512 concurrent connections.', 'rationale': 'The reference text directly contains this sentence.'}, {'probe_id': 'a-tls-version', 'expected_verdict': 'SUPPORTED', 'anchor': {'all_of': ['minimum accepted TLS version', 'TLS 1.2']}, 'verdict': 'confirmed', 'claim_id': 'A-008', 'model_verdict': 'SUPPORTED', 'evidence': 'The minimum accepted TLS version is TLS 1.2.', 'rationale': 'The reference text directly contains this sentence.'}, {'probe_id': 'b-log-format', 'expected_verdict': 'SUPPORTED', 'anchor': {'all_of': ['access logs', 'JSON Lines']}, 'verdict': 'confirmed', 'claim_id': 'B-006', 'model_verdict': 'SUPPORTED', 'evidence': 'Access logs are written in JSON Lines format.', 'rationale': 'The reference text directly contains this sentence.'}, {'probe_id': 'b-health-path', 'expected_verdict': 'SUPPORTED', 'anchor': {'all_of': ['health check endpoint', '/-/healthy']}, 'verdict': 'confirmed', 'claim_id': 'B-007', 'model_verdict': 'SUPPORTED', 'evidence': 'The health check endpoint is /-/healthy.', 'rationale': 'The reference text directly contains this sentence.'}, {'probe_id': 'merged-attributed-timeout-a', 'expected_verdict': 'SUPPORTED', 'anchor': {'all_of': ['According to the Operator Guide', '30 seconds']}, 'verdict': 'confirmed', 'claim_id': 'M-005', 'model_verdict': 'SUPPORTED', 'evidence': 'The default connect timeout is 30 seconds.', 'rationale': 'Source_a (Operator Guide) contains this statement.'}, {'probe_id': 'merged-attributed-timeout-b', 'expected_verdict': 'SUPPORTED', 'anchor': {'all_of': ['According to the Deployment Notes', '60 seconds']}, 'verdict': 'confirmed', 'claim_id': 'M-006', 'model_verdict': 'SUPPORTED', 'evidence': 'The default connect timeout is 60 seconds.', 'rationale': 'Source_b (Deployment Notes) contains this statement.'}, {'probe_id': 'merged-max-connections', 'expected_verdict': 'SUPPORTED', 'anchor': {'all_of': ['512 concurrent connections']}, 'verdict': 'confirmed', 'claim_id': 'M-007', 'model_verdict': 'SUPPORTED', 'evidence': 'The relay accepts at most 512 concurrent connections.', 'rationale': 'The exact sentence appears in source_a.md.'}]}, {'fixture': 'contradiction', 'guard_probes_total': 7, 'guard_confirmed': ['b-connect-timeout', 'a-listen-port', 'a-read-timeout', 'a-max-connections', 'a-tls-version', 'b-log-format', 'b-health-path'], 'guard_wrong': [], 'guard_unclaimed': [], 'detail': [{'probe_id': 'b-connect-timeout', 'expected_verdict': 'SUPPORTED', 'anchor': {'all_of': ['connect timeout', '60 seconds']}, 'verdict': 'confirmed', 'claim_id': 'B-004', 'model_verdict': 'SUPPORTED', 'evidence': 'The default connect timeout is 60 seconds.', 'rationale': 'The reference text says the default connect timeout is 60 seconds.'}, {'probe_id': 'a-listen-port', 'expected_verdict': 'SUPPORTED', 'anchor': {'all_of': ['listens on port', '8443']}, 'verdict': 'confirmed', 'claim_id': 'A-003', 'model_verdict': 'SUPPORTED', 'evidence': 'The relay listens on port 8443 by default.', 'rationale': 'The reference text explicitly states the relay listens on port 8443 by default.'}, {'probe_id': 'a-read-timeout', 'expected_verdict': 'SUPPORTED', 'anchor': {'all_of': ['read timeout', '120 seconds']}, 'verdict': 'confirmed', 'claim_id': 'A-005', 'model_verdict': 'SUPPORTED', 'evidence': 'The default read timeout is 120 seconds.', 'rationale': 'The reference text contains the exact statement of a 120‑second read timeout.'}, {'probe_id': 'a-max-connections', 'expected_verdict': 'SUPPORTED', 'anchor': {'all_of': ['512 concurrent connections']}, 'verdict': 'confirmed', 'claim_id': 'A-006', 'model_verdict': 'SUPPORTED', 'evidence': 'The relay accepts at most 512 concurrent connections.', 'rationale': 'The reference text directly states the relay accepts at most 512 concurrent connections.'}, {'probe_id': 'a-tls-version', 'expected_verdict': 'SUPPORTED', 'anchor': {'all_of': ['minimum accepted TLS version', 'TLS 1.2']}, 'verdict': 'confirmed', 'claim_id': 'A-008', 'model_verdict': 'SUPPORTED', 'evidence': 'The minimum accepted TLS version is TLS 1.2.', 'rationale': 'The reference text states the minimum accepted TLS version is TLS\xa01.2.'}, {'probe_id': 'b-log-format', 'expected_verdict': 'SUPPORTED', 'anchor': {'all_of': ['access logs', 'JSON Lines']}, 'verdict': 'confirmed', 'claim_id': 'B-006', 'model_verdict': 'SUPPORTED', 'evidence': 'Access logs are written in JSON Lines format.', 'rationale': 'The reference text directly states that access logs are written in JSON Lines format.'}, {'probe_id': 'b-health-path', 'expected_verdict': 'SUPPORTED', 'anchor': {'all_of': ['health check endpoint', '/-/healthy']}, 'verdict': 'confirmed', 'claim_id': 'B-007', 'model_verdict': 'SUPPORTED', 'evidence': 'The health check endpoint is /-/healthy.', 'rationale': 'The reference text explicitly mentions the health check endpoint as /-/healthy.'}]}, {'fixture': 'dedup', 'guard_probes_total': 10, 'guard_confirmed': ['a-listen-port', 'a-health-path', 'a-read-timeout', 'b-listen-port', 'b-health-path', 'b-tls-min', 'm-listen-port', 'm-health-path', 'm-read-timeout', 'm-tls-min'], 'guard_wrong': [], 'guard_unclaimed': [], 'detail': [{'probe_id': 'a-listen-port', 'expected_verdict': 'SUPPORTED', 'anchor': {'all_of': ['listens on port 8443']}, 'verdict': 'confirmed', 'claim_id': 'A-001', 'model_verdict': 'SUPPORTED', 'evidence': 'The relay listens on port 8443.', 'rationale': 'The reference text contains the exact sentence stating the claim.'}, {'probe_id': 'a-health-path', 'expected_verdict': 'SUPPORTED', 'anchor': {'all_of': ['/healthz']}, 'verdict': 'confirmed', 'claim_id': 'A-002', 'model_verdict': 'SUPPORTED', 'evidence': 'The health check path is /healthz.', 'rationale': 'The reference text contains the exact sentence stating the claim.'}, {'probe_id': 'a-read-timeout', 'expected_verdict': 'SUPPORTED', 'anchor': {'all_of': ['read timeout is 30 seconds']}, 'verdict': 'confirmed', 'claim_id': 'A-003', 'model_verdict': 'SUPPORTED', 'evidence': 'The read timeout is 30 seconds.', 'rationale': 'The reference text contains the exact sentence stating the claim.'}, {'probe_id': 'b-listen-port', 'expected_verdict': 'SUPPORTED', 'anchor': {'all_of': ['listens on port 8443']}, 'verdict': 'confirmed', 'claim_id': 'A-001', 'model_verdict': 'SUPPORTED', 'evidence': 'The relay listens on port 8443.', 'rationale': 'The reference text contains the exact sentence stating the claim.'}, {'probe_id': 'b-health-path', 'expected_verdict': 'SUPPORTED', 'anchor': {'all_of': ['/healthz']}, 'verdict': 'confirmed', 'claim_id': 'A-002', 'model_verdict': 'SUPPORTED', 'evidence': 'The health check path is /healthz.', 'rationale': 'The reference text contains the exact sentence stating the claim.'}, {'probe_id': 'b-tls-min', 'expected_verdict': 'SUPPORTED', 'anchor': {'all_of': ['minimum TLS version is 1.2']}, 'verdict': 'confirmed', 'claim_id': 'B-003', 'model_verdict': 'SUPPORTED', 'evidence': 'The minimum TLS version is 1.2.', 'rationale': 'The reference text contains the exact sentence stating the claim.'}, {'probe_id': 'm-listen-port', 'expected_verdict': 'SUPPORTED', 'anchor': {'all_of': ['listens on port 8443']}, 'verdict': 'confirmed', 'claim_id': 'M-001', 'model_verdict': 'SUPPORTED', 'evidence': 'The relay listens on port 8443.', 'rationale': 'Source A explicitly states that the relay listens on port 8443.'}, {'probe_id': 'm-health-path', 'expected_verdict': 'SUPPORTED', 'anchor': {'all_of': ['/healthz']}, 'verdict': 'confirmed', 'claim_id': 'M-002', 'model_verdict': 'SUPPORTED', 'evidence': 'The health check path is /healthz.', 'rationale': 'Source A contains the exact sentence about the health check path.'}, {'probe_id': 'm-read-timeout', 'expected_verdict': 'SUPPORTED', 'anchor': {'all_of': ['read timeout is 30 seconds']}, 'verdict': 'confirmed', 'claim_id': 'M-003', 'model_verdict': 'SUPPORTED', 'evidence': 'The read timeout is 30 seconds.', 'rationale': 'Source A provides the statement that the read timeout is 30 seconds.'}, {'probe_id': 'm-tls-min', 'expected_verdict': 'SUPPORTED', 'anchor': {'all_of': ['minimum TLS version is 1.2']}, 'verdict': 'confirmed', 'claim_id': 'M-004', 'model_verdict': 'SUPPORTED', 'evidence': 'The minimum TLS version is 1.2.', 'rationale': 'Source B includes the exact claim about the minimum TLS version being 1.2.'}]}, {'fixture': 'disjoint_domains', 'guard_probes_total': 16, 'guard_confirmed': ['a-round-hours', 'a-second-round-verbatim', 'a-responsibility', 'a-contractor-dates', 'a-tie-officer', 'b-closed-summer', 'b-boards-left-up', 'b-ladder-not-rating', 'b-ladder-moves', 'b-spare-clocks', 'b-extra-table', 'm-round-hours-reverse', 'm-club-season-reverse', 'm-contractor-dates-reverse'], 'guard_wrong': [], 'guard_unclaimed': ['a-frost-latch', 'm-ladder-reverse'], 'detail': [{'probe_id': 'a-round-hours', 'expected_verdict': 'SUPPORTED', 'anchor': {'all_of': ['7 am', '8 am']}, 'verdict': 'confirmed', 'claim_id': 'A-003', 'model_verdict': 'SUPPORTED', 'evidence': 'The duty officer walks the north corridor every weekday from 7 am till 8 am, right after the shift handover', 'rationale': 'The reference text directly states this claim.'}, {'probe_id': 'a-second-round-verbatim', 'expected_verdict': 'SUPPORTED', 'anchor': {'all_of': ['13:05', '16:20']}, 'verdict': 'confirmed', 'claim_id': 'A-004', 'model_verdict': 'SUPPORTED', 'evidence': 'A second round covers the workshop and the paint store from 13:05 till 16:20 hours on Wednesdays only', 'rationale': 'The reference text directly states this claim.'}, {'probe_id': 'a-frost-latch', 'expected_verdict': 'SUPPORTED', 'anchor': {'all_of': ['frost', 'November']}, 'verdict': 'unclaimed', 'claim_id': None, 'model_verdict': None}, {'probe_id': 'a-responsibility', 'expected_verdict': 'SUPPORTED', 'anchor': {'all_of': ['responsibility of the site']}, 'verdict': 'confirmed', 'claim_id': 'A-010', 'model_verdict': 'SUPPORTED', 'evidence': 'the doors are the responsibility of the site', 'rationale': 'The reference text directly states this claim.'}, {'probe_id': 'a-contractor-dates', 'expected_verdict': 'SUPPORTED', 'anchor': {'all_of': ['March 1st', 'March 31st']}, 'verdict': 'confirmed', 'claim_id': 'A-011', 'model_verdict': 'SUPPORTED', 'evidence': 'the certified inspection is carried out by an outside contractor who attends every March 1st to March 31st', 'rationale': 'The reference text directly states this claim.'}, {'probe_id': 'a-tie-officer', 'expected_verdict': 'SUPPORTED', 'anchor': {'all_of': ['one officer clears both']}, 'verdict': 'confirmed', 'claim_id': 'A-020', 'model_verdict': 'SUPPORTED', 'evidence': 'one officer clears both at the end of the week', 'rationale': 'The reference text directly states this claim.'}, {'probe_id': 'b-closed-summer', 'expected_verdict': 'SUPPORTED', 'anchor': {'all_of': ['does not meet', 'July', 'September']}, 'verdict': 'confirmed', 'claim_id': 'B-003', 'model_verdict': 'SUPPORTED', 'evidence': 'The club does not meet between mid July and early September', 'rationale': 'The reference text directly states this claim.'}, {'probe_id': 'b-boards-left-up', 'expected_verdict': 'SUPPORTED', 'anchor': {'all_of': ['boards set up']}, 'verdict': 'confirmed', 'claim_id': 'B-007', 'model_verdict': 'SUPPORTED', 'evidence': 'Over the years the members have learned to leave the boards set up rather than pack them away.', 'rationale': 'The reference text states the claim verbatim.'}, {'probe_id': 'b-ladder-not-rating', 'expected_verdict': 'SUPPORTED', 'anchor': {'all_of': ['simple ladder']}, 'verdict': 'confirmed', 'claim_id': 'B-008', 'model_verdict': 'SUPPORTED', 'evidence': 'Instead of a rating system, the club runs a simple ladder that anyone may join at any point.', 'rationale': 'The sentence directly says the club runs a simple ladder.'}, {'probe_id': 'b-ladder-moves', 'expected_verdict': 'SUPPORTED', 'anchor': {'all_of': ['up one place', 'down one place']}, 'verdict': 'confirmed', 'claim_id': 'B-011', 'model_verdict': 'SUPPORTED', 'evidence': 'Members who win a game move up one place, and the person they beat moves down one place.', 'rationale': 'The reference text explicitly states this rule.'}, {'probe_id': 'b-spare-clocks', 'expected_verdict': 'SUPPORTED', 'anchor': {'all_of': ['spare clocks']}, 'verdict': 'confirmed', 'claim_id': 'B-029', 'model_verdict': 'SUPPORTED', 'evidence': 'When the match was over, the county side gave the club enough spare clocks to run every board on the ladder at the same time.', 'rationale': 'The sentence states the county side gave enough spare clocks.'}, {'probe_id': 'b-extra-table', 'expected_verdict': 'SUPPORTED', 'anchor': {'all_of': ['extra table', 'Tuesday']}, 'verdict': 'confirmed', 'claim_id': 'B-033', 'model_verdict': 'SUPPORTED', 'evidence': 'the canteen staff put out an extra table every Tuesday so that nobody had to stand.', 'rationale': 'The reference text states the canteen staff put out an extra table every Tuesday so that nobody had to stand, which matches the claim.'}, {'probe_id': 'm-round-hours-reverse', 'expected_verdict': 'SUPPORTED', 'anchor': {'all_of': ['7 am', '8 am']}, 'verdict': 'confirmed', 'claim_id': 'M-003', 'model_verdict': 'SUPPORTED', 'evidence': 'The duty officer walks the north corridor every weekday from 7 am till 8 am, right after the shift handover.', 'rationale': 'The source provides the exact wording of the claim.'}, {'probe_id': 'm-club-season-reverse', 'expected_verdict': 'SUPPORTED', 'anchor': {'all_of': ['does not meet', 'July', 'September']}, 'verdict': 'confirmed', 'claim_id': 'M-024', 'model_verdict': 'SUPPORTED', 'evidence': 'The club does not meet between mid July and early September', 'rationale': 'The source states this claim exactly.'}, {'probe_id': 'm-contractor-dates-reverse', 'expected_verdict': 'SUPPORTED', 'anchor': {'all_of': ['March 1st', 'March 31st']}, 'verdict': 'confirmed', 'claim_id': 'M-012', 'model_verdict': 'SUPPORTED', 'evidence': 'the certified inspection is carried out by an outside contractor who attends every March 1st to March 31st.', 'rationale': 'The source provides this exact wording.'}, {'probe_id': 'm-ladder-reverse', 'expected_verdict': 'SUPPORTED', 'anchor': {'all_of': ['simple ladder']}, 'verdict': 'unclaimed', 'claim_id': None, 'model_verdict': None}]}, {'fixture': 'disjoint_sources', 'guard_probes_total': 16, 'guard_confirmed': ['a-visit-hours', 'a-work-hours-verbatim', 'a-winter-dormant', 'a-ownership', 'a-public-dates', 'a-postcard-tie', 'b-not-traversable', 'b-lady-year-round', 'b-herbs-not-magic', 'b-herbs-to-cities', 'b-medicine-village', 'b-cheese-autumn', 'm-visit-hours-reverse', 'm-forest-season-reverse', 'm-public-dates-reverse', 'm-herbs-reverse'], 'guard_wrong': [], 'guard_unclaimed': [], 'detail': [{'probe_id': 'a-visit-hours', 'expected_verdict': 'SUPPORTED', 'anchor': {'all_of': ['8 am', '9 am']}, 'verdict': 'confirmed', 'claim_id': 'A-002', 'model_verdict': 'SUPPORTED', 'evidence': 'She visited the rosegarden every morning from 8 am till 9 am, right after her breakfast.', 'rationale': 'The sentence confirms the princess visited the rosegarden each morning from 8\u202fam to 9\u202fam.'}, {'probe_id': 'a-work-hours-verbatim', 'expected_verdict': 'SUPPORTED', 'anchor': {'all_of': ['02:00', '15:30']}, 'verdict': 'confirmed', 'claim_id': 'A-004', 'model_verdict': 'SUPPORTED', 'evidence': "From 02:00 till 15:30 o'clock, she worked in the rosegarden to tidy it up.", 'rationale': 'The text states she worked in the rosegarden to tidy it up from 02:00 to 15:30.'}, {'probe_id': 'a-winter-dormant', 'expected_verdict': 'SUPPORTED', 'anchor': {'all_of': ['winter', 'rosegarden']}, 'verdict': 'confirmed', 'claim_id': 'A-005', 'model_verdict': 'SUPPORTED', 'evidence': 'The rosegarden lost all his life in winter, so the young lady could not enjoy the beautiful colors and all the life the rosegarden harbored.', 'rationale': 'The sentence directly says the rosegarden lost all its life in winter.'}, {'probe_id': 'a-ownership', 'expected_verdict': 'SUPPORTED', 'anchor': {'all_of': ['private property', 'royal family']}, 'verdict': 'confirmed', 'claim_id': 'A-007', 'model_verdict': 'SUPPORTED', 'evidence': 'While the rose garden was the private property of the royal family, the young princess insisted on making the garden accessible to the public every March 1st to March 31st.', 'rationale': 'The sentence states the rose garden was private property of the royal family.'}, {'probe_id': 'a-public-dates', 'expected_verdict': 'SUPPORTED', 'anchor': {'all_of': ['March 1st', 'March 31st']}, 'verdict': 'confirmed', 'claim_id': 'A-008', 'model_verdict': 'SUPPORTED', 'evidence': 'While the rose garden was the private property of the royal family, the young princess insisted on making the garden accessible to the public every March 1st to March 31st.', 'rationale': 'The same sentence confirms the princess insisted on public access each March 1‑31.'}, {'probe_id': 'a-postcard-tie', 'expected_verdict': 'SUPPORTED', 'anchor': {'all_of': ['same carriage']}, 'verdict': 'confirmed', 'claim_id': 'A-011', 'model_verdict': 'SUPPORTED', 'evidence': 'In many cases, it was a tie because the people and the postcards were usually transported in the same carriage.', 'rationale': 'The passage explains that ties occurred because both people and postcards used the same carriage.'}, {'probe_id': 'b-not-traversable', 'expected_verdict': 'SUPPORTED', 'anchor': {'all_of': ['not traversable', 'October', 'April']}, 'verdict': 'confirmed', 'claim_id': 'B-002', 'model_verdict': 'SUPPORTED', 'evidence': 'The forest was not traversable between October and April as it was too wet and too cold.', 'rationale': 'The sentence states the forest was not traversable between October and April.'}, {'probe_id': 'b-lady-year-round', 'expected_verdict': 'SUPPORTED', 'anchor': {'all_of': ['entire year', 'forest']}, 'verdict': 'confirmed', 'claim_id': 'B-004', 'model_verdict': 'SUPPORTED', 'evidence': 'The old lady nevertheless lived the entire year in the forest.', 'rationale': 'The text directly says the old lady lived there all year.'}, {'probe_id': 'b-herbs-not-magic', 'expected_verdict': 'SUPPORTED', 'anchor': {'all_of': ['wild herbs']}, 'verdict': 'confirmed', 'claim_id': 'B-008', 'model_verdict': 'SUPPORTED', 'evidence': 'she collected wild herbs.', 'rationale': 'The passage states she collected wild herbs.'}, {'probe_id': 'b-herbs-to-cities', 'expected_verdict': 'SUPPORTED', 'anchor': {'all_of': ['travelers', 'cities']}, 'verdict': 'confirmed', 'claim_id': 'B-012', 'model_verdict': 'SUPPORTED', 'evidence': 'who transported and sold the herbs into the cities.', 'rationale': 'The passage states wandering travelers transported and sold the herbs into the cities.'}, {'probe_id': 'b-medicine-village', 'expected_verdict': 'SUPPORTED', 'anchor': {'all_of': ['medicine', 'entire village']}, 'verdict': 'confirmed', 'claim_id': 'B-026', 'model_verdict': 'SUPPORTED', 'evidence': 'she gave him enough medicine made out of the herbs she collected in spring for the entire village', 'rationale': 'The reference text states she gave him enough medicine made out of the herbs she collected in spring for the entire village.'}, {'probe_id': 'b-cheese-autumn', 'expected_verdict': 'SUPPORTED', 'anchor': {'all_of': ['cheese', 'autumn']}, 'verdict': 'confirmed', 'claim_id': 'B-031', 'model_verdict': 'SUPPORTED', 'evidence': 'the villagers sent a huge basket of cheese and dried meat to the old lady every autumn', 'rationale': 'The reference text states that the villagers sent a huge basket of cheese and dried meat to the old lady every autumn.'}, {'probe_id': 'm-visit-hours-reverse', 'expected_verdict': 'SUPPORTED', 'anchor': {'all_of': ['8 am', '9 am']}, 'verdict': 'confirmed', 'claim_id': 'M-002', 'model_verdict': 'SUPPORTED', 'evidence': 'She visited the rosegarden every morning from 8 am till 9 am, right after her breakfast.', 'rationale': 'The source states this visitation schedule verbatim.'}, {'probe_id': 'm-forest-season-reverse', 'expected_verdict': 'SUPPORTED', 'anchor': {'all_of': ['not traversable', 'October', 'April']}, 'verdict': 'confirmed', 'claim_id': 'M-016', 'model_verdict': 'SUPPORTED', 'evidence': 'The forest was not traversable between October and April as it was too wet and too cold', 'rationale': 'The source states this traversability condition verbatim.'}, {'probe_id': 'm-public-dates-reverse', 'expected_verdict': 'SUPPORTED', 'anchor': {'all_of': ['March 1st', 'March 31st']}, 'verdict': 'confirmed', 'claim_id': 'M-010', 'model_verdict': 'SUPPORTED', 'evidence': 'the young princess insisted on making the garden accessible to the public every March 1st to March 31st', 'rationale': 'The source contains this exact statement of public access dates.'}, {'probe_id': 'm-herbs-reverse', 'expected_verdict': 'SUPPORTED', 'anchor': {'all_of': ['wild herbs']}, 'verdict': 'confirmed', 'claim_id': 'M-020', 'model_verdict': 'SUPPORTED', 'evidence': 'she collected wild herbs', 'rationale': 'The source states this activity exactly.'}]}, {'fixture': 'dropped_claim', 'guard_probes_total': 9, 'guard_confirmed': ['b-log-rotation', 'b-worker-per-core', 'b-health-path', 'b-health-failures', 'a-listen-port', 'a-connect-timeout', 'a-read-timeout', 'a-retry-limit', 'a-tls-version'], 'guard_wrong': [], 'guard_unclaimed': [], 'detail': [{'probe_id': 'b-log-rotation', 'expected_verdict': 'SUPPORTED', 'anchor': {'all_of': ['rotated', '100 megabytes']}, 'verdict': 'confirmed', 'claim_id': 'B-007', 'model_verdict': 'SUPPORTED', 'evidence': 'Log files are rotated when they reach 100 megabytes.', 'rationale': 'The reference text contains the exact sentence stating the claim.'}, {'probe_id': 'b-worker-per-core', 'expected_verdict': 'SUPPORTED', 'anchor': {'all_of': ['one worker process', 'per CPU core']}, 'verdict': 'confirmed', 'claim_id': 'B-005', 'model_verdict': 'SUPPORTED', 'evidence': 'One worker process is started per CPU core.', 'rationale': 'The reference text contains the exact sentence stating the claim.'}, {'probe_id': 'b-health-path', 'expected_verdict': 'SUPPORTED', 'anchor': {'all_of': ['health check endpoint', '/-/healthy']}, 'verdict': 'confirmed', 'claim_id': 'B-008', 'model_verdict': 'SUPPORTED', 'evidence': 'The health check endpoint is /-/healthy.', 'rationale': 'The reference text contains the exact sentence stating the claim.'}, {'probe_id': 'b-health-failures', 'expected_verdict': 'SUPPORTED', 'anchor': {'all_of': ['3 consecutive health checks', 'removed from rotation']}, 'verdict': 'confirmed', 'claim_id': 'B-009', 'model_verdict': 'SUPPORTED', 'evidence': 'A relay that fails 3 consecutive health checks is removed from rotation.', 'rationale': 'The reference text contains the exact sentence stating the claim.'}, {'probe_id': 'a-listen-port', 'expected_verdict': 'SUPPORTED', 'anchor': {'all_of': ['listens on port', '8443']}, 'verdict': 'confirmed', 'claim_id': 'A-003', 'model_verdict': 'SUPPORTED', 'evidence': 'The relay listens on port 8443 by default.', 'rationale': 'The reference text contains the exact sentence stating the claim.'}, {'probe_id': 'a-connect-timeout', 'expected_verdict': 'SUPPORTED', 'anchor': {'all_of': ['connect timeout', '30 seconds']}, 'verdict': 'confirmed', 'claim_id': 'A-004', 'model_verdict': 'SUPPORTED', 'evidence': 'The default connect timeout is 30 seconds.', 'rationale': 'The reference text contains the exact sentence stating the claim.'}, {'probe_id': 'a-read-timeout', 'expected_verdict': 'SUPPORTED', 'anchor': {'all_of': ['read timeout', '120 seconds']}, 'verdict': 'confirmed', 'claim_id': 'A-005', 'model_verdict': 'SUPPORTED', 'evidence': 'The default read timeout is 120 seconds.', 'rationale': 'The reference text contains the exact sentence stating the claim.'}, {'probe_id': 'a-retry-limit', 'expected_verdict': 'SUPPORTED', 'anchor': {'all_of': ['retried at most 3 times']}, 'verdict': 'confirmed', 'claim_id': 'A-007', 'model_verdict': 'SUPPORTED', 'evidence': 'Each request is retried at most 3 times.', 'rationale': 'The reference text contains the exact sentence stating the claim.'}, {'probe_id': 'a-tls-version', 'expected_verdict': 'SUPPORTED', 'anchor': {'all_of': ['minimum accepted TLS version', 'TLS 1.2']}, 'verdict': 'confirmed', 'claim_id': 'A-008', 'model_verdict': 'SUPPORTED', 'evidence': 'The minimum accepted TLS version is TLS 1.2.', 'rationale': 'The reference text contains the exact sentence stating the claim.'}]}, {'fixture': 'hallucination', 'guard_probes_total': 7, 'guard_confirmed': ['merged-connect-timeout', 'merged-max-connections', 'a-listen-port', 'a-tls-version', 'a-client-cert', 'b-log-format', 'b-health-path'], 'guard_wrong': [], 'guard_unclaimed': [], 'detail': [{'probe_id': 'merged-connect-timeout', 'expected_verdict': 'SUPPORTED', 'anchor': {'all_of': ['connect timeout', '30 seconds']}, 'verdict': 'confirmed', 'claim_id': 'M-004', 'model_verdict': 'SUPPORTED', 'evidence': 'The default connect timeout is 30 seconds.', 'rationale': 'The exact wording is found in source_a.md.'}, {'probe_id': 'merged-max-connections', 'expected_verdict': 'SUPPORTED', 'anchor': {'all_of': ['512 concurrent connections']}, 'verdict': 'confirmed', 'claim_id': 'M-005', 'model_verdict': 'SUPPORTED', 'evidence': 'The relay accepts at most 512 concurrent connections.', 'rationale': 'The sentence is stated verbatim in source_b.md.'}, {'probe_id': 'a-listen-port', 'expected_verdict': 'SUPPORTED', 'anchor': {'all_of': ['listens on port', '8443']}, 'verdict': 'confirmed', 'claim_id': 'A-003', 'model_verdict': 'SUPPORTED', 'evidence': 'The relay listens on port 8443 by default.', 'rationale': "The claim conveys the same meaning as the reference text, which states 'The relay listens on port 8443 by default'."}, {'probe_id': 'a-tls-version', 'expected_verdict': 'SUPPORTED', 'anchor': {'all_of': ['minimum accepted TLS version', 'TLS 1.2']}, 'verdict': 'confirmed', 'claim_id': 'A-005', 'model_verdict': 'SUPPORTED', 'evidence': 'The minimum accepted TLS version is TLS 1.2.', 'rationale': 'The claim matches the reference text verbatim.'}, {'probe_id': 'a-client-cert', 'expected_verdict': 'SUPPORTED', 'anchor': {'all_of': ['client certificate verification', 'disabled by default']}, 'verdict': 'confirmed', 'claim_id': 'A-006', 'model_verdict': 'SUPPORTED', 'evidence': 'Client certificate verification is disabled by default.', 'rationale': 'The claim matches the reference text verbatim.'}, {'probe_id': 'b-log-format', 'expected_verdict': 'SUPPORTED', 'anchor': {'all_of': ['access logs', 'JSON Lines']}, 'verdict': 'confirmed', 'claim_id': 'B-004', 'model_verdict': 'SUPPORTED', 'evidence': 'Access logs are written in JSON Lines format.', 'rationale': 'The claim matches the reference text verbatim.'}, {'probe_id': 'b-health-path', 'expected_verdict': 'SUPPORTED', 'anchor': {'all_of': ['health check endpoint', '/-/healthy']}, 'verdict': 'confirmed', 'claim_id': 'B-005', 'model_verdict': 'SUPPORTED', 'evidence': 'The health check endpoint is /-/healthy.', 'rationale': 'The claim matches the reference text verbatim.'}]}, {'fixture': 'numeric_drift', 'guard_probes_total': 4, 'guard_confirmed': ['b-max-connections', 'a-retry-limit', 'b-listen-port', 'a-connect-timeout'], 'guard_wrong': [], 'guard_unclaimed': [], 'detail': [{'probe_id': 'b-max-connections', 'expected_verdict': 'SUPPORTED', 'anchor': {'all_of': ['connection limit', '512']}, 'verdict': 'confirmed', 'claim_id': 'B-003', 'model_verdict': 'SUPPORTED', 'evidence': 'The connection limit is 512 concurrent connections.', 'rationale': 'The claim exactly matches the reference sentence.'}, {'probe_id': 'a-retry-limit', 'expected_verdict': 'SUPPORTED', 'anchor': {'all_of': ['retried at most 3 times']}, 'verdict': 'confirmed', 'claim_id': 'A-004', 'model_verdict': 'SUPPORTED', 'evidence': 'Each request is retried at most 3 times.', 'rationale': 'The claim matches the sentence in the reference text.'}, {'probe_id': 'b-listen-port', 'expected_verdict': 'SUPPORTED', 'anchor': {'all_of': ['listens on port', '8443']}, 'verdict': 'confirmed', 'claim_id': 'A-005', 'model_verdict': 'SUPPORTED', 'evidence': 'The relay listens on port 8443 by default.', 'rationale': 'The claim is exactly stated in the reference.'}, {'probe_id': 'a-connect-timeout', 'expected_verdict': 'SUPPORTED', 'anchor': {'all_of': ['connect timeout', '30 seconds']}, 'verdict': 'confirmed', 'claim_id': 'A-006', 'model_verdict': 'SUPPORTED', 'evidence': 'The default connect timeout is 30 seconds.', 'rationale': 'The claim matches the reference sentence verbatim.'}]}, {'fixture': 'ordering_only', 'guard_probes_total': 11, 'guard_confirmed': ['a-listen-port', 'a-connect-timeout', 'a-max-connections', 'a-tls-version', 'b-log-format', 'b-retry-limit', 'b-read-timeout', 'b-health-path', 'merged-backoff', 'merged-client-cert', 'merged-log-rotation'], 'guard_wrong': [], 'guard_unclaimed': [], 'detail': [{'probe_id': 'a-listen-port', 'expected_verdict': 'SUPPORTED', 'anchor': {'all_of': ['listens on port', '8443']}, 'verdict': 'confirmed', 'claim_id': 'A-003', 'model_verdict': 'SUPPORTED', 'evidence': 'The relay listens on port 8443 by default.', 'rationale': 'The reference text contains the exact statement.'}, {'probe_id': 'a-connect-timeout', 'expected_verdict': 'SUPPORTED', 'anchor': {'all_of': ['connect timeout', '30 seconds']}, 'verdict': 'confirmed', 'claim_id': 'A-004', 'model_verdict': 'SUPPORTED', 'evidence': 'The default connect timeout is 30 seconds.', 'rationale': 'The reference text contains the exact statement.'}, {'probe_id': 'a-max-connections', 'expected_verdict': 'SUPPORTED', 'anchor': {'all_of': ['512 concurrent connections']}, 'verdict': 'confirmed', 'claim_id': 'A-007', 'model_verdict': 'SUPPORTED', 'evidence': 'The relay accepts at most 512 concurrent connections.', 'rationale': 'The reference text contains the exact statement.'}, {'probe_id': 'a-tls-version', 'expected_verdict': 'SUPPORTED', 'anchor': {'all_of': ['minimum accepted TLS version', 'TLS 1.2']}, 'verdict': 'confirmed', 'claim_id': 'A-011', 'model_verdict': 'SUPPORTED', 'evidence': 'The minimum accepted TLS version is TLS 1.2.', 'rationale': 'The reference text contains the exact statement.'}, {'probe_id': 'b-log-format', 'expected_verdict': 'SUPPORTED', 'anchor': {'all_of': ['access logs', 'JSON Lines']}, 'verdict': 'confirmed', 'claim_id': 'A-013', 'model_verdict': 'SUPPORTED', 'evidence': 'Access logs are written in JSON Lines format.', 'rationale': 'The reference text contains the exact statement.'}, {'probe_id': 'b-retry-limit', 'expected_verdict': 'SUPPORTED', 'anchor': {'all_of': ['retried at most 3 times']}, 'verdict': 'confirmed', 'claim_id': 'A-008', 'model_verdict': 'SUPPORTED', 'evidence': 'Each request is retried at most 3 times.', 'rationale': 'The reference text contains the exact statement.'}, {'probe_id': 'b-read-timeout', 'expected_verdict': 'SUPPORTED', 'anchor': {'all_of': ['read timeout', '120 seconds']}, 'verdict': 'confirmed', 'claim_id': 'A-005', 'model_verdict': 'SUPPORTED', 'evidence': 'The default read timeout is 120 seconds.', 'rationale': 'The reference text contains the exact statement.'}, {'probe_id': 'b-health-path', 'expected_verdict': 'SUPPORTED', 'anchor': {'all_of': ['health check endpoint', '/-/healthy']}, 'verdict': 'confirmed', 'claim_id': 'A-006', 'model_verdict': 'SUPPORTED', 'evidence': 'The health check endpoint is /-/healthy.', 'rationale': 'The reference text contains the exact statement.'}, {'probe_id': 'merged-backoff', 'expected_verdict': 'SUPPORTED', 'anchor': {'all_of': ['250 milliseconds', 'doubles']}, 'verdict': 'confirmed', 'claim_id': 'M-004', 'model_verdict': 'SUPPORTED', 'evidence': 'Retry backoff starts at 250 milliseconds and doubles on every attempt.', 'rationale': 'The exact wording is found in the Limits sections of both documents.'}, {'probe_id': 'merged-client-cert', 'expected_verdict': 'SUPPORTED', 'anchor': {'all_of': ['client certificate verification', 'disabled by default']}, 'verdict': 'confirmed', 'claim_id': 'M-012', 'model_verdict': 'SUPPORTED', 'evidence': 'Client certificate verification is disabled by default.', 'rationale': 'The sentence appears unchanged in the Security sections of both files.'}, {'probe_id': 'merged-log-rotation', 'expected_verdict': 'SUPPORTED', 'anchor': {'all_of': ['rotated', '100 megabytes']}, 'verdict': 'confirmed', 'claim_id': 'M-006', 'model_verdict': 'SUPPORTED', 'evidence': 'Log files are rotated when they reach 100 megabytes.', 'rationale': 'The sentence appears exactly in the Logging sections of both sources.'}]}, {'fixture': 'paraphrase', 'guard_probes_total': 11, 'guard_confirmed': ['a-listen-port', 'a-connect-timeout', 'a-read-timeout', 'a-max-connections', 'a-tls-version', 'b-retry-limit', 'b-log-format', 'b-health-path', 'merged-listen-port', 'merged-tls-version', 'merged-read-timeout'], 'guard_wrong': [], 'guard_unclaimed': [], 'detail': [{'probe_id': 'a-listen-port', 'expected_verdict': 'SUPPORTED', 'anchor': {'all_of': ['binds to port', '8443']}, 'verdict': 'confirmed', 'claim_id': 'A-002', 'model_verdict': 'SUPPORTED', 'evidence': "The relay's default listen port is 8443.", 'rationale': 'The reference states the relay listens on port 8443, matching the claim.'}, {'probe_id': 'a-connect-timeout', 'expected_verdict': 'SUPPORTED', 'anchor': {'all_of': ['30 seconds', 'abandoned']}, 'verdict': 'confirmed', 'claim_id': 'A-003', 'model_verdict': 'SUPPORTED', 'evidence': 'Connections that do not complete a handshake inside 30 seconds are dropped.', 'rationale': 'The reference describes a 30‑second handshake timeout and dropping the connection, which aligns with the claim.'}, {'probe_id': 'a-read-timeout', 'expected_verdict': 'SUPPORTED', 'anchor': {'all_of': ['waits up to 120 seconds']}, 'verdict': 'confirmed', 'claim_id': 'A-004', 'model_verdict': 'SUPPORTED', 'evidence': 'The relay allows a response 120 seconds to arrive before it times out.', 'rationale': 'The reference indicates a 120‑second wait for a response, matching the claim.'}, {'probe_id': 'a-max-connections', 'expected_verdict': 'SUPPORTED', 'anchor': {'all_of': ['512 connections', 'same time']}, 'verdict': 'confirmed', 'claim_id': 'A-005', 'model_verdict': 'SUPPORTED', 'evidence': 'At most 512 connections may be held open concurrently.', 'rationale': 'The reference limits concurrent connections to 512, which supports the claim.'}, {'probe_id': 'a-tls-version', 'expected_verdict': 'SUPPORTED', 'anchor': {'all_of': ['older than TLS 1.2', 'rejected']}, 'verdict': 'confirmed', 'claim_id': 'A-007', 'model_verdict': 'SUPPORTED', 'evidence': 'The relay will not negotiate any protocol version below TLS 1.2.', 'rationale': 'The reference says versions below TLS\u202f1.2 are not negotiated, which fulfills the claim.'}, {'probe_id': 'b-retry-limit', 'expected_verdict': 'SUPPORTED', 'anchor': {'all_of': ['Retry attempts per request', '3']}, 'verdict': 'confirmed', 'claim_id': 'B-006', 'model_verdict': 'SUPPORTED', 'evidence': 'The relay makes up to 3 further attempts after a request fails.', 'rationale': 'The reference allows up to three retry attempts per request, which fulfills the claim.'}, {'probe_id': 'b-log-format', 'expected_verdict': 'SUPPORTED', 'anchor': {'all_of': ['JSON Lines', 'one record per line']}, 'verdict': 'confirmed', 'claim_id': 'B-008', 'model_verdict': 'SUPPORTED', 'evidence': 'Access log records are emitted one JSON object per line.', 'rationale': 'The reference describes the access log as JSON Lines with one record per line.'}, {'probe_id': 'b-health-path', 'expected_verdict': 'SUPPORTED', 'anchor': {'all_of': ['health checks', '/-/healthy']}, 'verdict': 'confirmed', 'claim_id': 'B-009', 'model_verdict': 'SUPPORTED', 'evidence': 'The endpoint /-/healthy answers health checks.', 'rationale': 'The reference states health checks are served at /-/healthy, satisfying the claim.'}, {'probe_id': 'merged-listen-port', 'expected_verdict': 'SUPPORTED', 'anchor': {'all_of': ['default listen port', '8443']}, 'verdict': 'confirmed', 'claim_id': 'M-003', 'model_verdict': 'SUPPORTED', 'evidence': '- Listen port: 8443', 'rationale': 'Source_b lists the listen port as 8443, which is the default port claimed.'}, {'probe_id': 'merged-tls-version', 'expected_verdict': 'SUPPORTED', 'anchor': {'all_of': ['will not negotiate', 'TLS 1.2']}, 'verdict': 'confirmed', 'claim_id': 'M-008', 'model_verdict': 'SUPPORTED', 'evidence': 'Anything older than TLS 1.2 is rejected.', 'rationale': 'The source says protocols older than TLS\u202f1.2 are rejected, which means none below are negotiated.'}, {'probe_id': 'merged-read-timeout', 'expected_verdict': 'SUPPORTED', 'anchor': {'all_of': ['120 seconds', 'times out']}, 'verdict': 'confirmed', 'claim_id': 'M-005', 'model_verdict': 'SUPPORTED', 'evidence': 'Once connected, the relay waits up to 120 seconds for a response.', 'rationale': 'Source_a explicitly says the relay waits up to 120\u202fseconds for a response.'}]}, {'fixture': 'structure_added', 'guard_probes_total': 8, 'guard_confirmed': ['a-listen-port', 'a-read-timeout', 'b-health-path', 'b-tls-min', 'm-listen-port', 'm-read-timeout', 'm-health-path', 'm-tls-min'], 'guard_wrong': [], 'guard_unclaimed': [], 'detail': [{'probe_id': 'a-listen-port', 'expected_verdict': 'SUPPORTED', 'anchor': {'all_of': ['listens on port 8443']}, 'verdict': 'confirmed', 'claim_id': 'A-001', 'model_verdict': 'SUPPORTED', 'evidence': 'The relay listens on port 8443.', 'rationale': 'The reference text contains the exact sentence stating the relay listens on port 8443.'}, {'probe_id': 'a-read-timeout', 'expected_verdict': 'SUPPORTED', 'anchor': {'all_of': ['read timeout is 30 seconds']}, 'verdict': 'confirmed', 'claim_id': 'A-002', 'model_verdict': 'SUPPORTED', 'evidence': 'The read timeout is 30 seconds.', 'rationale': 'The reference text explicitly states the read timeout is 30 seconds.'}, {'probe_id': 'b-health-path', 'expected_verdict': 'SUPPORTED', 'anchor': {'all_of': ['/healthz']}, 'verdict': 'confirmed', 'claim_id': 'B-002', 'model_verdict': 'SUPPORTED', 'evidence': 'The health check path is /healthz.', 'rationale': 'The reference text explicitly states the health check path is /healthz.'}, {'probe_id': 'b-tls-min', 'expected_verdict': 'SUPPORTED', 'anchor': {'all_of': ['minimum TLS version is 1.2']}, 'verdict': 'confirmed', 'claim_id': 'B-003', 'model_verdict': 'SUPPORTED', 'evidence': 'The minimum TLS version is 1.2.', 'rationale': 'The reference text states the minimum TLS version is 1.2.'}, {'probe_id': 'm-listen-port', 'expected_verdict': 'SUPPORTED', 'anchor': {'all_of': ['listens on port 8443']}, 'verdict': 'confirmed', 'claim_id': 'M-001', 'model_verdict': 'SUPPORTED', 'evidence': 'The relay listens on port 8443.', 'rationale': 'The exact sentence appears in source_a.md, directly supporting the claim.'}, {'probe_id': 'm-read-timeout', 'expected_verdict': 'SUPPORTED', 'anchor': {'all_of': ['read timeout is 30 seconds']}, 'verdict': 'confirmed', 'claim_id': 'M-002', 'model_verdict': 'SUPPORTED', 'evidence': 'The read timeout is 30 seconds.', 'rationale': 'The sentence is present verbatim in source_a.md, so the claim is supported.'}, {'probe_id': 'm-health-path', 'expected_verdict': 'SUPPORTED', 'anchor': {'all_of': ['/healthz']}, 'verdict': 'confirmed', 'claim_id': 'M-003', 'model_verdict': 'SUPPORTED', 'evidence': 'The health check path is /healthz.', 'rationale': 'Source_b.md contains this exact statement, providing direct support.'}, {'probe_id': 'm-tls-min', 'expected_verdict': 'SUPPORTED', 'anchor': {'all_of': ['minimum TLS version is 1.2']}, 'verdict': 'confirmed', 'claim_id': 'M-004', 'model_verdict': 'SUPPORTED', 'evidence': 'The minimum TLS version is 1.2.', 'rationale': 'The claim matches the exact wording in source_b.md, so it is supported.'}]}]}
\expandafter\def\csname ccval@regrade.guard-probes.amended\endcsname{2026-09-03}
\expandafter\def\csname ccval@regrade.inventions.A-27b.fixtures\endcsname{13}
\expandafter\def\csname ccval@regrade.inventions.A-27b.fixtures_measured\endcsname{13}
\expandafter\def\csname ccval@regrade.inventions.A-27b.invented_total\endcsname{1}
\expandafter\def\csname ccval@regrade.inventions.A-27b.plants_detected\endcsname{7}
\expandafter\def\csname ccval@regrade.inventions.A-27b.plants_missed\endcsname{0}
\expandafter\def\csname ccval@regrade.inventions.A-27b.plants_total\endcsname{7}
\expandafter\def\csname ccval@regrade.inventions.A-27b.plants_unmeasured\endcsname{0}
\expandafter\def\csname ccval@regrade.inventions.B-70b.fixtures\endcsname{13}
\expandafter\def\csname ccval@regrade.inventions.B-70b.fixtures_measured\endcsname{12}
\expandafter\def\csname ccval@regrade.inventions.B-70b.invented_total\endcsname{0}
\expandafter\def\csname ccval@regrade.inventions.B-70b.plants_detected\endcsname{6}
\expandafter\def\csname ccval@regrade.inventions.B-70b.plants_missed\endcsname{0}
\expandafter\def\csname ccval@regrade.inventions.B-70b.plants_total\endcsname{7}
\expandafter\def\csname ccval@regrade.inventions.B-70b.plants_unmeasured\endcsname{1}
\expandafter\def\csname ccval@regrade.inventions.C-120b.fixtures\endcsname{13}
\expandafter\def\csname ccval@regrade.inventions.C-120b.fixtures_measured\endcsname{13}
\expandafter\def\csname ccval@regrade.inventions.C-120b.invented_total\endcsname{2}
\expandafter\def\csname ccval@regrade.inventions.C-120b.plants_detected\endcsname{6}
\expandafter\def\csname ccval@regrade.inventions.C-120b.plants_missed\endcsname{1}
\expandafter\def\csname ccval@regrade.inventions.C-120b.plants_total\endcsname{7}
\expandafter\def\csname ccval@regrade.inventions.C-120b.plants_unmeasured\endcsname{0}
\expandafter\def\csname ccval@regrade.inventions.amended\endcsname{2026-09-03}
\expandafter\def\csname ccval@regrade.pair429.A1-27b-default.a_segments\endcsname{33}
\expandafter\def\csname ccval@regrade.pair429.A1-27b-default.a_undeclared_absence\endcsname{0}
\expandafter\def\csname ccval@regrade.pair429.A1-27b-default.a_undeclared_rewording\endcsname{0}
\expandafter\def\csname ccval@regrade.pair429.A1-27b-default.b_segments\endcsname{30}
\expandafter\def\csname ccval@regrade.pair429.A1-27b-default.b_undeclared_absence\endcsname{0}
\expandafter\def\csname ccval@regrade.pair429.A1-27b-default.b_undeclared_rewording\endcsname{0}
\expandafter\def\csname ccval@regrade.pair429.A1-27b-default.check2_findings\endcsname{0}
\expandafter\def\csname ccval@regrade.pair429.A1-27b-default.fidelity\endcsname{off}
\expandafter\def\csname ccval@regrade.pair429.A1-27b-default.merged_bytes\endcsname{5673}
\expandafter\def\csname ccval@regrade.pair429.A1-27b-default.near_match\endcsname{0.68}
\expandafter\def\csname ccval@regrade.pair429.A1-27b-default.recorded_undeclared_absence\endcsname{0}
\expandafter\def\csname ccval@regrade.pair429.A1-27b-default.segments\endcsname{63}
\expandafter\def\csname ccval@regrade.pair429.A1-27b-default.undeclared_absence\endcsname{0}
\expandafter\def\csname ccval@regrade.pair429.A1-27b-default.undeclared_rewording\endcsname{0}
\expandafter\def\csname ccval@regrade.pair429.A2-27b-think.a_segments\endcsname{33}
\expandafter\def\csname ccval@regrade.pair429.A2-27b-think.a_undeclared_absence\endcsname{0}
\expandafter\def\csname ccval@regrade.pair429.A2-27b-think.a_undeclared_rewording\endcsname{0}
\expandafter\def\csname ccval@regrade.pair429.A2-27b-think.b_segments\endcsname{30}
\expandafter\def\csname ccval@regrade.pair429.A2-27b-think.b_undeclared_absence\endcsname{0}
\expandafter\def\csname ccval@regrade.pair429.A2-27b-think.b_undeclared_rewording\endcsname{0}
\expandafter\def\csname ccval@regrade.pair429.A2-27b-think.check2_findings\endcsname{0}
\expandafter\def\csname ccval@regrade.pair429.A2-27b-think.fidelity\endcsname{off}
\expandafter\def\csname ccval@regrade.pair429.A2-27b-think.merged_bytes\endcsname{5673}
\expandafter\def\csname ccval@regrade.pair429.A2-27b-think.near_match\endcsname{0.68}
\expandafter\def\csname ccval@regrade.pair429.A2-27b-think.recorded_undeclared_absence\endcsname{0}
\expandafter\def\csname ccval@regrade.pair429.A2-27b-think.segments\endcsname{63}
\expandafter\def\csname ccval@regrade.pair429.A2-27b-think.undeclared_absence\endcsname{0}
\expandafter\def\csname ccval@regrade.pair429.A2-27b-think.undeclared_rewording\endcsname{0}
\expandafter\def\csname ccval@regrade.pair429.B1-70b-default.a_segments\endcsname{33}
\expandafter\def\csname ccval@regrade.pair429.B1-70b-default.a_undeclared_absence\endcsname{2}
\expandafter\def\csname ccval@regrade.pair429.B1-70b-default.a_undeclared_rewording\endcsname{0}
\expandafter\def\csname ccval@regrade.pair429.B1-70b-default.b_segments\endcsname{30}
\expandafter\def\csname ccval@regrade.pair429.B1-70b-default.b_undeclared_absence\endcsname{16}
\expandafter\def\csname ccval@regrade.pair429.B1-70b-default.b_undeclared_rewording\endcsname{6}
\expandafter\def\csname ccval@regrade.pair429.B1-70b-default.check2_findings\endcsname{24}
\expandafter\def\csname ccval@regrade.pair429.B1-70b-default.fidelity\endcsname{off}
\expandafter\def\csname ccval@regrade.pair429.B1-70b-default.merged_bytes\endcsname{4234}
\expandafter\def\csname ccval@regrade.pair429.B1-70b-default.near_match\endcsname{0.68}
\expandafter\def\csname ccval@regrade.pair429.B1-70b-default.recorded_undeclared_absence\endcsname{24}
\expandafter\def\csname ccval@regrade.pair429.B1-70b-default.segments\endcsname{63}
\expandafter\def\csname ccval@regrade.pair429.B1-70b-default.undeclared_absence\endcsname{18}
\expandafter\def\csname ccval@regrade.pair429.B1-70b-default.undeclared_rewording\endcsname{6}
\expandafter\def\csname ccval@regrade.pair429.B2-70b-think.a_segments\endcsname{33}
\expandafter\def\csname ccval@regrade.pair429.B2-70b-think.a_undeclared_absence\endcsname{2}
\expandafter\def\csname ccval@regrade.pair429.B2-70b-think.a_undeclared_rewording\endcsname{0}
\expandafter\def\csname ccval@regrade.pair429.B2-70b-think.b_segments\endcsname{30}
\expandafter\def\csname ccval@regrade.pair429.B2-70b-think.b_undeclared_absence\endcsname{16}
\expandafter\def\csname ccval@regrade.pair429.B2-70b-think.b_undeclared_rewording\endcsname{6}
\expandafter\def\csname ccval@regrade.pair429.B2-70b-think.check2_findings\endcsname{24}
\expandafter\def\csname ccval@regrade.pair429.B2-70b-think.fidelity\endcsname{off}
\expandafter\def\csname ccval@regrade.pair429.B2-70b-think.merged_bytes\endcsname{4234}
\expandafter\def\csname ccval@regrade.pair429.B2-70b-think.near_match\endcsname{0.68}
\expandafter\def\csname ccval@regrade.pair429.B2-70b-think.recorded_undeclared_absence\endcsname{24}
\expandafter\def\csname ccval@regrade.pair429.B2-70b-think.segments\endcsname{63}
\expandafter\def\csname ccval@regrade.pair429.B2-70b-think.undeclared_absence\endcsname{18}
\expandafter\def\csname ccval@regrade.pair429.B2-70b-think.undeclared_rewording\endcsname{6}
\expandafter\def\csname ccval@regrade.pair429.C2-120b-think.a_segments\endcsname{33}
\expandafter\def\csname ccval@regrade.pair429.C2-120b-think.a_undeclared_absence\endcsname{0}
\expandafter\def\csname ccval@regrade.pair429.C2-120b-think.a_undeclared_rewording\endcsname{1}
\expandafter\def\csname ccval@regrade.pair429.C2-120b-think.b_segments\endcsname{30}
\expandafter\def\csname ccval@regrade.pair429.C2-120b-think.b_undeclared_absence\endcsname{2}
\expandafter\def\csname ccval@regrade.pair429.C2-120b-think.b_undeclared_rewording\endcsname{0}
\expandafter\def\csname ccval@regrade.pair429.C2-120b-think.check2_findings\endcsname{3}
\expandafter\def\csname ccval@regrade.pair429.C2-120b-think.fidelity\endcsname{off}
\expandafter\def\csname ccval@regrade.pair429.C2-120b-think.merged_bytes\endcsname{5692}
\expandafter\def\csname ccval@regrade.pair429.C2-120b-think.near_match\endcsname{0.68}
\expandafter\def\csname ccval@regrade.pair429.C2-120b-think.recorded_undeclared_absence\endcsname{3}
\expandafter\def\csname ccval@regrade.pair429.C2-120b-think.segments\endcsname{63}
\expandafter\def\csname ccval@regrade.pair429.C2-120b-think.undeclared_absence\endcsname{2}
\expandafter\def\csname ccval@regrade.pair429.C2-120b-think.undeclared_rewording\endcsname{1}
\expandafter\def\csname ccval@rerun.A-27b.attribution_invented.anchored\endcsname{11}
\expandafter\def\csname ccval@rerun.A-27b.attribution_invented.claims_total\endcsname{11}
\expandafter\def\csname ccval@rerun.A-27b.attribution_invented.graded\endcsname{11}
\expandafter\def\csname ccval@rerun.A-27b.attribution_invented.m-tls-attributed.claim_id\endcsname{M-004}
\expandafter\def\csname ccval@rerun.A-27b.attribution_invented.m-tls-attributed.verdict\endcsname{detected}
\expandafter\def\csname ccval@rerun.A-27b.attribution_swapped.anchored\endcsname{10}
\expandafter\def\csname ccval@rerun.A-27b.attribution_swapped.claims_total\endcsname{10}
\expandafter\def\csname ccval@rerun.A-27b.attribution_swapped.graded\endcsname{10}
\expandafter\def\csname ccval@rerun.A-27b.attribution_swapped.m-notes-timeout-attribution.claim_id\endcsname{M-003}
\expandafter\def\csname ccval@rerun.A-27b.attribution_swapped.m-notes-timeout-attribution.verdict\endcsname{detected}
\expandafter\def\csname ccval@rerun.A-27b.conflict_surfaced.anchored\endcsname{23}
\expandafter\def\csname ccval@rerun.A-27b.conflict_surfaced.claims_total\endcsname{23}
\expandafter\def\csname ccval@rerun.A-27b.conflict_surfaced.graded\endcsname{23}
\expandafter\def\csname ccval@rerun.A-27b.contradiction.a-connect-timeout.claim_id\endcsname{A-003}
\expandafter\def\csname ccval@rerun.A-27b.contradiction.a-connect-timeout.verdict\endcsname{detected}
\expandafter\def\csname ccval@rerun.A-27b.contradiction.anchored\endcsname{22}
\expandafter\def\csname ccval@rerun.A-27b.contradiction.claims_total\endcsname{22}
\expandafter\def\csname ccval@rerun.A-27b.contradiction.graded\endcsname{22}
\expandafter\def\csname ccval@rerun.A-27b.dedup.anchored\endcsname{10}
\expandafter\def\csname ccval@rerun.A-27b.dedup.claims_total\endcsname{10}
\expandafter\def\csname ccval@rerun.A-27b.dedup.graded\endcsname{10}
\expandafter\def\csname ccval@rerun.A-27b.disjoint_domains.anchored\endcsname{99}
\expandafter\def\csname ccval@rerun.A-27b.disjoint_domains.claims_total\endcsname{99}
\expandafter\def\csname ccval@rerun.A-27b.disjoint_domains.graded\endcsname{99}
\expandafter\def\csname ccval@rerun.A-27b.disjoint_sources.anchored\endcsname{74}
\expandafter\def\csname ccval@rerun.A-27b.disjoint_sources.claims_total\endcsname{74}
\expandafter\def\csname ccval@rerun.A-27b.disjoint_sources.graded\endcsname{74}
\expandafter\def\csname ccval@rerun.A-27b.disqualified_count\endcsname{1}
\expandafter\def\csname ccval@rerun.A-27b.dropped_claim.anchored\endcsname{26}
\expandafter\def\csname ccval@rerun.A-27b.dropped_claim.b-log-format.claim_id\endcsname{B-005}
\expandafter\def\csname ccval@rerun.A-27b.dropped_claim.b-log-format.verdict\endcsname{detected}
\expandafter\def\csname ccval@rerun.A-27b.dropped_claim.claims_total\endcsname{26}
\expandafter\def\csname ccval@rerun.A-27b.dropped_claim.graded\endcsname{25}
\expandafter\def\csname ccval@rerun.A-27b.exit_code_matched\endcsname{12}
\expandafter\def\csname ccval@rerun.A-27b.exit_code_missed\endcsname{1}
\expandafter\def\csname ccval@rerun.A-27b.exit_code_unmeasured\endcsname{0}
\expandafter\def\csname ccval@rerun.A-27b.fixtures\endcsname{13}
\expandafter\def\csname ccval@rerun.A-27b.fixtures_measured\endcsname{13}
\expandafter\def\csname ccval@rerun.A-27b.hallucination.anchored\endcsname{19}
\expandafter\def\csname ccval@rerun.A-27b.hallucination.claims_total\endcsname{19}
\expandafter\def\csname ccval@rerun.A-27b.hallucination.graded\endcsname{18}
\expandafter\def\csname ccval@rerun.A-27b.hallucination.merged-socks5.claim_id\endcsname{M-005}
\expandafter\def\csname ccval@rerun.A-27b.hallucination.merged-socks5.verdict\endcsname{detected}
\expandafter\def\csname ccval@rerun.A-27b.invented_on_unmeasured\endcsname{0}
\expandafter\def\csname ccval@rerun.A-27b.invented_total\endcsname{2}
\expandafter\def\csname ccval@rerun.A-27b.model\endcsname{qwen3.8:27b}
\expandafter\def\csname ccval@rerun.A-27b.numeric_drift.a-max-connections.claim_id\endcsname{A-002}
\expandafter\def\csname ccval@rerun.A-27b.numeric_drift.a-max-connections.verdict\endcsname{detected}
\expandafter\def\csname ccval@rerun.A-27b.numeric_drift.anchored\endcsname{13}
\expandafter\def\csname ccval@rerun.A-27b.numeric_drift.claims_total\endcsname{13}
\expandafter\def\csname ccval@rerun.A-27b.numeric_drift.graded\endcsname{13}
\expandafter\def\csname ccval@rerun.A-27b.ordering_only.anchored\endcsname{39}
\expandafter\def\csname ccval@rerun.A-27b.ordering_only.claims_total\endcsname{39}
\expandafter\def\csname ccval@rerun.A-27b.ordering_only.graded\endcsname{39}
\expandafter\def\csname ccval@rerun.A-27b.paraphrase.anchored\endcsname{28}
\expandafter\def\csname ccval@rerun.A-27b.paraphrase.claims_total\endcsname{28}
\expandafter\def\csname ccval@rerun.A-27b.paraphrase.graded\endcsname{28}
\expandafter\def\csname ccval@rerun.A-27b.plants_detected\endcsname{6}
\expandafter\def\csname ccval@rerun.A-27b.plants_missed\endcsname{0}
\expandafter\def\csname ccval@rerun.A-27b.plants_total\endcsname{6}
\expandafter\def\csname ccval@rerun.A-27b.plants_unmeasured\endcsname{0}
\expandafter\def\csname ccval@rerun.A-27b.structure_added.anchored\endcsname{9}
\expandafter\def\csname ccval@rerun.A-27b.structure_added.claims_total\endcsname{9}
\expandafter\def\csname ccval@rerun.A-27b.structure_added.graded\endcsname{9}
\expandafter\def\csname ccval@rerun.B-70b.attribution_invented.anchored\endcsname{11}
\expandafter\def\csname ccval@rerun.B-70b.attribution_invented.claims_total\endcsname{11}
\expandafter\def\csname ccval@rerun.B-70b.attribution_invented.graded\endcsname{11}
\expandafter\def\csname ccval@rerun.B-70b.attribution_invented.m-tls-attributed.claim_id\endcsname{M-004}
\expandafter\def\csname ccval@rerun.B-70b.attribution_invented.m-tls-attributed.verdict\endcsname{detected}
\expandafter\def\csname ccval@rerun.B-70b.attribution_swapped.anchored\endcsname{10}
\expandafter\def\csname ccval@rerun.B-70b.attribution_swapped.claims_total\endcsname{10}
\expandafter\def\csname ccval@rerun.B-70b.attribution_swapped.graded\endcsname{10}
\expandafter\def\csname ccval@rerun.B-70b.attribution_swapped.m-notes-timeout-attribution.claim_id\endcsname{M-003}
\expandafter\def\csname ccval@rerun.B-70b.attribution_swapped.m-notes-timeout-attribution.verdict\endcsname{detected}
\expandafter\def\csname ccval@rerun.B-70b.conflict_surfaced.anchored\endcsname{20}
\expandafter\def\csname ccval@rerun.B-70b.conflict_surfaced.claims_total\endcsname{20}
\expandafter\def\csname ccval@rerun.B-70b.conflict_surfaced.graded\endcsname{20}
\expandafter\def\csname ccval@rerun.B-70b.contradiction.a-connect-timeout.claim_id\endcsname{A-002}
\expandafter\def\csname ccval@rerun.B-70b.contradiction.a-connect-timeout.verdict\endcsname{detected}
\expandafter\def\csname ccval@rerun.B-70b.contradiction.anchored\endcsname{19}
\expandafter\def\csname ccval@rerun.B-70b.contradiction.claims_total\endcsname{19}
\expandafter\def\csname ccval@rerun.B-70b.contradiction.graded\endcsname{19}
\expandafter\def\csname ccval@rerun.B-70b.dedup.anchored\endcsname{10}
\expandafter\def\csname ccval@rerun.B-70b.dedup.claims_total\endcsname{10}
\expandafter\def\csname ccval@rerun.B-70b.dedup.graded\endcsname{10}
\expandafter\def\csname ccval@rerun.B-70b.disjoint_domains.anchored\endcsname{57}
\expandafter\def\csname ccval@rerun.B-70b.disjoint_domains.claims_total\endcsname{57}
\expandafter\def\csname ccval@rerun.B-70b.disjoint_domains.graded\endcsname{57}
\expandafter\def\csname ccval@rerun.B-70b.disjoint_sources.anchored\endcsname{58}
\expandafter\def\csname ccval@rerun.B-70b.disjoint_sources.claims_total\endcsname{58}
\expandafter\def\csname ccval@rerun.B-70b.disjoint_sources.graded\endcsname{58}
\expandafter\def\csname ccval@rerun.B-70b.disqualified_count\endcsname{0}
\expandafter\def\csname ccval@rerun.B-70b.dropped_claim.anchored\endcsname{23}
\expandafter\def\csname ccval@rerun.B-70b.dropped_claim.b-log-format.claim_id\endcsname{B-004}
\expandafter\def\csname ccval@rerun.B-70b.dropped_claim.b-log-format.verdict\endcsname{detected}
\expandafter\def\csname ccval@rerun.B-70b.dropped_claim.claims_total\endcsname{23}
\expandafter\def\csname ccval@rerun.B-70b.dropped_claim.graded\endcsname{22}
\expandafter\def\csname ccval@rerun.B-70b.exit_code_matched\endcsname{12}
\expandafter\def\csname ccval@rerun.B-70b.exit_code_missed\endcsname{0}
\expandafter\def\csname ccval@rerun.B-70b.exit_code_unmeasured\endcsname{1}
\expandafter\def\csname ccval@rerun.B-70b.fixtures\endcsname{13}
\expandafter\def\csname ccval@rerun.B-70b.fixtures_measured\endcsname{12}
\expandafter\def\csname ccval@rerun.B-70b.hallucination.anchored\endcsname{16}
\expandafter\def\csname ccval@rerun.B-70b.hallucination.claims_total\endcsname{16}
\expandafter\def\csname ccval@rerun.B-70b.hallucination.graded\endcsname{15}
\expandafter\def\csname ccval@rerun.B-70b.hallucination.merged-socks5.claim_id\endcsname{M-004}
\expandafter\def\csname ccval@rerun.B-70b.hallucination.merged-socks5.verdict\endcsname{detected}
\expandafter\def\csname ccval@rerun.B-70b.invented_on_unmeasured\endcsname{0}
\expandafter\def\csname ccval@rerun.B-70b.invented_total\endcsname{1}
\expandafter\def\csname ccval@rerun.B-70b.model\endcsname{deepseek-r1:70b}
\expandafter\def\csname ccval@rerun.B-70b.numeric_drift.a-max-connections.claim_id\endcsname{A-001}
\expandafter\def\csname ccval@rerun.B-70b.numeric_drift.a-max-connections.model_verdict\endcsname{SUPPORTED}
\expandafter\def\csname ccval@rerun.B-70b.numeric_drift.a-max-connections.verdict\endcsname{acceptable_no_finding}
\expandafter\def\csname ccval@rerun.B-70b.numeric_drift.anchored\endcsname{12}
\expandafter\def\csname ccval@rerun.B-70b.numeric_drift.claims_total\endcsname{12}
\expandafter\def\csname ccval@rerun.B-70b.numeric_drift.graded\endcsname{12}
\expandafter\def\csname ccval@rerun.B-70b.ordering_only.anchored\endcsname{33}
\expandafter\def\csname ccval@rerun.B-70b.ordering_only.claims_total\endcsname{33}
\expandafter\def\csname ccval@rerun.B-70b.ordering_only.graded\endcsname{33}
\expandafter\def\csname ccval@rerun.B-70b.paraphrase.anchored\endcsname{24}
\expandafter\def\csname ccval@rerun.B-70b.paraphrase.claims_total\endcsname{24}
\expandafter\def\csname ccval@rerun.B-70b.paraphrase.graded\endcsname{24}
\expandafter\def\csname ccval@rerun.B-70b.plants_detected\endcsname{5}
\expandafter\def\csname ccval@rerun.B-70b.plants_missed\endcsname{0}
\expandafter\def\csname ccval@rerun.B-70b.plants_total\endcsname{6}
\expandafter\def\csname ccval@rerun.B-70b.plants_unmeasured\endcsname{1}
\expandafter\def\csname ccval@rerun.B-70b.structure_added.anchored\endcsname{9}
\expandafter\def\csname ccval@rerun.B-70b.structure_added.claims_total\endcsname{9}
\expandafter\def\csname ccval@rerun.B-70b.structure_added.graded\endcsname{9}
\expandafter\def\csname ccval@rerun.C-120b.attribution_invented.anchored\endcsname{11}
\expandafter\def\csname ccval@rerun.C-120b.attribution_invented.claims_total\endcsname{11}
\expandafter\def\csname ccval@rerun.C-120b.attribution_invented.graded\endcsname{11}
\expandafter\def\csname ccval@rerun.C-120b.attribution_invented.m-tls-attributed.claim_id\endcsname{B-002}
\expandafter\def\csname ccval@rerun.C-120b.attribution_invented.m-tls-attributed.model_verdict\endcsname{SUPPORTED}
\expandafter\def\csname ccval@rerun.C-120b.attribution_invented.m-tls-attributed.verdict\endcsname{missed}
\expandafter\def\csname ccval@rerun.C-120b.attribution_swapped.anchored\endcsname{11}
\expandafter\def\csname ccval@rerun.C-120b.attribution_swapped.claims_total\endcsname{11}
\expandafter\def\csname ccval@rerun.C-120b.attribution_swapped.graded\endcsname{11}
\expandafter\def\csname ccval@rerun.C-120b.attribution_swapped.m-notes-timeout-attribution.claim_id\endcsname{M-004}
\expandafter\def\csname ccval@rerun.C-120b.attribution_swapped.m-notes-timeout-attribution.verdict\endcsname{detected}
\expandafter\def\csname ccval@rerun.C-120b.conflict_surfaced.anchored\endcsname{26}
\expandafter\def\csname ccval@rerun.C-120b.conflict_surfaced.claims_total\endcsname{26}
\expandafter\def\csname ccval@rerun.C-120b.conflict_surfaced.graded\endcsname{26}
\expandafter\def\csname ccval@rerun.C-120b.contradiction.a-connect-timeout.claim_id\endcsname{A-004}
\expandafter\def\csname ccval@rerun.C-120b.contradiction.a-connect-timeout.verdict\endcsname{detected}
\expandafter\def\csname ccval@rerun.C-120b.contradiction.anchored\endcsname{25}
\expandafter\def\csname ccval@rerun.C-120b.contradiction.claims_total\endcsname{25}
\expandafter\def\csname ccval@rerun.C-120b.contradiction.graded\endcsname{25}
\expandafter\def\csname ccval@rerun.C-120b.dedup.anchored\endcsname{10}
\expandafter\def\csname ccval@rerun.C-120b.dedup.claims_total\endcsname{10}
\expandafter\def\csname ccval@rerun.C-120b.dedup.graded\endcsname{10}
\expandafter\def\csname ccval@rerun.C-120b.disjoint_domains.anchored\endcsname{102}
\expandafter\def\csname ccval@rerun.C-120b.disjoint_domains.claims_total\endcsname{104}
\expandafter\def\csname ccval@rerun.C-120b.disjoint_domains.graded\endcsname{104}
\expandafter\def\csname ccval@rerun.C-120b.disjoint_sources.anchored\endcsname{75}
\expandafter\def\csname ccval@rerun.C-120b.disjoint_sources.claims_total\endcsname{77}
\expandafter\def\csname ccval@rerun.C-120b.disjoint_sources.graded\endcsname{77}
\expandafter\def\csname ccval@rerun.C-120b.disqualified_count\endcsname{2}
\expandafter\def\csname ccval@rerun.C-120b.dropped_claim.anchored\endcsname{29}
\expandafter\def\csname ccval@rerun.C-120b.dropped_claim.b-log-format.claim_id\endcsname{B-006}
\expandafter\def\csname ccval@rerun.C-120b.dropped_claim.b-log-format.verdict\endcsname{detected}
\expandafter\def\csname ccval@rerun.C-120b.dropped_claim.claims_total\endcsname{29}
\expandafter\def\csname ccval@rerun.C-120b.dropped_claim.graded\endcsname{28}
\expandafter\def\csname ccval@rerun.C-120b.exit_code_matched\endcsname{11}
\expandafter\def\csname ccval@rerun.C-120b.exit_code_missed\endcsname{2}
\expandafter\def\csname ccval@rerun.C-120b.exit_code_unmeasured\endcsname{0}
\expandafter\def\csname ccval@rerun.C-120b.fixtures\endcsname{13}
\expandafter\def\csname ccval@rerun.C-120b.fixtures_measured\endcsname{13}
\expandafter\def\csname ccval@rerun.C-120b.hallucination.anchored\endcsname{21}
\expandafter\def\csname ccval@rerun.C-120b.hallucination.claims_total\endcsname{21}
\expandafter\def\csname ccval@rerun.C-120b.hallucination.graded\endcsname{20}
\expandafter\def\csname ccval@rerun.C-120b.hallucination.merged-socks5.claim_id\endcsname{M-006}
\expandafter\def\csname ccval@rerun.C-120b.hallucination.merged-socks5.verdict\endcsname{detected}
\expandafter\def\csname ccval@rerun.C-120b.invented_on_unmeasured\endcsname{0}
\expandafter\def\csname ccval@rerun.C-120b.invented_total\endcsname{3}
\expandafter\def\csname ccval@rerun.C-120b.model\endcsname{gpt-oss:120b}
\expandafter\def\csname ccval@rerun.C-120b.numeric_drift.a-max-connections.claim_id\endcsname{A-003}
\expandafter\def\csname ccval@rerun.C-120b.numeric_drift.a-max-connections.verdict\endcsname{detected}
\expandafter\def\csname ccval@rerun.C-120b.numeric_drift.anchored\endcsname{18}
\expandafter\def\csname ccval@rerun.C-120b.numeric_drift.claims_total\endcsname{18}
\expandafter\def\csname ccval@rerun.C-120b.numeric_drift.graded\endcsname{18}
\expandafter\def\csname ccval@rerun.C-120b.ordering_only.anchored\endcsname{38}
\expandafter\def\csname ccval@rerun.C-120b.ordering_only.claims_total\endcsname{39}
\expandafter\def\csname ccval@rerun.C-120b.ordering_only.graded\endcsname{39}
\expandafter\def\csname ccval@rerun.C-120b.paraphrase.anchored\endcsname{28}
\expandafter\def\csname ccval@rerun.C-120b.paraphrase.claims_total\endcsname{28}
\expandafter\def\csname ccval@rerun.C-120b.paraphrase.graded\endcsname{28}
\expandafter\def\csname ccval@rerun.C-120b.plants_detected\endcsname{5}
\expandafter\def\csname ccval@rerun.C-120b.plants_missed\endcsname{1}
\expandafter\def\csname ccval@rerun.C-120b.plants_total\endcsname{6}
\expandafter\def\csname ccval@rerun.C-120b.plants_unmeasured\endcsname{0}
\expandafter\def\csname ccval@rerun.C-120b.structure_added.anchored\endcsname{9}
\expandafter\def\csname ccval@rerun.C-120b.structure_added.claims_total\endcsname{9}
\expandafter\def\csname ccval@rerun.C-120b.structure_added.graded\endcsname{9}
\expandafter\def\csname ccval@safemode.figures\endcsname{18}
\expandafter\def\csname ccval@safemode.haiku.rate_limits.silent.new\endcsname{1}
\expandafter\def\csname ccval@safemode.haiku.rate_limits.silent.old\endcsname{34}
\expandafter\def\csname ccval@safemode.moved\endcsname{1}
\expandafter\def\csname ccval@safemode.pairs\endcsname{3}
\expandafter\def\csname ccval@sweep.pair429.A1-27b-default.declared_drops\endcsname{0}
\expandafter\def\csname ccval@sweep.pair429.A1-27b-default.draws_graded\endcsname{3}
\expandafter\def\csname ccval@sweep.pair429.A1-27b-default.draws_identical\endcsname{3}
\expandafter\def\csname ccval@sweep.pair429.A1-27b-default.draws_k\endcsname{3}
\expandafter\def\csname ccval@sweep.pair429.A1-27b-default.exit_code\endcsname{0}
\expandafter\def\csname ccval@sweep.pair429.A1-27b-default.findings_grounded\endcsname{0}
\expandafter\def\csname ccval@sweep.pair429.A1-27b-default.findings_ungrounded\endcsname{0}
\expandafter\def\csname ccval@sweep.pair429.A1-27b-default.graded\endcsname{129}
\expandafter\def\csname ccval@sweep.pair429.A1-27b-default.grounded\endcsname{129}
\expandafter\def\csname ccval@sweep.pair429.A1-27b-default.merged_bytes\endcsname{5673}
\expandafter\def\csname ccval@sweep.pair429.A1-27b-default.model\endcsname{qwen3.8:27b}
\expandafter\def\csname ccval@sweep.pair429.A1-27b-default.segments\endcsname{63}
\expandafter\def\csname ccval@sweep.pair429.A1-27b-default.structural_findings\endcsname{0}
\expandafter\def\csname ccval@sweep.pair429.A1-27b-default.undeclared_absence\endcsname{0}
\expandafter\def\csname ccval@sweep.pair429.A1-27b-default.union_coverage_merged\endcsname{0.7864}
\expandafter\def\csname ccval@sweep.pair429.A1-27b-default.verbatim_violation\endcsname{0}
\expandafter\def\csname ccval@sweep.pair429.A1-27b-default.verify_findings\endcsname{0}
\expandafter\def\csname ccval@sweep.pair429.A2-27b-think.declared_drops\endcsname{0}
\expandafter\def\csname ccval@sweep.pair429.A2-27b-think.draws_graded\endcsname{3}
\expandafter\def\csname ccval@sweep.pair429.A2-27b-think.draws_identical\endcsname{3}
\expandafter\def\csname ccval@sweep.pair429.A2-27b-think.draws_k\endcsname{3}
\expandafter\def\csname ccval@sweep.pair429.A2-27b-think.exit_code\endcsname{1}
\expandafter\def\csname ccval@sweep.pair429.A2-27b-think.findings_grounded\endcsname{1}
\expandafter\def\csname ccval@sweep.pair429.A2-27b-think.findings_ungrounded\endcsname{1}
\expandafter\def\csname ccval@sweep.pair429.A2-27b-think.graded\endcsname{122}
\expandafter\def\csname ccval@sweep.pair429.A2-27b-think.grounded\endcsname{121}
\expandafter\def\csname ccval@sweep.pair429.A2-27b-think.merged_bytes\endcsname{5673}
\expandafter\def\csname ccval@sweep.pair429.A2-27b-think.model\endcsname{qwen3.8:27b}
\expandafter\def\csname ccval@sweep.pair429.A2-27b-think.segments\endcsname{63}
\expandafter\def\csname ccval@sweep.pair429.A2-27b-think.structural_findings\endcsname{0}
\expandafter\def\csname ccval@sweep.pair429.A2-27b-think.undeclared_absence\endcsname{0}
\expandafter\def\csname ccval@sweep.pair429.A2-27b-think.union_coverage_merged\endcsname{0.5193}
\expandafter\def\csname ccval@sweep.pair429.A2-27b-think.verbatim_violation\endcsname{0}
\expandafter\def\csname ccval@sweep.pair429.A2-27b-think.verify_findings\endcsname{2}
\expandafter\def\csname ccval@sweep.pair429.B1-70b-default.declared_drops\endcsname{0}
\expandafter\def\csname ccval@sweep.pair429.B1-70b-default.draws_graded\endcsname{3}
\expandafter\def\csname ccval@sweep.pair429.B1-70b-default.draws_identical\endcsname{3}
\expandafter\def\csname ccval@sweep.pair429.B1-70b-default.draws_k\endcsname{3}
\expandafter\def\csname ccval@sweep.pair429.B1-70b-default.exit_code\endcsname{1}
\expandafter\def\csname ccval@sweep.pair429.B1-70b-default.findings_grounded\endcsname{1}
\expandafter\def\csname ccval@sweep.pair429.B1-70b-default.findings_ungrounded\endcsname{7}
\expandafter\def\csname ccval@sweep.pair429.B1-70b-default.graded\endcsname{65}
\expandafter\def\csname ccval@sweep.pair429.B1-70b-default.grounded\endcsname{58}
\expandafter\def\csname ccval@sweep.pair429.B1-70b-default.merged_bytes\endcsname{4234}
\expandafter\def\csname ccval@sweep.pair429.B1-70b-default.model\endcsname{deepseek-r1:70b}
\expandafter\def\csname ccval@sweep.pair429.B1-70b-default.segments\endcsname{63}
\expandafter\def\csname ccval@sweep.pair429.B1-70b-default.structural_findings\endcsname{29}
\expandafter\def\csname ccval@sweep.pair429.B1-70b-default.undeclared_absence\endcsname{24}
\expandafter\def\csname ccval@sweep.pair429.B1-70b-default.union_coverage_merged\endcsname{0.8158}
\expandafter\def\csname ccval@sweep.pair429.B1-70b-default.verbatim_violation\endcsname{5}
\expandafter\def\csname ccval@sweep.pair429.B1-70b-default.verify_findings\endcsname{8}
\expandafter\def\csname ccval@sweep.pair429.B2-70b-think.declared_drops\endcsname{0}
\expandafter\def\csname ccval@sweep.pair429.B2-70b-think.draws_graded\endcsname{3}
\expandafter\def\csname ccval@sweep.pair429.B2-70b-think.draws_identical\endcsname{3}
\expandafter\def\csname ccval@sweep.pair429.B2-70b-think.draws_k\endcsname{3}
\expandafter\def\csname ccval@sweep.pair429.B2-70b-think.exit_code\endcsname{1}
\expandafter\def\csname ccval@sweep.pair429.B2-70b-think.findings_grounded\endcsname{1}
\expandafter\def\csname ccval@sweep.pair429.B2-70b-think.findings_ungrounded\endcsname{7}
\expandafter\def\csname ccval@sweep.pair429.B2-70b-think.graded\endcsname{65}
\expandafter\def\csname ccval@sweep.pair429.B2-70b-think.grounded\endcsname{58}
\expandafter\def\csname ccval@sweep.pair429.B2-70b-think.merged_bytes\endcsname{4234}
\expandafter\def\csname ccval@sweep.pair429.B2-70b-think.model\endcsname{deepseek-r1:70b}
\expandafter\def\csname ccval@sweep.pair429.B2-70b-think.segments\endcsname{63}
\expandafter\def\csname ccval@sweep.pair429.B2-70b-think.structural_findings\endcsname{29}
\expandafter\def\csname ccval@sweep.pair429.B2-70b-think.undeclared_absence\endcsname{24}
\expandafter\def\csname ccval@sweep.pair429.B2-70b-think.union_coverage_merged\endcsname{0.8158}
\expandafter\def\csname ccval@sweep.pair429.B2-70b-think.verbatim_violation\endcsname{5}
\expandafter\def\csname ccval@sweep.pair429.B2-70b-think.verify_findings\endcsname{8}
\expandafter\def\csname ccval@sweep.pair429.C1-120b-default.draws_errored\endcsname{5}
\expandafter\def\csname ccval@sweep.pair429.C1-120b-default.draws_identical\endcsname{5}
\expandafter\def\csname ccval@sweep.pair429.C1-120b-default.draws_k\endcsname{5}
\expandafter\def\csname ccval@sweep.pair429.C1-120b-default.exit_code\endcsname{2}
\expandafter\def\csname ccval@sweep.pair429.C1-120b-default.model\endcsname{gpt-oss:120b}
\expandafter\def\csname ccval@sweep.pair429.C2-120b-think.declared_drops\endcsname{0}
\expandafter\def\csname ccval@sweep.pair429.C2-120b-think.draws_graded\endcsname{5}
\expandafter\def\csname ccval@sweep.pair429.C2-120b-think.draws_identical\endcsname{1}
\expandafter\def\csname ccval@sweep.pair429.C2-120b-think.draws_k\endcsname{5}
\expandafter\def\csname ccval@sweep.pair429.C2-120b-think.exit_code\endcsname{1}
\expandafter\def\csname ccval@sweep.pair429.C2-120b-think.findings_grounded\endcsname{0}
\expandafter\def\csname ccval@sweep.pair429.C2-120b-think.findings_ungrounded\endcsname{0}
\expandafter\def\csname ccval@sweep.pair429.C2-120b-think.graded\endcsname{85}
\expandafter\def\csname ccval@sweep.pair429.C2-120b-think.grounded\endcsname{80}
\expandafter\def\csname ccval@sweep.pair429.C2-120b-think.merged_bytes\endcsname{5690}
\expandafter\def\csname ccval@sweep.pair429.C2-120b-think.model\endcsname{gpt-oss:120b}
\expandafter\def\csname ccval@sweep.pair429.C2-120b-think.segments\endcsname{63}
\expandafter\def\csname ccval@sweep.pair429.C2-120b-think.structural_findings\endcsname{3}
\expandafter\def\csname ccval@sweep.pair429.C2-120b-think.undeclared_absence\endcsname{3}
\expandafter\def\csname ccval@sweep.pair429.C2-120b-think.union_coverage_merged\endcsname{0.3873}
\expandafter\def\csname ccval@sweep.pair429.C2-120b-think.verbatim_violation\endcsname{0}
\expandafter\def\csname ccval@sweep.pair429.C2-120b-think.verify_findings\endcsname{0}
\expandafter\def\csname ccval@tool.checks\endcsname{9}
\expandafter\def\csname ccval@tool.default_depth\endcsname{full}
\expandafter\def\csname ccval@tool.default_level\endcsname{high}
\expandafter\def\csname ccval@tool.fan_in\endcsname{3}
\expandafter\def\csname ccval@tool.finding_kinds\endcsname{13}
\expandafter\def\csname ccval@tool.levels\endcsname{6}
\expandafter\def\csname ccval@tool.loss_budget_percent\endcsname{3}
\expandafter\def\csname ccval@tool.max_sources\endcsname{12}
\expandafter\def\csname ccval@tool.near_match\endcsname{0.68}
\expandafter\def\csname ccval@tool.verify_batch_command\endcsname{100}
\expandafter\def\csname ccval@tool.verify_batch_http\endcsname{25}
\makeatother
}{%
  \PackageError{llossless}{generated/numbers.tex is missing}%
    {Run `make numbers` first. Result sections cannot be typeset without it.}}
